% Déséquilibres énergétiques des organismes et leurs conséquences — SVT 1ère TI — Cours signé OnBuch+
% Programme officiel MINESEC. Compiler avec : tectonic NCT02-desequilibres-energetiques-organismes.tex
\newcommand{\DOCMATIERE}{SVT}
\newcommand{\DOCNIVEAU}{1ère TI}
\newcommand{\DOCMODULE}{Module 1 : Le Monde Vivant}
\newcommand{\DOCLECON}{NCT02}
\newcommand{\DOCTITRE}{Déséquilibres énergétiques des organismes et leurs conséquences}
% OnBuch+ — préambule « cours signés » ENRICHI (compilé avec tectonic / XeLaTeX)
% Système visuel « Pop » d'OnBuch+ : fond crème (jamais blanc), bordures encre
% épaisses, ombres « dures » décalées, titres Archivo Black en capitales.
% Version MATHÉMATIQUES : graphes (pgfplots/tikz), boîtes pédagogiques
% (définition, propriété, méthode, attention, le savais-tu, exercice résolu,
% exercice, TP, bilan), page de garde et signature OnBuch+ ancrée en bas.
%
% Macros attendues dans main.tex AVANT \input{preamble} :
%   \DOCMATIERE  \DOCNIVEAU  \DOCMODULE  \DOCLECON (numéro)  \DOCTITRE
\documentclass[11pt]{article}

\usepackage{fontspec}
\usepackage[french]{babel}
\usepackage[a4paper,margin=2cm,top=3.4cm,bottom=2.8cm,headheight=1.4cm,footskip=1.3cm]{geometry}
\usepackage{amsmath,amssymb,amsfonts}
\usepackage{mathtools} % \xmapsto, \coloneqq... souvent utilisés par les modèles
\usepackage{cancel} % simplifications barrées (bilans de forces)
% \abs{z} (module, valeur absolue) et \norm{u} (norme), utilisés par les modèles
\providecommand{\abs}{}\let\abs\relax\DeclarePairedDelimiter{\abs}{\lvert}{\rvert}
\providecommand{\norm}{}\let\norm\relax\DeclarePairedDelimiter{\norm}{\lVert}{\rVert}
\usepackage[table]{xcolor} % [table] : fournit aussi \rowcolors (alternance de lignes)
\usepackage{pagecolor}
\usepackage{tikz}
\usetikzlibrary{arrows.meta,positioning,calc,shapes.geometric,shapes.misc,shapes.arrows,decorations.pathmorphing,decorations.pathreplacing,decorations.markings,patterns,shadows,mindmap,trees,fit,backgrounds,matrix,intersections,angles,quotes}
\usepackage{pgfplots}
\usepackage[version=4]{mhchem} % équations nucléaires \ce{^{235}_{92}U}
\usepackage[european,siunitx]{circuitikz} % schémas électriques
\pgfplotsset{compat=1.18}
\usepgfplotslibrary{fillbetween,groupplots} % aires entre courbes (calcul intégral)
\usepackage{enumitem}
\usepackage{fancyhdr}
% [most] charge toutes les bibliothèques tcolorbox (listings, minted, raster,
% documentation...) et gonfle inutilement la mémoire TeX sur un document long
% (~300 boîtes) : on ne charge que ce qui est réellement utilisé.
\usepackage[skins,breakable]{tcolorbox}
\usepackage{graphicx}
\usepackage{colortbl}
\usepackage{booktabs}
\usepackage{tabularx}
\usepackage{diagbox} % case d'en-tête diagonale des tableaux à double entrée
% Colonnes extensibles centrée / gauche / droite (C, L, R), souvent utilisées par les modèles
\newcolumntype{C}{>{\centering\arraybackslash}X}
\newcolumntype{L}{>{\raggedright\arraybackslash}X}
\newcolumntype{R}{>{\raggedleft\arraybackslash}X}
\usepackage{array}
\usepackage{multicol}
\usepackage{siunitx}
% parse-numbers=false : accepte aussi \SI{2,5 \times 10^{-3}}{...}
\sisetup{locale=FR,output-decimal-marker={,},group-minimum-digits=4,parse-numbers=false}
\DeclareSIUnit\ohm{\ensuremath{\symup{\Omega}}} % U+2126 absent de la police
\DeclareSIUnit\division{div} % réglages d'oscilloscope (ms/div, V/div)
\DeclareSIUnit\tr{tr} % tours (vitesse de rotation, ex. \SI{1500}{\tr\per\minute})
\DeclareSIUnit\ch{ch} % cheval-vapeur (puissance moteur, unité usuelle hors SI)
\DeclareSIUnit\franccfa{FCFA} % franc CFA (coûts, contexte camerounais)
\DeclareSIUnit\dioptre{\ensuremath{\delta}} % vergence d'une lentille (dioptrie)
\usepackage{qrcode}
% \degree : commande fréquemment utilisée par les modèles pour un angle ou une
% température en dehors de \SI{}{} (non fournie par les paquets chargés).
\providecommand{\degree}{^{\circ}}
% \qed (fin de démonstration), utilisé par les modèles sans amsthm
\providecommand{\qed}{\hfill$\square$}
% \barre : double barre des valeurs interdites dans un tableau de variations (tkz-tab)
\providecommand{\barre}{\ensuremath{\|}}
% environnement proof (amsthm non chargé) utilisé par les modèles
\ifdefined\proof\else\newenvironment{proof}[1][Démonstration]{\par\noindent\textit{#1.}\ }{\qed\par}\fi
\usepackage{lastpage}
\usepackage{hyperref}
\hypersetup{hidelinks}

% ============================================================
% Palette « Pop » OnBuch+ — mêmes hex que la classe P (plus_ui.dart)
% ============================================================
\definecolor{poporange}{HTML}{F59321}
\definecolor{popdark}{HTML}{E07A0C}
\definecolor{popcream}{HTML}{FFF6E8}
\definecolor{popink}{HTML}{1C1714}
\definecolor{poppurple}{HTML}{7A5AE0}
\definecolor{popgreen}{HTML}{2E9E6A}
\definecolor{poppink}{HTML}{E0457B}
\definecolor{popblue}{HTML}{2F7DE1}
\definecolor{popgold}{HTML}{C98A12}
\definecolor{popmuted}{HTML}{8A7F76}
\definecolor{popline}{HTML}{EADFD2}
\definecolor{popgray}{HTML}{8A7F76}
\colorlet{popgrey}{popgray}
% Alias tolérants (le modèle nomme parfois les couleurs différemment)
\colorlet{popwhite}{white}
\colorlet{popblack}{popink}
\colorlet{popred}{poppink}
\colorlet{popyellow}{popgold}
% Teintes claires pour fonds de tableaux / zones
\colorlet{popblueL}{popblue!10!white}
\colorlet{poporangeL}{poporange!14!white}
\colorlet{popgreenL}{popgreen!12!white}
\colorlet{poppurpleL}{poppurple!12!white}
\colorlet{poppinkL}{poppink!12!white}
\colorlet{popgoldL}{popgold!14!white}

\pagecolor{popcream}

% ============================================================
% Typographie
% ============================================================
\defaultfontfeatures{Ligatures=TeX}
\setmainfont{PlusJakartaSans-Medium.ttf}[Path=../../../fonts/,BoldFont=PlusJakartaSans-Bold.ttf]
% Maths en sans-serif (Fira Math) avec chiffres et lettres droites de Plus
% Jakarta, pour que les formules se fondent dans le texte.
\usepackage[math-style=ISO,bold-style=ISO]{unicode-math}
\setmathfont{FiraMath-Regular.otf}[Path=../../../fonts/]
\setmathfont{PlusJakartaSans-Medium.ttf}[Path=../../../fonts/,range={up/num,up/{Latin,latin},\mathup}]
\usepackage{icomma} % virgule décimale sans espace en mode math
% Caractères mathématiques Unicode absents des polices de texte (Plus Jakarta,
% Archivo Black) : rendus par leur équivalent mathématique, en texte comme en
% formule — sinon ils s'affichent en carré vide (ex. « Arithmétique dans ℤ »).
\usepackage{newunicodechar}
\newunicodechar{ℝ}{\ensuremath{\mathbb{R}}}
\newunicodechar{ℤ}{\ensuremath{\mathbb{Z}}}
\newunicodechar{ℂ}{\ensuremath{\mathbb{C}}}
\newunicodechar{ℕ}{\ensuremath{\mathbb{N}}}
\newunicodechar{ℚ}{\ensuremath{\mathbb{Q}}}
\newunicodechar{𝔻}{\ensuremath{\mathbb{D}}}
\newunicodechar{α}{\ensuremath{\alpha}}
\newunicodechar{β}{\ensuremath{\beta}}
\newunicodechar{γ}{\ensuremath{\gamma}}
\newunicodechar{δ}{\ensuremath{\delta}}
\newunicodechar{ε}{\ensuremath{\varepsilon}}
\newunicodechar{θ}{\ensuremath{\theta}}
\newunicodechar{λ}{\ensuremath{\lambda}}
\newunicodechar{μ}{\ensuremath{\mu}}
\newunicodechar{σ}{\ensuremath{\sigma}}
\newunicodechar{τ}{\ensuremath{\tau}}
\newunicodechar{φ}{\ensuremath{\varphi}}
\newunicodechar{ψ}{\ensuremath{\psi}}
\newunicodechar{ω}{\ensuremath{\omega}}
\newunicodechar{Σ}{\ensuremath{\Sigma}}
% majuscules produites par \MakeUppercase dans les titres (α-aminés -> Α)
\newunicodechar{Α}{\ensuremath{\alpha}}
\newunicodechar{Β}{\ensuremath{\beta}}
\newunicodechar{Γ}{\ensuremath{\Gamma}}
\newunicodechar{Δ}{\ensuremath{\Delta}}
\newunicodechar{Ε}{\ensuremath{\varepsilon}}
\newunicodechar{Θ}{\ensuremath{\Theta}}
\newunicodechar{Λ}{\ensuremath{\Lambda}}
\newunicodechar{Μ}{\ensuremath{\mu}}
\newunicodechar{Τ}{\ensuremath{\tau}}
\newunicodechar{Φ}{\ensuremath{\Phi}}
\newunicodechar{Ψ}{\ensuremath{\Psi}}
\newunicodechar{Ω}{\ensuremath{\Omega}}
\newunicodechar{↦}{\ensuremath{\mapsto}}
\newunicodechar{∈}{\ensuremath{\in}}
\newunicodechar{∉}{\ensuremath{\notin}}
\newunicodechar{∧}{\ensuremath{\wedge}}
\newunicodechar{⇔}{\ensuremath{\Leftrightarrow}}
\newunicodechar{⟺}{\ensuremath{\Longleftrightarrow}}
\newunicodechar{⇒}{\ensuremath{\Rightarrow}}
\newunicodechar{⟹}{\ensuremath{\Longrightarrow}}
\newunicodechar{∘}{\ensuremath{\circ}}
\newunicodechar{∪}{\ensuremath{\cup}}
\newunicodechar{∩}{\ensuremath{\cap}}
\newunicodechar{≡}{\ensuremath{\equiv}}
\newunicodechar{⊂}{\ensuremath{\subset}}
\newunicodechar{⊥}{\ensuremath{\perp}}
\newunicodechar{∃}{\ensuremath{\exists}}
\newunicodechar{∀}{\ensuremath{\forall}}
\newunicodechar{∛}{\ensuremath{\sqrt[3]{\,}}}
\newunicodechar{∜}{\ensuremath{\sqrt[4]{\,}}}
\newunicodechar{⃗}{\ensuremath{{}^{\to}}}
\newunicodechar{ⁿ}{\ifmmode^{n}\else\textsuperscript{\ensuremath{n}}\fi}
\newunicodechar{ˣ}{\ifmmode^{x}\else\textsuperscript{\ensuremath{x}}\fi}
\newunicodechar{ᵇ}{\ifmmode^{b}\else\textsuperscript{\ensuremath{b}}\fi}
\newunicodechar{ʳ}{\ifmmode^{r}\else\textsuperscript{\ensuremath{r}}\fi}
\newunicodechar{ᵘ}{\ifmmode^{u}\else\textsuperscript{\ensuremath{u}}\fi}
\newunicodechar{ʸ}{\ifmmode^{y}\else\textsuperscript{\ensuremath{y}}\fi}
\newunicodechar{ᵃ}{\ifmmode^{a}\else\textsuperscript{\ensuremath{a}}\fi}
\newunicodechar{ᵉ}{\ifmmode^{e}\else\textsuperscript{\ensuremath{e}}\fi}
\newunicodechar{ᵏ}{\ifmmode^{k}\else\textsuperscript{\ensuremath{k}}\fi}
\newunicodechar{ᵖ}{\ifmmode^{p}\else\textsuperscript{\ensuremath{p}}\fi}
\newunicodechar{ᴺ}{\ifmmode^{N}\else\textsuperscript{\ensuremath{N}}\fi}
\newunicodechar{ᶜ}{\ifmmode^{c}\else\textsuperscript{\ensuremath{c}}\fi}
\newunicodechar{ʰ}{\ifmmode^{h}\else\textsuperscript{\ensuremath{h}}\fi}
\newunicodechar{ᵠ}{\ifmmode^{q}\else\textsuperscript{\ensuremath{q}}\fi}
\newunicodechar{ₙ}{\ifmmode_{n}\else\textsubscript{\ensuremath{n}}\fi}
\newunicodechar{ₐ}{\ifmmode_{a}\else\textsubscript{\ensuremath{a}}\fi}
\newunicodechar{ᵢ}{\ifmmode_{i}\else\textsubscript{\ensuremath{i}}\fi}
\newunicodechar{ᵣ}{\ifmmode_{r}\else\textsubscript{\ensuremath{r}}\fi}
\newunicodechar{ₖ}{\ifmmode_{k}\else\textsubscript{\ensuremath{k}}\fi}
\newunicodechar{ᵦ}{\ifmmode_{\beta}\else\textsubscript{\ensuremath{\beta}}\fi}
\newunicodechar{ₓ}{\ifmmode_{x}\else\textsubscript{\ensuremath{x}}\fi}
\newfontfamily\popdisplay{ArchivoBlack-Regular.ttf}[Path=../../../fonts/]
\newfontfamily\poplogo{SpaceGrotesk-Bold.ttf}[Path=../../../fonts/]
\newfontfamily\popsemi{PlusJakartaSans-SemiBold.ttf}[Path=../../../fonts/]
\newfontfamily\popxbold{PlusJakartaSans-ExtraBold.ttf}[Path=../../../fonts/]
\newfontfamily\popxboldls{PlusJakartaSans-ExtraBold.ttf}[Path=../../../fonts/,LetterSpace=3.0]

\setlist[enumerate]{leftmargin=1.7em,itemsep=5pt,topsep=4pt}
\setlist[itemize]{leftmargin=1.7em,itemsep=4pt,topsep=4pt,label={\color{poporange}\textbullet}}
\setlength{\parindent}{0pt}
\setlength{\parskip}{5pt}
\renewcommand{\arraystretch}{1.25}

% pgfplots : style Pop par défaut
\pgfplotsset{
  every axis/.append style={axis line style={popink,line width=1pt},tick style={popink},
    grid style={popline,line width=0.5pt},label style={font=\small},tick label style={font=\footnotesize}},
  popcurve/.style={poporange,line width=1.8pt,no markers,smooth},
}
% Même style disponible dans un \draw TikZ simple (hors pgfplots)
\tikzset{popcurve/.style={poporange,line width=1.8pt}}
\tikzset{popgrid/.style={popline,very thin}}
\colorlet{popcurve}{poporange} % le nom du style est parfois utilisé comme couleur

% ============================================================
% Pill / squiggle
% ============================================================
\newcommand{\poppill}[3]{%
  \tikz[baseline=-0.55ex]\node[draw=popink,line width=1.2pt,rounded corners=2.6pt,%
    fill=#2,inner xsep=6.5pt,inner ysep=3pt]{%
    \popxboldls\fontsize{7}{7}\selectfont\color{#3}\MakeUppercase{#1}};%
}
\newcommand{\popsquiggle}[1]{%
  \tikz[baseline=-0.4ex]{
    \draw[#1,line width=2.6pt,line cap=round]
      plot [smooth] coordinates {(0,0.09) (0.34,0.01) (0.68,0.21) (1.02,0.01) (1.36,0.10)};
  }%
}
% Mot-clé surligné (marqueur orange sous le texte)
\newcommand{\cle}[1]{{\popxbold\color{popdark}#1}}

% ============================================================
% En-tête / pied
% ============================================================
\pagestyle{fancy}
\fancyhf{}
\renewcommand{\headrulewidth}{0pt}
\renewcommand{\footrulewidth}{0pt}
\lhead{\raisebox{-0.15ex}{\poplogo\fontsize{13}{13}\selectfont\color{popink}OnBuch\color{poporange}+}}
\rhead{\poppill{\DOCMATIERE\ \textbullet\ \DOCNIVEAU\ \textbullet\ Leçon \DOCLECON}{white}{popink}}
\lfoot{\popxbold\fontsize{7.6}{7.6}\selectfont\color{popmuted}\MakeUppercase{Cours signé OnBuch+}\ \textbullet\ {\popsemi\DOCTITRE}}
\rfoot{\tikz[baseline=-0.6ex]\node[rounded corners=8.5pt,fill=popink,minimum height=17pt,minimum width=17pt,inner xsep=5pt,inner sep=0]{\popxbold\fontsize{8}{8}\selectfont\color{popcream}\thepage\,/\,\pageref*{LastPage}};}

% ============================================================
% Titres
% ============================================================
\newcommand{\chaptitre}[2]{%
  \begin{tcolorbox}[enhanced,colback=poporange,colframe=popink,boxrule=2.6pt,arc=11pt,
      drop shadow=popink,left=15pt,right=15pt,top=12pt,bottom=11pt,before skip=3pt,after skip=18pt]
    \poppill{#2}{popink}{popcream}\par\vspace{9pt}
    {\raggedright\hyphenpenalty=10000\exhyphenpenalty=10000\popdisplay\fontsize{22}{24}\selectfont\color{popcream}\MakeUppercase{#1}\par}
  \end{tcolorbox}
}
% Section numérotée : pastille orange avec le numéro + titre Archivo
\newcounter{popsec}
\newcommand{\coursec}[1]{%
  \stepcounter{popsec}\setcounter{popsub}{0}%
  \par\vspace{16pt}\noindent
  \begin{minipage}{\linewidth}
  \tikz[baseline=-0.7ex]\node[draw=popink,line width=1.6pt,fill=poporange,rounded corners=4pt,minimum size=22pt,inner sep=0]{\popdisplay\fontsize{12}{12}\selectfont\color{popcream}\thepopsec};%
  \hspace{8pt}{\popdisplay\fontsize{15}{17}\selectfont\color{popink}\MakeUppercase{#1}}\par
  \vspace{4pt}\noindent{\color{poporange}\rule{36pt}{3.6pt}}
  \end{minipage}\par\nopagebreak\vspace{8pt}
}
\newcounter{popsub}
\newcommand{\courssub}[1]{%
  \stepcounter{popsub}%
  \par\vspace{9pt}\noindent{\popxbold\fontsize{11.5}{13}\selectfont\color{popink}{\color{poporange}\thepopsec.\thepopsub}\hspace{6pt}#1}\par\nopagebreak\vspace{4pt}
}

% ============================================================
% Boîtes pédagogiques — même recette : fond blanc, bordure épaisse colorée,
% ombre dure encre, badge plein. Titre optionnel [#1].
% ============================================================
\tcbset{popbase/.style={breakable,enhanced,colback=white,boxrule=2.3pt,arc=9pt,
  shadow={3pt}{-3pt}{0pt}{popink},left=14pt,right=14pt,top=11pt,bottom=11pt,before skip=11pt,after skip=15pt,
  fontupper=\popsemi\fontsize{10.6}{14.5}\selectfont\color{popink},
  fontlower=\popsemi\fontsize{10.6}{14.5}\selectfont\color{popink}}}
% Coche dessinée (le glyphe ✓ n'existe pas dans les polices du document)
\newcommand{\popcheck}[1]{\tikz[baseline=-0.5ex]\draw[#1,line width=1.6pt,line cap=round,line join=round](-0.13,0.0)--(-0.04,-0.09)--(0.13,0.1);}
\providecommand{\coche}{\popcheck{popgreen}}
\providecommand{\resultat}[1]{\par\smallskip\noindent\fcolorbox{popgreen}{popgreenL}{\parbox{\dimexpr\linewidth-2\fboxsep-2\fboxrule\relax}{\textbf{Résultat.} #1}}\par\smallskip}
\newcommand{\popboxhead}[3]{\poppill{#1}{#2}{white}\ifx\relax#3\relax\else\hspace{6pt}{\popxbold\fontsize{10.6}{12}\selectfont\color{popink}#3}\fi\par\vspace{7pt}}

\newtcolorbox{aretenir}[1][]{popbase,colframe=popgreen,before upper={\popboxhead{À retenir}{popgreen}{#1}}}
\newtcolorbox{exemplebox}[1][]{popbase,colframe=poporange,before upper={\popboxhead{Exemple}{poporange}{#1}}}
\newtcolorbox{definition}[1][]{popbase,colframe=popblue,before upper={\popboxhead{Définition}{popblue}{#1}}}
\newtcolorbox{propriete}[1][]{popbase,colframe=poppurple,before upper={\popboxhead{Propriété}{poppurple}{#1}}}
\newtcolorbox{methode}[1][]{popbase,colframe=popgold,before upper={\popboxhead{Méthode}{popgold}{#1}}}
\newtcolorbox{attention}[1][]{popbase,colframe=poppink,before upper={\popboxhead{Attention}{poppink}{#1}}}
\newtcolorbox{experience}[1][]{popbase,colframe=popblue,colback=popblueL,before upper={\popboxhead{Expérience}{popblue}{#1}}}
% « Le savais-tu ? » : carte violette pleine (contexte, culture, Cameroun)
\newtcolorbox{savaistu}[1][]{popbase,colback=poppurple,colframe=popink,
  fontupper=\popsemi\fontsize{10.4}{14.2}\selectfont\color{white},
  before upper={\poppill{Le savais-tu\,?}{white}{poppurple}\ifx\relax#1\relax\else\hspace{6pt}{\popxbold\color{white}#1}\fi\par\vspace{7pt}}}
% Exercice résolu : énoncé en haut, \tcblower puis la solution sur fond crème
\newcounter{popexo}\newcounter{popexr}
\newtcolorbox{exoresolu}[1][]{popbase,colframe=poporange,colbacklower=popcream,
  segmentation style={popink,line width=1.2pt,dashed},
  before upper={\stepcounter{popexr}\popboxhead{Exercice résolu \thepopexr}{poporange}{#1}},
  before lower={\poppill{Solution}{popink}{popcream}\par\vspace{6pt}}}
% Exercice (non résolu en place) — la correction est dans la partie Corrigés
\newtcolorbox{exercice}[1][]{popbase,colframe=popink,
  before upper={\stepcounter{popexo}\popboxhead{Exercice \thepopexo}{popink}{#1}}}
% Corrigé
\newtcolorbox{corrige}[1][]{popbase,colframe=popgreen,colback=popgreenL,
  before upper={\popboxhead{Corrigé}{popgreen}{#1}}}
% Objectifs / prérequis / bilan
\newtcolorbox{objectifs}{popbase,colframe=popink,colback=poporangeL,
  before upper={\popboxhead{Objectifs de la leçon}{popink}{}}}
\newtcolorbox{prerequis}{popbase,colframe=popink,colback=popblueL,
  before upper={\popboxhead{Prérequis}{popblue}{}}}
\newtcolorbox{bilan}{popbase,colframe=popink,colback=white,boxrule=2.8pt,
  before upper={\popboxhead{Fiche bilan}{poporange}{L'essentiel à connaître pour l'examen}}}
% Figure encadrée avec légende
\newtcolorbox{popfigure}[1][]{enhanced,breakable=false,colback=white,colframe=popline,boxrule=1.4pt,arc=8pt,
  left=10pt,right=10pt,top=10pt,bottom=8pt,before skip=10pt,after skip=12pt,halign=center,
  lower separated=false,halign lower=center,
  fontlower=\popsemi\fontsize{9}{11.5}\selectfont\color{popmuted},
  #1}
\newcommand{\legende}[1]{\par\vspace{6pt}{\popsemi\fontsize{9}{11.5}\selectfont\color{popmuted}#1}}

% ============================================================
% Signature OnBuch+
% ============================================================
\newcommand{\poplogoinline}[1]{{\poplogo\fontsize{#1}{#1}\selectfont\color{popink}OnBuch\color{poporange}+}}
\newcommand{\signatureonbuch}{%
  \begin{tcolorbox}[enhanced,colback=popink,colframe=popink,boxrule=2.2pt,arc=10pt,
      drop shadow=poporange,left=14pt,right=14pt,top=10pt,bottom=10pt,before skip=0pt,after skip=0pt]
    \tikz[baseline=-0.7ex]\node[circle,fill=popgreen,draw=popcream,line width=1.3pt,minimum size=22pt,inner sep=0]{\popcheck{white}};%
    \hspace{9pt}\begin{minipage}[c]{0.62\linewidth}
      {\popxbold\fontsize{12}{13}\selectfont\color{popcream}Signé\ \textbullet\ {\poplogo OnBuch\color{poporange}+}}\par\vspace{2pt}
      {\raggedright\popsemi\fontsize{8}{10.5}\selectfont\color{popline}\MakeUppercase{Équipe pédagogique OnBuch+}\\\MakeUppercase{Conforme au programme officiel MINESEC}\par}
    \end{minipage}\hfill
    {\poplogo\fontsize{22}{22}\selectfont\color{popcream}OnBuch\color{poporange}+}
  \end{tcolorbox}%
}

% ============================================================
% Page de garde
% #1 titre · #2 matière · #3 niveau · #4 module · #5 n° leçon · #6 durée ·
% #7 description / objectifs (texte ou itemize)
% ============================================================
\newcommand{\pagegarde}[7]{%
  \thispagestyle{empty}
  \newgeometry{a4paper,margin=1.7cm,top=1.6cm,bottom=1.4cm}%
  \begin{tikzpicture}[remember picture,overlay]
    \fill[poporange!18] ([xshift=-3.2cm,yshift=-3.6cm]current page.north east) circle (4.6cm);
    \fill[poppurple!14] ([xshift=2.4cm,yshift=7.6cm]current page.south west) circle (3.8cm);
  \end{tikzpicture}%
  {\poplogo\fontsize{30}{30}\selectfont\color{popink}OnBuch\color{poporange}+}\hfill
  \raisebox{0.6ex}{\poppill{Cours signé \textbullet\ Premium}{popink}{popcream}}\par
  \vspace{1pt}\popsquiggle{poppurple}\par
  \vspace{8pt}
  \begin{tcolorbox}[enhanced,colback=poporange,colframe=popink,boxrule=2.8pt,arc=13pt,
      drop shadow=popink,left=18pt,right=18pt,top=12pt,bottom=13pt,after skip=10pt]
    \poppill{#2 \textbullet\ #3}{popink}{popcream}\hspace{5pt}\poppill{#4}{white}{popink}\par\vspace{8pt}
    {\popdisplay\fontsize{14}{14}\selectfont\color{popink}LEÇON #5}\par\vspace{4pt}
    {\raggedright\hyphenpenalty=10000\exhyphenpenalty=10000\popdisplay\fontsize{25}{27}\selectfont\color{popcream}\MakeUppercase{#1}\par}
    \vspace{8pt}
    {\popxbold\fontsize{9.5}{9.5}\selectfont\color{popink}\MakeUppercase{Durée conseillée : #6}}
  \end{tcolorbox}
  \begin{tcolorbox}[enhanced,colback=white,colframe=popink,boxrule=2.2pt,arc=10pt,
      drop shadow=popink,left=16pt,right=16pt,top=10pt,bottom=10pt,after skip=8pt]
    \poppill{Dans cette leçon}{poppurple}{white}\par\vspace{5pt}
    \popsemi\fontsize{9.3}{12.6}\selectfont\color{popink}#7
  \end{tcolorbox}
  \noindent\begin{minipage}[t]{0.487\linewidth}
    \begin{tcolorbox}[enhanced,colback=popgreen,colframe=popink,boxrule=2.2pt,arc=10pt,
        drop shadow=popink,left=11pt,right=11pt,top=10pt,bottom=10pt,height=3.1cm,valign=center]
      \tcbox[colback=white,colframe=popink,boxrule=1.3pt,arc=5pt,boxsep=0pt,nobeforeafter,
        left=3pt,right=3pt,top=3pt,bottom=3pt,valign=center]{\qrcode[height=1.9cm]{https://play.google.com/store/apps/details?id=com.luvvix.onbuch&hl=fr}}%
      \hspace{8pt}\begin{minipage}[c]{0.5\linewidth}
        {\popdisplay\fontsize{11}{12.5}\selectfont\color{white}\MakeUppercase{Télécharge OnBuch}}\par\vspace{4pt}
        {\raggedright\popsemi\fontsize{8.4}{11}\selectfont\color{white}Gratuit \textbullet\ Google Play\\Tuteur IA, annales, concours, résultats\par}
      \end{minipage}
    \end{tcolorbox}
  \end{minipage}\hfill
  \begin{minipage}[t]{0.487\linewidth}
    \begin{tcolorbox}[enhanced,colback=poppurple,colframe=popink,boxrule=2.2pt,arc=10pt,
        drop shadow=popink,left=11pt,right=11pt,top=10pt,bottom=10pt,height=3.1cm,valign=center]
      \tikz[baseline=-0.7ex]\node[circle,fill=white,draw=popink,line width=1.3pt,minimum size=1.6cm,inner sep=0]{\popdisplay\fontsize{24}{24}\selectfont\color{poppurple}+};%
      \hspace{8pt}\begin{minipage}[c]{0.6\linewidth}
        {\popdisplay\fontsize{11}{12.5}\selectfont\color{white}\MakeUppercase{Passe à OnBuch+}}\par\vspace{4pt}
        {\raggedright\popsemi\fontsize{8.4}{11}\selectfont\color{white}Tous les cours signés, Tuteur IA illimité, fiches PDF — onglet OnBuch+ de l'app\par}
      \end{minipage}
    \end{tcolorbox}
  \end{minipage}
  \vspace{8pt}
  \popcontenu
  \vfill
  \signatureonbuch
  \restoregeometry
  \newpage
}

% Rangée « Ce cours contient » (page de garde)
\newcommand{\poptile}[3]{% #1 couleur #2 gros texte #3 légende
  \begin{tcolorbox}[enhanced,colback=#1,colframe=popink,boxrule=2pt,arc=9pt,shadow={2.5pt}{-2.5pt}{0pt}{popink},
      left=6pt,right=6pt,top=9pt,bottom=9pt,halign=center,valign=center,height=1.75cm,width=\linewidth,nobeforeafter]
    {\popdisplay\fontsize{12.5}{13}\selectfont\color{white}\MakeUppercase{#2}}\par\vspace{3pt}
    {\centering\popsemi\fontsize{7.8}{9.5}\selectfont\color{white}#3\par}
  \end{tcolorbox}}
\newcommand{\popcontenu}{%
  \par\noindent{\popxboldls\fontsize{7.5}{7.5}\selectfont\color{popmuted}CE COURS SIGNÉ CONTIENT}\par\vspace{7pt}
  \noindent\begin{minipage}[t]{0.235\linewidth}\poptile{popblue}{Cours}{complet, illustré, pas à pas}\end{minipage}\hfill
  \begin{minipage}[t]{0.235\linewidth}\poptile{poporange}{Méthodes}{et exercices résolus}\end{minipage}\hfill
  \begin{minipage}[t]{0.235\linewidth}\poptile{poppink}{Exercices}{type examen + corrigés}\end{minipage}\hfill
  \begin{minipage}[t]{0.235\linewidth}\poptile{popgreen}{Fiche bilan}{l'essentiel à retenir}\end{minipage}\par}

% Sommaire
\newtcolorbox{sommaire}{enhanced,colback=white,colframe=popink,boxrule=2.2pt,arc=10pt,
  drop shadow=popink,left=18pt,right=18pt,top=13pt,bottom=13pt,before skip=6pt,after skip=20pt,
  before upper={\poppill{Sommaire}{poppurple}{white}\par\vspace{9pt}\popsemi\fontsize{10.6}{14.5}\selectfont\color{popink}}}

% Signature de fin de cours — poussée tout en bas de la dernière page
\newcommand{\findecours}{%
  \par\vfill\nopagebreak
  \begin{center}\popsquiggle{poporange}\end{center}\vspace{4pt}
  \signatureonbuch
}
\AtBeginDocument{\def\llbracket{\mathopen{[\mkern-3.5mu[}}\def\rrbracket{\mathclose{]\mkern-3.5mu]}}} % intervalles d'entiers (glyphe ⟦ absent de la police)
% Glyphes absents de Fira Math : redéfinis à partir de glyphes disponibles
\AtBeginDocument{%
  \def\setminus{\mathbin{\backslash}}\let\smallsetminus\setminus
  \def\vdots{\mathord{\vbox{\baselineskip3.2pt\lineskiplimit0pt\kern4pt\hbox{.}\hbox{.}\hbox{.}}}}%
  \def\ddots{\mathinner{\mkern1mu\raise6.5pt\hbox{.}\mkern2mu\raise3.6pt\hbox{.}\mkern2mu\raise0.7pt\hbox{.}\mkern1mu}}%
  \def\triangle{\mathord{\scalebox{0.9}{$\symup{\Delta}$}}}%
  \def\nabla{\mathord{\raisebox{\depth}{\scalebox{1}[-1]{$\symup{\Delta}$}}}}%
  \def\checkmark{\popcheck{popdark}}%
  \def\star{\mathbin{\ast}}\def\bigstar{\mathord{\ast}}%
  \def\diamond{\mathbin{\scalebox{0.6}{$\Diamond$}}}%
  \def\bigcup{\mathop{\mathchoice{\raisebox{-0.3em}{\scalebox{1.7}{$\cup$}}}{\scalebox{1.25}{$\cup$}}{\cup}{\cup}}\displaylimits}%
  \def\bigcap{\mathop{\mathchoice{\raisebox{-0.3em}{\scalebox{1.7}{$\cap$}}}{\scalebox{1.25}{$\cap$}}{\cap}{\cap}}\displaylimits}%
  \def\lesseqqgtr{\mathrel{\lessgtr}}\def\bigoplus{\mathop{\oplus}\displaylimits}%
}
\providecommand{\ssi}{si et seulement si} % abréviation fréquente
\AtBeginDocument{\providecommand{\wideparen}{\overparen}} % arc de cercle

%% FIGFIX-BEGIN v3 — correctifs globaux des figures (docs/onbuch-generation/tools/figfix)
\usepackage{adjustbox}\usepackage{etoolbox}
% toute figure TikZ/pgfplots plus large que la ligne est réduite à la largeur disponible
\BeforeBeginEnvironment{tikzpicture}{\begin{adjustbox}{max width=\linewidth}}
\AfterEndEnvironment{tikzpicture}{\end{adjustbox}}
% légendes pgfplots lisibles par-dessus les courbes
\pgfplotsset{every axis/.append style={legend style={fill=white,fill opacity=0.92,text opacity=1,draw=black!15,inner sep=3pt}}}
% étiquettes d'arêtes (syntaxe edge["..."]) avec fond blanc
\tikzset{every edge quotes/.append style={fill=white,inner sep=1.2pt}}
% bulles de cartes mentales : texte à la ligne dans la bulle, bulle un peu plus grande
\tikzset{every concept/.append style={text width=2.0cm,align=center,minimum size=2.3cm,inner sep=2pt}}
%% FIGFIX-END
\begin{document}

\pagegarde{Déséquilibres énergétiques des organismes et leurs conséquences}{SVT}{1ère TI}{Module 1 : Le Monde Vivant}{NCT02}{4 h}{Cette leçon étudie le bilan énergétique de l'organisme, c'est-à-dire l'équilibre entre les apports alimentaires et les dépenses énergétiques. Elle explique les conséquences des déséquilibres (obésité, dénutrition) et le rôle de l'activité physique dans la régulation de ce bilan.
\vspace{4pt}
\begin{multicols}{2}\small
\begin{itemize}
\item À la fin de cette leçon, je sais définir le bilan énergétique et identifier ses deux composantes (apports et dépenses)
\item À la fin de cette leçon, je sais citer les principales sources d'apports énergétiques et les postes de dépenses de l'organisme
\item À la fin de cette leçon, je sais établir un bilan énergétique simple à partir de données chiffrées (apports/dépenses)
\item À la fin de cette leçon, je sais interpréter les conséquences physiologiques d'un déséquilibre énergétique positif (surpoids, obésité)
\item À la fin de cette leçon, je sais interpréter les conséquences physiologiques d'un déséquilibre énergétique négatif (dénutrition, épuisement)
\item À la fin de cette leçon, je sais expliquer le rôle de l'activité physique dans la régulation du bilan énergétique
\end{itemize}\end{multicols}}

\begin{sommaire}
\begin{multicols}{2}
\begin{enumerate}
\item Les besoins énergétiques de l'organisme
\item Le bilan énergétique : apports et dépenses
\item Le déséquilibre énergétique positif : surpoids et obésité
\item Le déséquilibre énergétique négatif : dénutrition et épuisement
\item Rôle de l'activité physique dans la régulation du bilan énergétique
\item Rééquilibrage du bilan énergétique et mesures préventives
\item Activité d'intégration
\item Méthodes et exercices résolus
\item Exercices
\item Corrigés des exercices
\item Fiche bilan
\end{enumerate}
\end{multicols}
\end{sommaire}

\begin{objectifs}
\begin{itemize}
\item À la fin de cette leçon, je sais définir le bilan énergétique et identifier ses deux composantes (apports et dépenses)
\item À la fin de cette leçon, je sais citer les principales sources d'apports énergétiques et les postes de dépenses de l'organisme
\item À la fin de cette leçon, je sais établir un bilan énergétique simple à partir de données chiffrées (apports/dépenses)
\item À la fin de cette leçon, je sais interpréter les conséquences physiologiques d'un déséquilibre énergétique positif (surpoids, obésité)
\item À la fin de cette leçon, je sais interpréter les conséquences physiologiques d'un déséquilibre énergétique négatif (dénutrition, épuisement)
\item À la fin de cette leçon, je sais expliquer le rôle de l'activité physique dans la régulation du bilan énergétique
\item À la fin de cette leçon, je sais proposer des mesures de rééquilibrage adaptées à une situation concrète
\item À la fin de cette leçon, je sais utiliser l'IMC comme indicateur simple de l'état nutritionnel (complément utile au Probatoire)
\end{itemize}
\end{objectifs}

\begin{prerequis}
\begin{itemize}
\item Notion de nutriments (glucides, lipides, protéines) et d'aliments énergétiques vue en classe de Seconde
\item Notion de digestion et d'absorption des nutriments
\item Notion de respiration cellulaire et de production d'ATP (début d'année de Première)
\item Calculs simples : additions, soustractions, pourcentages et lecture de tableaux
\end{itemize}
\end{prerequis}

\newpage

\coursec{Les besoins énergétiques de l'organisme}

\textbf{Situation d'accroche.} Amina, élève de Première C au lycée de Yaoundé, mange chaque midi un grand plat de ndolé accompagné de miondo, et grignote le soir des beignets de haricots bien huilés. Pourtant, ses journées se déroulent presque entièrement assise : le matin en classe, l'après-midi sur son téléphone. En un an, elle a pris 12 kg et s'essouffle dès qu'elle monte les escaliers du bâtiment C. À l'inverse, son cousin Rodrigue, apprenti menuisier à Douala, porte des planches, scie et ponce toute la journée sous la chaleur, mais saute souvent le repas de midi : il maigrit, s'affaiblit et se fatigue de plus en plus vite. Deux situations opposées, un même phénomène biologique : la gestion de l'\cle{énergie} par l'organisme.

\textbf{Problématique.} Comment expliquer scientifiquement ces deux situations opposées ? D'où vient l'énergie dont notre corps a besoin, à quoi sert-elle, et comment l'organisme gère-t-il l'équilibre entre ce qu'il reçoit et ce qu'il dépense ?

\courssub{Pourquoi l'organisme a-t-il besoin d'énergie ?}

\textbf{Observation de départ.} Même lorsque tu dors profondément, ton corps n'est pas « éteint » : ton cœur bat environ 60 à 80 fois par minute, tes poumons se gonflent et se dégonflent 12 à 20 fois par minute, ta température reste stable autour de 37 °C alors que l'air ambiant est souvent plus frais, et ton cerveau continue de traiter des informations. Toutes ces activités consomment de l'énergie, en permanence, jour et nuit.

\begin{definition}[Besoin énergétique]
Le \cle{besoin énergétique} d'un individu est la quantité d'énergie nécessaire à son organisme pour assurer, d'une part, ses \textbf{fonctions vitales} (maintien de la température corporelle à 37 °C, battements cardiaques, respiration, activité cérébrale, renouvellement cellulaire, croissance chez l'enfant et l'adolescent) et, d'autre part, ses \textbf{activités} (marcher, travailler, étudier, pratiquer un sport).
\end{definition}

On distingue classiquement deux grandes composantes de la dépense énergétique :

\begin{itemize}
\item le \cle{métabolisme de base} (ou métabolisme basal) : c'est la dépense énergétique \emph{minimale}, celle d'un organisme au repos complet, à jeun, dans un environnement thermiquement neutre. Il correspond au fonctionnement des organes vitaux : cœur, poumons, cerveau, foie, reins, et au maintien de la température à 37 °C. Il représente à lui seul environ \textbf{60 à 75 \%} de la dépense totale d'une personne peu active ;
\item la \cle{dépense d'activité} : elle correspond à tous les mouvements et travaux musculaires supplémentaires (se déplacer, porter, cultiver un champ, jouer au football), ainsi qu'à la digestion (thermic effect of food) et à l'adaptation au climat.
\end{itemize}

\begin{attention}[Erreur fréquente]
Beaucoup d'élèves croient que seules les activités sportives consomment de l'énergie. C'est faux : le \textbf{métabolisme de base consomme de l'énergie en permanence, même au repos}, même pendant le sommeil. Un élève qui resterait couché toute la journée dépenserait tout de même environ 5 000 à 7 000 kJ pour faire fonctionner ses organes vitaux. C'est d'ailleurs l'erreur d'Amina : elle pense « ne rien faire », mais son corps, lui, reçoit de l'énergie excédentaire qu'il ne dépense pas.
\end{attention}

\begin{savaistu}[Un cerveau très gourmand]
Le cerveau humain ne représente que \textbf{2 \%} de la masse corporelle (environ 1,4 kg chez l'adulte), mais il consomme à lui seul environ \textbf{20 \%} de l'énergie de l'organisme au repos, essentiellement sous forme de glucose. C'est pourquoi une journée de révisions intensives ou une composition de quatre heures peut donner une véritable sensation de faim : réfléchir coûte de l'énergie !
\end{savaistu}

\courssub{Les unités de l'énergie en nutrition : joule, kilojoule et kilocalorie}

En physique, l'unité d'énergie du Système international est le \cle{joule} (J). En nutrition, les quantités mises en jeu sont grandes : on utilise donc le \cle{kilojoule} (kJ), avec $1~\text{kJ} = 1\,000~\text{J}$. Une unité plus ancienne, encore très présente sur les étiquettes des produits alimentaires, est la \cle{kilocalorie} (kcal), parfois appelée abusivement « calorie » dans le langage courant.

\begin{aretenir}[Correspondance entre kcal et kJ]
\[ 1~\text{kcal} \approx 4{,}18~\text{kJ} \qquad \text{et donc} \qquad 1~\text{kJ} \approx 0{,}24~\text{kcal}. \]
\end{aretenir}

\begin{methode}[Convertir des kilocalories en kilojoules]
Pour convertir une valeur énergétique exprimée en kilocalories en kilojoules :
\begin{enumerate}
\item repérer la valeur en kcal ;
\item la \textbf{multiplier par 4,18} ;
\item exprimer le résultat en kJ, en arrondissant si nécessaire au nombre de chiffres significatifs imposé par l'énoncé.
\end{enumerate}
\emph{Exemple.} Un adolescent dont les besoins sont estimés à 2 500 kcal par jour a des besoins de :
\[ 2\,500 \times 4{,}18 = 10\,450~\text{kJ} \approx 10\,500~\text{kJ par jour}. \]
Pour la conversion inverse (kJ vers kcal), on \textbf{divise} par 4,18.
\end{methode}

\begin{attention}[Piège de conversion]
Ne jamais additionner ou comparer directement des valeurs exprimées dans des unités différentes. Dire « 2 500 kcal $+$ 2 000 kJ » n'a aucun sens : il faut d'abord tout convertir dans la même unité. Au Probatoire, vérifie toujours l'unité demandée dans l'énoncé avant de conclure.
\end{attention}

\courssub{D'où vient l'énergie ? Les valeurs énergétiques des nutriments}

L'énergie de l'organisme provient exclusivement des \cle{nutriments énergétiques} contenus dans les aliments : les \textbf{glucides}, les \textbf{lipides} et les \textbf{protéines}. Lors de la respiration cellulaire, ces molécules sont oxydées dans les mitochondries des cellules et l'énergie libérée est captée sous forme d'ATP, la « monnaie énergétique » de la cellule. Mais tous les nutriments ne fournissent pas la même quantité d'énergie.

\begin{experience}[Mesurer l'énergie contenue dans un aliment : calorimétrie rudimentaire]
\textbf{Principe.} Brûler complètement un aliment sous une masse connue d'eau et mesurer l'élévation de température de l'eau. L'énergie libérée par la combustion est transférée (partiellement) à l'eau.\par
\textbf{Matériel.} Support à creuset, tube à essai, thermomètre, balance de précision, arachide sèche, allumettes, eau, pince à creuset.\par
\textbf{Protocole.}
\begin{enumerate}
\item Peser 1 g d'arachide séchée et la placer sur le support.
\item Verser 20 mL d'eau (soit 20 g) dans le tube à essai, mesurer la température initiale $T_i$ (ex. 25,0 °C).
\item Enflammer l'arachide et la placer immédiatement sous le tube à essai jusqu'à combustion complète.
\item Mesurer la température finale $T_f$ de l'eau (ex. 65,0 °C).
\end{enumerate}
\textbf{Observations.} La température de l'eau s'élève nettement ($ \Delta T = 40,0~°\text{C} $). La combustion d'un gramme d'arachide dégage beaucoup de chaleur. La même expérience avec 1 g de sucre (glucide pur) élève moins la température de l'eau.\par
\textbf{Interprétation.} La chaleur reçue par l'eau ($Q = m \times c \times \Delta T$) provient de l'énergie chimique stockée dans les molécules de l'aliment. Les lipides de l'arachide libèrent plus d'énergie par gramme que les glucides du sucre. Les pertes de chaleur vers l'air expliquent que la valeur mesurée soit inférieure à la valeur théorique.\par
\textbf{Conclusion.} Les aliments contiennent de l'énergie chimique, et leur valeur énergétique dépend de leur composition en nutriments.
\end{experience}

Les mesures calorimétriques de référence (facteurs d'Atwater) donnent les valeurs suivantes :

\begin{aretenir}[Valeurs énergétiques des nutriments (facteurs d'Atwater)]
\begin{itemize}
\item $1$ g de \textbf{glucides} $\approx 17$ kJ (soit environ 4 kcal) ;
\item $1$ g de \textbf{lipides} $\approx 38$ kJ (soit environ 9 kcal) ;
\item $1$ g de \textbf{protéines} $\approx 17$ kJ (soit environ 4 kcal).
\end{itemize}
Les lipides sont donc les nutriments les plus énergétiques : ils fournissent \textbf{plus du double} d'énergie par gramme que les glucides ou les protéines.
\end{aretenir}

\begin{popfigure}
\begin{center}
\begin{tikzpicture}
\pgfplotsset{compat=1.18}
\begin{axis}[
    ybar,
    width=0.85\linewidth,
    height=6.5cm,
    bar width=1.4cm,
    ymin=0, ymax=45,
    ylabel={Valeur énergétique (kJ/g)},
    symbolic x coords={Glucides, Lipides, Protéines},
    xtick=data,
    nodes near coords,
    every node near coord/.append style={font=\small\bfseries, color=popdark},
    ymajorgrids=true,
    grid style={popline!40},
    axis line style={popdark},
    tick label style={font=\small},
    enlarge x limits=0.25,
]
\addplot[fill=popblue, draw=popdark] coordinates {(Glucides,17)};
\addplot[fill=poporange, draw=popdark] coordinates {(Lipides,38)};
\addplot[fill=popgreen, draw=popdark] coordinates {(Protéines,17)};
\end{axis}
\end{tikzpicture}
\end{center}
\legende{Comparaison des valeurs énergétiques des trois nutriments énergétiques, en kilojoules par gramme. Les lipides (38 kJ/g) fournissent plus du double de l'énergie apportée par les glucides ou les protéines (17 kJ/g).}
\end{popfigure}

\courssub{Les aliments énergétiques de l'alimentation camerounaise}

L'alimentation quotidienne au Cameroun repose largement sur des aliments énergétiques : le \textbf{manioc} (consommé sous forme de bâton de manioc, de miondo ou de gari), l'\textbf{igname}, le \textbf{plantain}, le \textbf{maïs} et le \textbf{riz} sont riches en glucides ; l'\textbf{huile de palme}, l'\textbf{arachide} et le \textbf{karité} sont riches en lipides. Le tableau ci-dessous rassemble quelques valeurs énergétiques approximatives, à connaître en ordre de grandeur.

\begin{center}
\begin{tabularx}{\linewidth}{|l|X|c|c|}
\hline
\rowcolor{popblueL} \textbf{Aliment (100 g)} & \textbf{Nutriment dominant} & \textbf{Énergie (kJ)} & \textbf{Énergie (kcal)} \\
\hline
Plantain cuit & Glucides & $\approx 500$ & $\approx 120$ \\
\hline
Bâton de manioc & Glucides & $\approx 640$ & $\approx 153$ \\
\hline
Maïs (grains secs) & Glucides & $\approx 1\,500$ & $\approx 360$ \\
\hline
Arachides grillées & Lipides + protéines & $\approx 2\,400$ & $\approx 575$ \\
\hline
Huile de palme & Lipides & $\approx 3\,800$ & $\approx 900$ \\
\hline
\end{tabularx}
\end{center}

\begin{exemplebox}[L'huile de palme, un concentré d'énergie]
L'huile de palme, très utilisée dans la cuisine camerounaise (ndolé, eru, okok, sauces), est presque uniquement constituée de lipides. Une portion de 50 g d'huile de palme apporte environ :
\[ 50 \times 38 = 1\,900~\text{kJ}, \]
soit près du \textbf{cinquième} des besoins quotidiens d'un adolescent (10 000 à 12 000 kJ). Une sauce très huilée peut donc, à elle seule, couvrir une fraction considérable des besoins de la journée, sans même compter le miondo ou le plantain qui l'accompagnent. C'est l'une des clés pour comprendre la prise de poids d'Amina.
\end{exemplebox}

\courssub{De quoi dépendent les besoins énergétiques ?}

Les besoins énergétiques ne sont pas identiques pour tous : ils varient selon plusieurs facteurs qu'il faut savoir citer et justifier.

\begin{itemize}
\item \textbf{L'âge.} L'enfant et l'adolescent ont des besoins élevés car ils doivent assurer, en plus des fonctions vitales, la \textbf{croissance} (fabrication de nouveaux tissus). Les besoins par kilogramme de masse corporelle diminuent ensuite avec l'âge.
\item \textbf{Le sexe.} À âge et activité égaux, un garçon dépense en général plus d'énergie qu'une fille, car sa masse musculaire (tissu très consommateur d'énergie au repos) est en moyenne plus développée.
\item \textbf{La masse corporelle.} Plus le corps est lourd, plus le maintien des fonctions vitales et les déplacements coûtent d'énergie.
\item \textbf{L'activité physique.} Un apprenti menuisier comme Rodrigue, un cultivateur ou un footballeur dépensent beaucoup plus qu'un élève sédentaire : l'activité musculaire peut faire varier la dépense totale du simple au double.
\item \textbf{Le climat.} Le maintien de la température corporelle à 37 °C coûte davantage d'énergie en milieu froid (par exemple sur les hauts plateaux de l'Ouest ou dans la nuit fraîche du Mont Cameroun) qu'en zone chaude, car l'organisme doit produire de la chaleur pour compenser les pertes (thermorégulation).
\end{itemize}

\begin{aretenir}[Ordres de grandeur à retenir pour le Probatoire]
Un \textbf{adolescent actif} (15--18 ans) dépense environ \textbf{10 000 à 12 000 kJ par jour} (soit 2 400 à 2 900 kcal). Un adulte sédentaire dépense environ 8 000 kJ, tandis qu'un travailleur de force ou un sportif de haut niveau peut dépasser 15 000 kJ par jour.
\end{aretenir}

\courssub{Application : calculer l'énergie apportée par un repas}

\begin{exoresolu}[Énergie d'un goûter d'élève]
\textbf{Énoncé.} Pour son goûter, Amina mange 30 g d'arachides grillées et 100 g de plantain bouilli. On donne : 100 g d'arachides grillées apportent environ 2 400 kJ ; 100 g de plantain bouilli apportent environ 500 kJ. (1) Calculer l'énergie totale apportée par ce goûter, en kJ. (2) Convertir cette énergie en kcal (on donne $1~\text{kcal} \approx 4{,}18~\text{kJ}$). (3) Comparer cette énergie aux besoins quotidiens d'un adolescent actif (10 000 à 12 000 kJ) et commenter.
\tcblower
\textbf{Solution détaillée.}\par
(1) Énergie apportée par 30 g d'arachides (règle de trois) :
\[ E_1 = \frac{2\,400 \times 30}{100} = 720~\text{kJ}. \]
Énergie apportée par 100 g de plantain : $E_2 = 500$ kJ. Énergie totale du goûter :
\[ E = E_1 + E_2 = 720 + 500 = 1\,220~\text{kJ}. \]
(2) Conversion en kilocalories (on divise par 4,18) :
\[ E = \frac{1\,220}{4{,}18} \approx 292~\text{kcal}. \]
(3) Comparaison : $\dfrac{1\,220}{10\,000} \approx 0{,}12$, soit environ \textbf{12 \%} des besoins quotidiens d'un adolescent actif. Ce simple goûter, pris en plus des trois repas de la journée, représente donc un apport non négligeable. Si les apports totaux dépassent régulièrement les dépenses, l'excédent est stocké dans l'organisme sous forme de graisses (triglycérides dans le tissu adipeux) : c'est précisément le mécanisme du déséquilibre énergétique positif que nous étudierons dans la suite de la leçon.
\end{exoresolu}

\begin{exercice}[À toi de jouer]
Rodrigue, l'apprenti menuisier, dépense environ 14 000 kJ par jour en raison de son travail physique intense. Un jour où il saute le repas de midi, ses seuls apports sont : au petit-déjeuner, 100 g de beignets (environ 1 300 kJ) ; le soir, 200 g de bâton de manioc (640 kJ pour 100 g) avec une sauce contenant 30 g d'huile de palme (38 kJ/g). (1) Calcule ses apports énergétiques totaux de la journée. (2) Compare ces apports à ses dépenses et conclus sur la situation de Rodrigue.
\end{exercice}

\begin{corrige}[Exercice]
(1) Apports : beignets : 1 300 kJ ; manioc : $2 \times 640 = 1\,280$ kJ ; huile de palme : $30 \times 38 = 1\,140$ kJ. Total : $1\,300 + 1\,280 + 1\,140 = 3\,720$ kJ.\par
(2) Dépenses : 14 000 kJ. Les apports (3 720 kJ) ne couvrent qu'environ $3\,720/14\,000 \approx 27~\%$ des dépenses. Rodrigue est en \textbf{déséquilibre énergétique négatif} : son organisme doit puiser dans ses réserves (graisses du tissu adipeux, puis protéines musculaires en cas de prolongation) pour compenser le déficit, ce qui explique son amaigrissement et son affaiblissement progressif.
\end{corrige}

\textbf{Bilan de la section.} L'organisme est comparable à un moteur qui fonctionne en continu : même au repos, il dépense de l'énergie pour ses fonctions vitales (métabolisme de base), et davantage encore lors de l'activité physique. Cette énergie provient des nutriments — glucides et protéines (17 kJ/g), lipides (38 kJ/g) — apportés par les aliments. Les besoins, de l'ordre de 10 000 à 12 000 kJ par jour pour un adolescent actif, varient selon l'âge, le sexe, la masse corporelle, le climat et l'activité. Comparer les apports aux dépenses, c'est établir le \cle{bilan énergétique}, objet de la section suivante, qui permettra d'expliquer rigoureusement les situations d'Amina et de Rodrigue.

\coursec{Le bilan énergétique : apports et dépenses}

Dans la section précédente, nous avons établi que l'organisme a des \cle{besoins énergétiques} quotidiens, variables selon l'âge, le sexe, la taille et le mode de vie. Mais connaître ses besoins ne suffit pas : encore faut-il comparer ce que l'on \emph{apporte} à l'organisme (par l'alimentation) à ce qu'il \emph{dépense} réellement. C'est toute l'idée du \cle{bilan énergétique}, une notion simple dans son principe — une balance entre deux flux — mais dont la maîtrise est indispensable pour comprendre le surpoids, la dénutrition et le rôle de l'activité physique.

\courssub{Intuition : l'organisme fonctionne comme un budget}

Considérons deux élèves de Première à Yaoundé. Le premier, Nkoulou, pratique le football trois fois par semaine et marche chaque jour jusqu'au lycée ; il mange copieusement mais reste mince. La seconde, Amina, mange des quantités comparables mais passe ses journées assise, entre la salle de classe et les réseaux sociaux ; elle prend progressivement du poids. Même observation dans les villages : le cultivateur de la région de l'Ouest, qui travaille ses champs de cocoyam du matin au soir, peut consommer de grandes quantités de foutou sans grossir, alors qu'un employé de bureau sédentaire grossit avec des rations bien moindres.

\begin{experience}[Observation de départ]
On constate que deux personnes consommant des rations alimentaires identiques n'évoluent pas de la même façon : l'une maintient sa masse corporelle, l'autre grossit ou maigrit. \textbf{Hypothèse} : la masse corporelle dépend non pas des seuls apports, mais de la \emph{comparaison} entre les apports et les dépenses. \textbf{Vérification} : on mesure chez des volontaires, sur plusieurs semaines, l'énergie totale des aliments ingérés et l'énergie totale dépensée (calorimétrie). \textbf{Résultat} : les sujets dont les apports dépassent les dépenses prennent de la masse grasse ; ceux dont les dépenses dépassent les apports en perdent. \textbf{Conclusion} : c'est bien le \emph{bilan} (la différence) qui détermine l'évolution des réserves corporelles.
\end{experience}

L'intuition est donc celle d'un \textbf{budget} : l'énergie alimentaire est la « recette », les dépenses de l'organisme sont les « dépenses ». Si les recettes dépassent les dépenses, on épargne (stockage de graisses) ; si les dépenses dépassent les recettes, on puise dans l'épargne (utilisation des réserves).

\courssub{Définition du bilan énergétique}

\begin{definition}[Bilan énergétique]
Le \cle{bilan énergétique} d'un organisme est la différence, sur une période donnée (généralement une journée), entre les \cle{apports énergétiques} fournis par les aliments et boissons consommés et les \cle{dépenses énergétiques} totales de l'organisme. Il s'exprime en kilojoules par jour (kJ/jour) :
\[ \text{Bilan énergétique} = \text{Apports totaux} - \text{Dépenses totales} \]
\end{definition}

\begin{attention}[Unités]
L'unité légale d'énergie est le \textbf{joule} (J) et son multiple le \textbf{kilojoule} (kJ). On rencontre encore souvent la \textbf{kilocalorie} (kcal) sur les étiquettes alimentaires : $1\ \text{kcal} \approx 4{,}18\ \text{kJ}$. Dans les calculs, ne jamais mélanger les deux unités : convertir d'abord toutes les valeurs dans la même unité.
\end{attention}

\courssub{Les apports énergétiques : exclusivement l'alimentation}

Les apports énergétiques proviennent \textbf{exclusivement} des aliments et des boissons énergétiques consommés. L'organisme humain ne sait pas capter l'énergie solaire (contrairement aux plantes chlorophylliennes) : toute son énergie vient de l'oxydation des nutriments.

Rappel des valeurs énergétiques des nutriments (vues en classe de Seconde et rappelées dans la section 1) :

\begin{center}
\begin{tabularx}{\linewidth}{|l|X|X|}
\hline
\rowcolor{popblueL} \textbf{Nutriment} & \textbf{Valeur énergétique} & \textbf{Exemples de sources (Cameroun)} \\
\hline
Glucides & $17\ \text{kJ/g}$ & manioc, igname, riz, maïs, banane plantain \\
\hline
Lipides & $38\ \text{kJ/g}$ & huile de palme, arachide, avocat \\
\hline
Protides & $17\ \text{kJ/g}$ & poisson, viande, haricots niébé, arachides \\
\hline
Alcool (éthanol) & $29\ \text{kJ/g}$ & boissons alcoolisées (énergie « vide », sans nutriments) \\
\hline
\end{tabularx}
\end{center}

\begin{exemplebox}[Calcul des apports d'une journée]
Un élève consomme en une journée : $400\ \text{g}$ de glucides, $80\ \text{g}$ de lipides et $90\ \text{g}$ de protides. Ses apports sont :
\begin{align*}
E_{\text{glucides}} &= 400 \times 17 = 6\,800\ \text{kJ} \\
E_{\text{lipides}} &= 80 \times 38 = 3\,040\ \text{kJ} \\
E_{\text{protides}} &= 90 \times 17 = 1\,530\ \text{kJ} \\
E_{\text{total}} &= 6\,800 + 3\,040 + 1\,530 = 11\,370\ \text{kJ/jour}
\end{align*}
\end{exemplebox}

\begin{attention}[Les boissons comptent aussi]
Erreur fréquente : ne compter que les aliments solides. Or les boissons sucrées (jus industriels, sodas, bières) apportent une énergie non négligeable : une bouteille de soda de $50\ \text{cL}$ contient environ $50\ \text{g}$ de sucre, soit environ $850\ \text{kJ}$ — l'équivalent d'une heure de marche rapide. L'eau, elle, n'apporte \emph{aucune} énergie.
\end{attention}

\courssub{Les dépenses énergétiques : quatre composantes}

Contrairement à une idée répandue, l'organisme ne dépense pas de l'énergie seulement lorsqu'il « fait du sport ». La dépense énergétique totale quotidienne comprend \textbf{quatre composantes} :

\begin{enumerate}
\item \textbf{Le métabolisme de base} : c'est la dépense minimale, au repos complet, nécessaire au maintien des fonctions vitales (battements cardiaques, respiration, maintien de la température corporelle, renouvellement cellulaire, activité cérébrale).
\item \textbf{La thermorégulation} : maintien du corps à environ $37\ ^\circ\text{C}$ quelle que soit la température extérieure (frissons en cas de froid, sudation en cas de chaleur).
\item \textbf{L'activité physique} : toute contraction musculaire, de la marche jusqu'au sport intense, en passant par les travaux domestiques et agricoles.
\item \textbf{La digestion} (ou thermogenèse alimentaire) : l'énergie consommée par les processus de digestion, d'absorption et de stockage des nutriments.
\end{enumerate}

\begin{definition}[Métabolisme de base]
Le \cle{métabolisme de base} (ou métabolisme basal) est la dépense énergétique \emph{minimale} de l'organisme au repos, à jeun, dans des conditions de température neutre. Il correspond à l'énergie indispensable au maintien des \textbf{fonctions vitales} : activité du cœur et des muscles respiratoires, fonctionnement du cerveau, des reins et du foie, maintien de la température corporelle. Il représente environ \textbf{60 à 70 \% de la dépense énergétique totale} chez une personne sédentaire.
\end{definition}

\begin{aretenir}[Ordres de grandeur]
Chez un adulte sédentaire, le métabolisme de base représente environ \textbf{65 \%} de la dépense totale, l'activité physique environ \textbf{25 \%}, et la digestion + thermorégulation environ \textbf{10 \%}. Le métabolisme de base est plus élevé chez l'homme que chez la femme (à masse égale), plus élevé chez le jeune en croissance, et il diminue avec l'âge.
\end{aretenir}

\begin{popfigure}
\begin{center}
\begin{tikzpicture}[font=\small]
% Camembert tracé manuellement (rayon 2 cm)
\def\R{2}
% 1. Métabolisme de base 65% : 0° à 234°
\filldraw[fill=popblue!70, draw=popdark, thick] (0,0) -- (0:\R) arc (0:234:\R) -- cycle;
% 2. Activité physique 25% : 234° à 324°
\filldraw[fill=poporange!75, draw=popdark, thick] (0,0) -- (234:\R) arc (234:324:\R) -- cycle;
% 3. Digestion + Thermorégulation 10% : 324° à 360°
\filldraw[fill=popgreen!70, draw=popdark, thick] (0,0) -- (324:\R) arc (324:360:\R) -- cycle;

% Légende à droite du camembert
\filldraw[fill=popblue!70, draw=popblue] (3, 1.5) rectangle (3.4, 1.9);
\node[anchor=west] at (3.5, 1.7) {Métabolisme de base (65 \%)};
\filldraw[fill=poporange!75, draw=poporange] (3, 0.8) rectangle (3.4, 1.2);
\node[anchor=west] at (3.5, 1.0) {Activité physique (25 \%)};
\filldraw[fill=popgreen!70, draw=popgreen] (3, 0.1) rectangle (3.4, 0.5);
\node[anchor=west] at (3.5, 0.3) {Digestion et thermorégulation (10 \%)};
\end{tikzpicture}
\end{center}
\legende{Répartition approximative des dépenses énergétiques quotidiennes chez une personne sédentaire. Le métabolisme de base, invisible mais permanent, constitue de loin la plus grande part : même immobile toute la journée, l'organisme dépense l'essentiel de son énergie.}
\end{popfigure}

\begin{attention}[Le piège classique du calcul des dépenses]
L'erreur la plus fréquente au Probatoire est d'\textbf{oublier le métabolisme de base} et de ne compter que l'exercice physique dans les dépenses. Un élève qui écrit « dépenses = 30 min de sport = 800 kJ » commet une erreur majeure : même sans aucun sport, l'organisme dépense chaque jour $6\,000$ à $8\,000\ \text{kJ}$ rien que pour ses fonctions vitales. \textbf{Toujours} additionner : métabolisme de base + thermorégulation + activité physique + digestion.
\end{attention}

\courssub{Les trois cas possibles du bilan énergétique}

Comme toute différence de deux grandeurs, le bilan énergétique peut être nul, positif ou négatif :

\begin{center}
\begin{tabularx}{\linewidth}{|l|X|X|}
\hline
\rowcolor{popblueL} \textbf{Cas} & \textbf{Relation} & \textbf{Conséquence immédiate} \\
\hline
Bilan \textbf{équilibré} & Apports $=$ Dépenses & Masse corporelle stable : les réserves sont maintenues. \\
\hline
Bilan \textbf{positif} & Apports $>$ Dépenses & L'excédent est \textbf{stocké}, principalement sous forme de graisses (triglycérides) dans le tissu adipeux $\rightarrow$ prise de masse grasse. \\
\hline
Bilan \textbf{négatif} & Apports $<$ Dépenses & L'organisme \textbf{puise dans ses réserves} (glycogène puis graisses, voire protéines musculaires) $\rightarrow$ amaigrissement, puis épuisement si le déficit persiste. \\
\hline
\end{tabularx}
\end{center}

\begin{popfigure}
\begin{center}
\begin{tikzpicture}[scale=1, every node/.style={font=\small}]
% --- Cas 1 : équilibre ---
\begin{scope}[shift={(-5,0)}]
\draw[popdark, thick] (0,0) -- (0,1.4);
\draw[popdark, thick] (-1.6,1.4) -- (1.6,1.4);
\draw[popblue, thick] (-1.6,1.4) -- (-1.6,0.9);
\draw[poporange, thick] (1.6,1.4) -- (1.6,0.9);
\draw[fill=popblueL, draw=popblue] (-2.1,0.9) rectangle (-1.1,0.55);
\draw[fill=poporangeL, draw=poporange] (1.1,0.9) rectangle (2.1,0.55);
\node[below] at (-1.6,0.5) {\textbf{Apports}};
\node[below] at (1.6,0.5) {\textbf{Dépenses}};
\node[below, align=center] at (0,-0.35) {\textbf{Bilan équilibré}\\ masse stable};
\end{scope}
% --- Cas 2 : positif ---
\begin{scope}[shift={(0,0)}]
\draw[popdark, thick] (0,0) -- (0,1.4);
\draw[popdark, thick, rotate around={-12:(0,1.4)}] (-1.6,1.4) -- (1.6,1.4);
\draw[popblue, thick] (-1.6,1.73) -- (-1.6,1.23);
\draw[poporange, thick] (1.6,1.07) -- (1.6,0.57);
\draw[fill=popblueL, draw=popblue] (-2.1,1.23) rectangle (-1.1,0.88);
\draw[fill=poporangeL, draw=poporange] (1.1,0.57) rectangle (2.1,0.22);
\node[below] at (-1.9,0.83) {\textbf{Apports}};
\node[below] at (1.9,0.17) {\textbf{Dépenses}};
\node[below, align=center] at (0,-0.35) {\textbf{Bilan positif}\\ stockage des graisses};
\end{scope}
% --- Cas 3 : négatif ---
\begin{scope}[shift={(5,0)}]
\draw[popdark, thick] (0,0) -- (0,1.4);
\draw[popdark, thick, rotate around={12:(0,1.4)}] (-1.6,1.4) -- (1.6,1.4);
\draw[popblue, thick] (-1.6,1.07) -- (-1.6,0.57);
\draw[poporange, thick] (1.6,1.73) -- (1.6,1.23);
\draw[fill=popblueL, draw=popblue] (-2.1,0.57) rectangle (-1.1,0.22);
\draw[fill=poporangeL, draw=poporange] (1.1,1.23) rectangle (2.1,0.88);
\node[below] at (-1.9,0.17) {\textbf{Apports}};
\node[below] at (1.9,0.83) {\textbf{Dépenses}};
\node[below, align=center] at (0,-0.35) {\textbf{Bilan négatif}\\ utilisation des réserves};
\end{scope}
\end{tikzpicture}
\end{center}
\legende{La balance énergétique de l'organisme. À gauche, apports égaux aux dépenses : la balance est horizontale, la masse corporelle est stable. Au centre, les apports l'emportent : l'excédent est stocké dans le tissu adipeux. À droite, les dépenses l'emportent : l'organisme mobilise puis épuise ses réserves.}
\end{popfigure}

\begin{propriete}[Devenir de l'excédent énergétique]
Lorsque le bilan est positif, l'excédent de glucose est d'abord stocké sous forme de \cle{glycogène} dans le foie et les muscles (réserve limitée, environ $400\ \text{g}$), puis l'excédent restant — ainsi que les lipides alimentaires — est transformé en \cle{triglycérides} et stocké dans les \textbf{adipocytes} du tissu adipeux, réserve pratiquement illimitée. C'est ce stockage chronique qui conduit au surpoids puis à l'obésité (section suivante).
\end{propriete}

\courssub{Méthode : établir un bilan énergétique simple}

\begin{methode}[Calculer un bilan énergétique en trois étapes]
\begin{enumerate}
\item \textbf{Calculer les apports totaux} : multiplier la masse de chaque nutriment consommé par sa valeur énergétique ($17\ \text{kJ/g}$ pour glucides et protides, $38\ \text{kJ/g}$ pour les lipides), puis additionner. Ne pas oublier les boissons énergétiques.
\item \textbf{Estimer les dépenses totales} : additionner le métabolisme de base (donné ou estimé), la thermorégulation, l'activité physique et la digestion. Si l'énoncé donne une répartition en pourcentages, l'utiliser.
\item \textbf{Soustraire et conclure} : Bilan $=$ Apports $-$ Dépenses. Si le résultat est $> 0$ : bilan positif (stockage). Si $< 0$ : bilan négatif (utilisation des réserves). Si $\approx 0$ : équilibre.
\end{enumerate}
\end{methode}

\begin{exoresolu}[Bilan énergétique d'un employé de bureau]
Énoncé : M. Etoundi, employé de bureau à Douala, consomme chaque jour des aliments et boissons totalisant $11\,500\ \text{kJ}$. Sa dépense énergétique totale (métabolisme de base + thermorégulation + activités + digestion) est estimée à $10\,000\ \text{kJ/jour}$. \textbf{1)} Calculer son bilan énergétique journalier. \textbf{2)} Conclure sur l'évolution prévisible de sa masse corporelle.
\tcblower
Solution :
\textbf{1)} On applique la formule du bilan :
\[ \text{Bilan} = \text{Apports} - \text{Dépenses} = 11\,500 - 10\,000 = +1\,500\ \text{kJ/jour}. \]
\textbf{2)} Le bilan est \textbf{positif} : les apports dépassent les dépenses de $1\,500\ \text{kJ}$ chaque jour. Cet excédent est stocké sous forme de triglycérides dans le tissu adipeux. Sachant que $1\ \text{kg}$ de graisse corporelle correspond à environ $30\,000\ \text{kJ}$, un tel excédent maintenu pendant un mois ($30 \times 1\,500 = 45\,000\ \text{kJ}$) correspond théoriquement à une prise d'environ $1{,}5\ \text{kg}$ de masse grasse par mois. M. Etoundi s'expose donc, à terme, au \textbf{surpoids} puis à l'\textbf{obésité}, sauf s'il réduit ses apports ou augmente ses dépenses (activité physique).
\end{exoresolu}

\begin{exoresolu}[Un bilan négatif]
Énoncé : Pendant la saison des travaux champêtres, une cultivatrice de la région du Nord dépense $13\,000\ \text{kJ/jour}$ mais ne consomme que $10\,500\ \text{kJ/jour}$ (période de soudure, avant les récoltes). Calculer son bilan et prévoir les conséquences si cette situation dure.
\tcblower
Solution :
\[ \text{Bilan} = 10\,500 - 13\,000 = -2\,500\ \text{kJ/jour}. \]
Le bilan est \textbf{négatif} : l'organisme comble le déficit en mobilisant d'abord le glycogène, puis les graisses de réserve. Si le déficit persiste plusieurs semaines, il y a \textbf{amaigrissement} ; s'il s'aggrave, l'organisme dégrade ses propres protéines musculaires, conduisant à la \textbf{dénutrition} et à l'\textbf{épuisement} (fonte musculaire, fatigue chronique, baisse des défenses immunitaires). Cette situation, fréquente en période de soudure dans les zones rurales, justifie les programmes d'aide alimentaire saisonniers.
\end{exoresolu}

\begin{savaistu}[Mesurer les dépenses : la calorimétrie]
Comment connaît-on la dépense énergétique d'une personne ? Par \textbf{calorimétrie directe} (on mesure la chaleur dégagée par le corps dans une chambre calorimétrique) ou, plus souvent, par \textbf{calorimétrie indirecte} : on mesure le volume de dioxygène consommé, car l'oxydation des nutriments est proportionnelle à l'énergie dépensée (environ $20\ \text{kJ}$ par litre de $\ce{O2}$ consommé). C'est ainsi qu'ont été établies les tables de dépense énergétique par activité : dormir coûte environ $4\ \text{kJ/min}$, marcher $15$ à $20\ \text{kJ/min}$, et un travail agricole intense ou un match de football peut dépasser $40\ \text{kJ/min}$.
\end{savaistu}

\begin{aretenir}[L'essentiel de la section]
\begin{itemize}
\item Le \cle{bilan énergétique} $=$ Apports alimentaires $-$ Dépenses totales (en kJ/jour).
\item Les apports proviennent \textbf{exclusivement} des aliments et boissons énergétiques.
\item Les dépenses comprennent \textbf{quatre} composantes : métabolisme de base (60--70 \% du total), thermorégulation, activité physique, digestion.
\item Bilan $= 0$ : équilibre, masse stable. Bilan $> 0$ : stockage de graisses. Bilan $< 0$ : mobilisation puis épuisement des réserves.
\item Ne jamais oublier le métabolisme de base dans le calcul des dépenses.
\end{itemize}
\end{aretenir}

Maintenant que nous savons établir et interpréter un bilan énergétique, nous pouvons étudier en détail ses deux grands types de déséquilibre. Commençons par le cas où la balance penche du côté des apports : le \textbf{déséquilibre énergétique positif}, responsable du surpoids et de l'obésité.

\coursec{Le déséquilibre énergétique positif : surpoids et obésité}

Dans la section précédente, nous avons appris à établir un \cle{bilan énergétique} en comparant les \cle{apports} (énergie fournie par les aliments) aux \cle{dépenses} (métabolisme de base, activité physique, thermorégulation, effet thermique des aliments). Nous avons vu que lorsque apports et dépenses s'équilibrent, la masse corporelle reste stable. Mais que se passe-t-il lorsque, jour après jour, les apports dépassent durablement les dépenses ? C'est la situation que nous allons maintenant étudier : le \cle{déséquilibre énergétique positif}, dont les manifestations les plus visibles sont le surpoids puis l'obésité.

\courssub{Que devient l'énergie consommée en excès ?}

\textbf{Une observation de la vie courante.}\par

Dans un quartier de Douala, un employé de bureau mange chaque midi un copieux plat de ndolé accompagné de miondo, complété le soir par des beignets et une boisson sucrée. Il se déplace en moto-taxi, travaille assis toute la journée et ne pratique aucune activité sportive. En trois ans, sa masse est passée de 70 kg à 88 kg. D'où vient cette masse supplémentaire ? Où est-elle stockée dans l'organisme ?

\textbf{Hypothèse.}\par L'énergie alimentaire consommée mais non dépensée n'est pas éliminée : elle est \emph{stockée} dans l'organisme.

\textbf{Données et interprétation.}\par Les nutriments énergétiques (glucides, lipides) absorbés en excès sont transformés par le métabolisme en \cle{triglycérides} (graisses de réserve). Ces triglycérides s'accumulent dans des cellules spécialisées, les \cle{adipocytes}, qui constituent le \cle{tissu adipeux}, situé principalement sous la peau (hypoderme : abdomen, cuisses, hanches) et autour de certains organes (tissu adipeux viscéral). Un adipocyte peut voir son volume multiplié par plusieurs dizaines lorsqu'il se remplit de graisse (hypertrophie adipocytaire) ; au-delà d'un certain seuil, le nombre d'adipocytes peut aussi augmenter (hyperplasie). C'est le tissu adipeux tout entier qui s'épaissit, et la masse corporelle augmente.

\begin{aretenir}[Stockage de l'excès d'énergie]
Lorsque le bilan énergétique est \cle{positif} (apports $>$ dépenses) de façon prolongée, l'excès d'énergie est stocké sous forme de \cle{graisses} (triglycérides) dans les \cle{adipocytes} du \cle{tissu adipeux}. La masse corporelle augmente alors progressivement.
\end{aretenir}

\begin{popfigure}
\begin{center}
\begin{tikzpicture}[scale=1]
% Flèche entrée : excès de nutriments
\draw[->, very thick, poporange] (-4.5,0) -- (-2.2,0);
\node[above, align=center, text=popdark] at (-3.35,0.15) {\small Excès de nutriments\\ \small (glucides, lipides)};
% Transformation
\node[draw=popblue, rounded corners, fill=popblueL, align=center, text width=3.2cm] at (0,0) {\small Métabolisme : transformation en \textbf{triglycérides}};
\draw[->, very thick, poporange] (1.8,0) -- (3.4,0);
% Adipocytes
\begin{scope}[shift={(5.6,0)}]
  \draw[fill=poppinkL, draw=poppink, thick] (0,0) circle (1.5);
  \draw[fill=popgoldL, draw=popgold] (0,0) circle (1.1);
  \draw[fill=popdark] (0.95,0.75) circle (0.12);
  \draw[fill=poppinkL, draw=poppink, thick] (2.6,0.4) circle (0.9);
  \draw[fill=popgoldL, draw=popgold] (2.6,0.4) circle (0.6);
  \draw[fill=popdark] (3.15,0.85) circle (0.09);
  \draw[fill=poppinkL, draw=poppink, thick] (2.2,-1.1) circle (0.75);
  \draw[fill=popgoldL, draw=popgold] (2.2,-1.1) circle (0.45);
  \draw[fill=popdark] (2.6,-0.7) circle (0.08);
\end{scope}
% Légendes pointées
\draw[thin, popmuted] (5.6,-1.5) -- (5.6,-2.2);
\node[below, align=center] at (5.6,-2.2) {\small Adipocytes gorgés de graisse\\ \small (tissu adipeux)};
\draw[thin, popmuted] (6.55,0.75) -- (7.6,1.6);
\node[right, align=left] at (7.6,1.6) {\small noyau};
\draw[thin, popmuted] (6.7,0) -- (7.6,0.6);
\node[right, align=left] at (7.6,0.4) {\small gouttelette de\\ \small triglycérides};
\end{tikzpicture}
\end{center}
\legende{Schéma simplifié du stockage de l'énergie en excès : les nutriments non dépensés sont transformés en triglycérides qui s'accumulent dans les adipocytes, cellules du tissu adipeux, dont le volume augmente (hypertrophie).}
\end{popfigure}

\courssub{Du surpoids à l'obésité : une accumulation excessive de masse grasse}

L'augmentation de la masse grasse ne se fait pas brutalement : elle suit une progression continue. Tant que l'excédent énergétique persiste, le tissu adipeux s'épaissit. On distingue deux stades de cette accumulation excessive : le \cle{surpoids}, stade modéré, puis l'\cle{obésité}, stade avancé.

\begin{definition}[Obésité]
L'\cle{obésité} est une \textbf{accumulation excessive de graisses corporelles} résultant d'un \textbf{bilan énergétique positif prolongé}, avec des \textbf{conséquences néfastes sur la santé}. Le \cle{surpoids} est un stade intermédiaire, moins avancé, entre l'état normal et l'obésité.
\end{definition}

\textbf{Une évolution mesurable.}\par Un excédent de seulement 500 kJ par jour (l'équivalent d'un beignet et d'un verre de boisson sucrée) peut sembler négligeable. Pourtant, répété chaque jour, il représente environ \num{15000} kJ par mois, soit l'équivalent énergétique d'environ 0,5 kg de graisse corporelle stockée. Mois après mois, la masse corporelle s'élève régulièrement, comme le montre la courbe ci-dessous.

\begin{popfigure}
\begin{center}
\begin{tikzpicture}
\begin{axis}[
  width=0.85\linewidth, height=6.5cm,
  xlabel={Temps (mois)}, ylabel={Masse corporelle (kg)},
  xmin=0, xmax=24, ymin=68, ymax=84,
  xtick={0,3,6,9,12,15,18,21,24},
  ytick={70,72,74,76,78,80,82},
  grid=both, grid style={popline!40},
  axis line style={popdark},
]
\addplot[popcurve, domain=0:24, samples=100] {70 + 0.5*x};
\addplot[mark=*, only marks, poporange, mark size=1.5pt] coordinates {(0,70) (6,73) (12,76) (18,79) (24,82)};
\node[align=center, text=popdark, font=\small] at (axis cs:6,80.5) {Bilan énergétique positif\\ constant : $+500$ kJ/jour};
\draw[->, thin, popmuted] (axis cs:9,79.8) -- (axis cs:15,77.6);
\end{axis}
\end{tikzpicture}
\end{center}
\legende{Évolution de la masse corporelle d'une personne sédentaire dont le bilan énergétique reste constamment positif ($+500$ kJ/jour, soit environ $0,5$ kg de graisse stockée par mois) : la masse passe de 70 kg à 82 kg en deux ans.}
\end{popfigure}

\courssub{Un indicateur simple : l'Indice de Masse Corporelle (IMC)}

Comment savoir objectivement si une personne est en surpoids ou obèse ? La seule masse ne suffit pas : 80 kg n'a pas la même signification pour une personne mesurant 1,60 m ou 1,90 m. Les médecins utilisent donc un indicateur qui rapporte la masse à la taille : l'\cle{Indice de Masse Corporelle} (IMC).

\begin{methode}[Calculer et interpréter l'IMC]
\begin{enumerate}
  \item Mesurer la \textbf{masse} $m$ de la personne en kilogrammes (kg).
  \item Mesurer sa \textbf{taille} $t$ en mètres (m), puis calculer le carré de la taille : $t^2$ (en m$^2$).
  \item Calculer l'IMC :
  \[ \mathrm{IMC} = \dfrac{m}{t^2} \qquad \text{(en kg/m}^2\text{)} \]
  \item Comparer le résultat aux valeurs de référence de l'OMS :
\end{enumerate}
\begin{center}
\begin{tabularx}{\linewidth}{|c|X|}\hline
\rowcolor{popblueL} \textbf{IMC (kg/m$^2$)} & \textbf{Interprétation} \\ \hline
$< 18,5$ & Maigreur (déficit énergétique probable) \\ \hline
$18,5$ à $25$ & Corpulence normale \\ \hline
$25$ à $30$ & \textbf{Surpoids} \\ \hline
$> 30$ & \textbf{Obésité} \\ \hline
\end{tabularx}
\end{center}
\end{methode}

\begin{exemplebox}[Calcul de l'IMC d'un employé de bureau]
Une personne a une masse $m = 80$ kg et une taille $t = 1,70$ m.
\[ \mathrm{IMC} = \dfrac{80}{(1,70)^2} = \dfrac{80}{2,89} \approx 27,7 \ \text{kg/m}^2 \]
Comme $25 < 27,7 < 30$, cette personne est en situation de \textbf{surpoids}. Elle n'est pas encore obèse, mais elle doit réagir avant que son IMC ne dépasse 30.
\end{exemplebox}

\begin{attention}[Ne pas confondre surpoids et obésité]
Erreur fréquente au Probatoire : employer les deux termes comme des synonymes. Le \cle{surpoids} correspond à un IMC compris entre 25 et 30 ; l'\cle{obésité} est un stade \textbf{plus avancé} de l'accumulation de graisse, correspondant à un IMC \textbf{supérieur à 30}. Autre piège : oublier d'élever la taille \emph{au carré} dans le calcul de l'IMC, ou exprimer la taille en centimètres au lieu de mètres.
\end{attention}

\courssub{Les causes fréquentes du déséquilibre énergétique positif}

Pourquoi une personne consomme-t-elle durablement plus d'énergie qu'elle n'en dépense ? Trois causes se combinent le plus souvent :

\begin{itemize}
  \item \textbf{Une alimentation trop riche en graisses et en sucres.} Les lipides sont les nutriments les plus énergétiques (environ 38 kJ/g, contre 17 kJ/g pour les glucides et les protéines). Les fritures (beignets, puff-puff, frites de plantain), les huiles en abondance, les boissons sucrées et les friandises apportent beaucoup d'énergie en petit volume, sans procurer une satiété durable.
  \item \textbf{La sédentarité.} Travail assis, déplacements motorisés (voitures, motos-taxis), loisirs devant les écrans : les dépenses énergétiques quotidiennes diminuent fortement alors que les apports restent élevés.
  \item \textbf{L'urbanisation des modes de vie.} En ville, l'alimentation traditionnelle riche en tubercules, légumes et céréales complètes cède souvent la place à une restauration rapide, grasse et sucrée, tandis que les activités physiques liées aux travaux champêtres disparaissent. Le déséquilibre positif s'installe alors insidieusement.
\end{itemize}

\begin{savaistu}[Une épidémie mondiale, y compris en Afrique]
Selon l'Organisation mondiale de la santé (OMS), le nombre de personnes obèses a \textbf{presque triplé dans le monde depuis 1975}, et cette progression touche désormais l'Afrique subsaharienne. Dans les grandes villes camerounaises comme Yaoundé et Douala, l'obésité devient un \textbf{problème de santé publique croissant} : elle touche les adultes mais aussi, de plus en plus, les enfants et les adolescents, en parallèle de la dénutrition qui persiste dans certaines zones rurales. On parle de « double fardeau » nutritionnel.
\end{savaistu}

\courssub{Les conséquences physiologiques du surpoids et de l'obésité}

\textbf{Des conséquences immédiates et perceptibles.}

L'excès de masse grasse n'est pas inerte : il gêne le fonctionnement de l'organisme au quotidien.

\begin{itemize}
  \item \textbf{Essoufflement à l'effort.} Le corps, plus lourd, exige un travail musculaire et cardiaque accru pour le moindre déplacement ; de plus, la graisse abdominale comprime la cage thoracique et limite l'amplitude respiratoire. Monter quelques étages suffit à essouffler.
  \item \textbf{Fatigue des articulations.} Les articulations des genoux, des hanches et de la colonne vertébrale supportent une charge excessive à chaque pas. À la longue, les cartilages s'usent prématurément, provoquant douleurs et difficultés de déplacement (arthrose).
  \item \textbf{Hypertension artérielle.} Le cœur doit pomper le sang avec plus de force pour irriguer une masse corporelle augmentée ; la pression artérielle s'élève durablement, fatiguant le cœur et les parois des vaisseaux.
\end{itemize}

\textbf{Des conséquences à long terme, plus graves.}

Lorsque l'obésité s'installe pendant des années, elle favorise l'apparition de maladies chroniques :

\begin{itemize}
  \item le \cle{diabète de type 2} : l'organisme devient moins sensible à l'insuline et la glycémie (taux de glucose sanguin) reste durablement trop élevée ;
  \item les \cle{maladies cardiovasculaires} : dépôts de graisses sur les parois des artères (athérosclérose), accidents vasculaires cérébraux, infarctus du myocarde.
\end{itemize}

\begin{attention}[Limite du programme]
Les mécanismes précis du diabète de type 2 (action de l'insuline, régulation de la glycémie) et des maladies cardiovasculaires (formation des plaques d'athérome) ne sont \textbf{pas exigibles} en Première C/TI : ils seront étudiés dans les classes ultérieures. Retiens simplement que l'obésité est un \emph{facteur de risque majeur} de ces maladies.
\end{attention}

\courssub{Exercice résolu et bilan de la section}

\begin{exoresolu}[Diagnostic d'un déséquilibre énergétique]
\textbf{Énoncé.} M. Etoundi, commerçant à Yaoundé, mesure 1,75 m et pèse 96 kg. Ses apports alimentaires quotidiens sont estimés à \num{12500} kJ, alors que ses dépenses totales (métabolisme de base + activité) n'atteignent que \num{10500} kJ.
\begin{enumerate}
  \item Calcule son IMC et interprète le résultat.
  \item Calcule son excédent énergétique quotidien, puis estime sa prise de masse sur un mois (30 jours), sachant que 1 kg de graisse corporelle correspond à environ \num{30000} kJ.
  \item Cite deux conséquences physiologiques qu'il risque de ressentir.
\end{enumerate}
\tcblower
\textbf{Solution détaillée.}
\begin{enumerate}
  \item $\mathrm{IMC} = \dfrac{96}{(1,75)^2} = \dfrac{96}{3,0625} \approx 31,3$ kg/m$^2$. Comme $\mathrm{IMC} > 30$, M. Etoundi est en situation d'\textbf{obésité}.
  \item Excédent quotidien : $\num{12500} - \num{10500} = \num{2000}$ kJ/jour. Sur 30 jours : $\num{2000} \times 30 = \num{60000}$ kJ, soit $\dfrac{\num{60000}}{\num{30000}} = 2$ kg de graisse stockés en un mois si rien ne change.
  \item Il risque notamment un \textbf{essoufflement à l'effort}, une \textbf{fatigue des articulations} (genoux, hanches) et une \textbf{hypertension artérielle}.
\end{enumerate}
\end{exoresolu}

\begin{exercice}[À toi de jouer]
Mme Mbappe mesure 1,60 m et pèse 72 kg. Calcule son IMC, interprète-le, puis propose une hypothèse expliquant comment elle a pu atteindre cette corpulence. Quelle conséquence à long terme doit-elle chercher à éviter en priorité ?
\end{exercice}

\begin{corrige}[Exercice 1]
$\mathrm{IMC} = \dfrac{72}{(1,60)^2} = \dfrac{72}{2,56} \approx 28,1$ kg/m$^2$ : Mme Mbappe est en \textbf{surpoids} ($25 < 28,1 < 30$), pas encore obèse. Hypothèse : un bilan énergétique positif prolongé, par exemple une alimentation riche en graisses et sucres associée à la sédentarité. Conséquence à éviter en priorité : la progression vers l'obésité et ses complications à long terme, notamment le \textbf{diabète de type 2} et les \textbf{maladies cardiovasculaires}.
\end{corrige}

\begin{aretenir}[L'essentiel de la section]
\begin{itemize}
  \item Un bilan énergétique \textbf{positif prolongé} entraîne le stockage de l'excès d'énergie sous forme de \textbf{graisses} dans le \textbf{tissu adipeux}.
  \item L'accumulation excessive de masse grasse conduit d'abord au \textbf{surpoids} ($25 < \mathrm{IMC} < 30$), puis à l'\textbf{obésité} ($\mathrm{IMC} > 30$).
  \item $\mathrm{IMC} = \dfrac{\text{masse (kg)}}{\text{taille}^2\ (\text{m}^2)}$ ; maigreur : $< 18,5$ ; normal : $18,5$ à $25$.
  \item Causes fréquentes : alimentation trop riche en graisses et sucres, sédentarité, urbanisation des modes de vie.
  \item Conséquences : essoufflement à l'effort, fatigue des articulations, hypertension artérielle ; à long terme, diabète de type 2 et maladies cardiovasculaires.
\end{itemize}
\end{aretenir}

Heureusement, ce déséquilibre n'est pas une fatalité : l'organisme peut revenir à l'équilibre. Mais avant d'étudier les mesures de rééquilibrage, nous examinerons dans la section suivante le déséquilibre inverse — le \cle{déséquilibre énergétique négatif} — qui conduit à la dénutrition et à l'épuisement.

\coursec{Le déséquilibre énergétique négatif : dénutrition et épuisement}

Dans la section précédente, nous avons vu ce qui se passe lorsque l'organisme reçoit \emph{plus} d'énergie qu'il n'en dépense : l'excédent est stocké sous forme de graisses, conduisant au surpoids puis à l'obésité. Mais la balance énergétique peut pencher de l'autre côté. Que se passe-t-il lorsque les apports sont durablement \emph{inférieurs} aux dépenses ? L'organisme, pour survivre, doit alors puiser dans ses propres réserves — et ce mécanisme de survie, s'il se prolonge, devient destructeur. C'est tout l'objet de cette section : comprendre le \cle{déséquilibre énergétique négatif} et ses conséquences, de la simple fatigue à la dénutrition sévère.

\courssub{Situation de départ : le cas d'Amadou, portefaix au marché Mokolo}

Observons une situation concrète. Amadou, 25 ans, travaille comme portefaix au marché Mokolo à Yaoundé. Son activité physique intense lui fait dépenser environ \SI{13000}{\kilo\joule} par jour. Or, pour des raisons économiques, ses repas quotidiens (un peu de bâton de manioc le matin, un plat de riz le soir) ne lui apportent que \SI{9000}{\kilo\joule}. Son bilan énergétique quotidien est donc :

\[ \text{Bilan} = \text{Apports} - \text{Dépenses} = \num{9000} - \num{13000} = -\SI{4000}{\kilo\joule/jour} \]

Après trois mois, Amadou a perdu 8 kg, se sent constamment fatigué, attrape des rhumes à répétition et peine à soulever les charges qu'il portait facilement auparavant.

\textbf{Question-problème :} d'où l'organisme d'Amadou tire-t-il les \SI{4000}{\kilo\joule} manquants chaque jour, et pourquoi son état se dégrade-t-il progressivement ?

\courssub{La mobilisation des réserves : un ordre précis}

\textbf{Hypothèse.} L'organisme doit posséder des réserves d'énergie qu'il utilise en cas de déficit. Mais les utilise-t-il au hasard, ou selon un ordre déterminé ?

\textbf{Données et observations.} Le suivi médical de personnes en jeûne prolongé (ou en grève de la faim encadrée médicalement) montre une succession de phases bien caractérisées :

\begin{itemize}
\item \textbf{Premières heures (0 à 24 h) :} la glycémie est maintenue grâce à la dégradation du \cle{glycogène} hépatique et musculaire (réserve de glucose, environ 400 à 500 g chez l'adulte, soit à peine une journée d'énergie).
\item \textbf{Premiers jours et semaines :} ce sont les \cle{graisses} (triglycérides du tissu adipeux) qui deviennent la principale source d'énergie. L'amaigrissement devient visible.
\item \textbf{Jeûne prolongé (au-delà de plusieurs semaines) :} lorsque les graisses s'épuisent, l'organisme dégrade ses propres \cle{protéines}, notamment les protéines musculaires, pour produire du glucose (néoglucogenèse). C'est la \cle{fonte musculaire}.
\end{itemize}

\textbf{Interprétation.} L'organisme utilise ses réserves dans un ordre précis, du plus « économique » au plus « coûteux » : d'abord le glycogène (réserve rapide mais limitée), puis les graisses (réserve abondante et très énergétique), enfin les protéines musculaires (tissus fonctionnels dont la destruction compromet la survie).

\begin{aretenir}[Ordre d'utilisation des réserves énergétiques]
En cas de bilan énergétique négatif prolongé, l'organisme puise dans ses réserves selon l'ordre suivant :
\begin{enumerate}
\item le \cle{glycogène} (foie et muscles) : épuisé en 24 à 48 heures ;
\item les \cle{graisses} (tissu adipeux) : mobilisées pendant des semaines, c'est la phase d'amaigrissement ;
\item les \cle{protéines musculaires} : utilisées en dernier recours, c'est la phase d'épuisement et de fonte musculaire.
\end{enumerate}
\end{aretenir}

\begin{popfigure}
\begin{center}
\begin{tikzpicture}[scale=1, every node/.style={font=\small}]
% Marches de l'escalier
\fill[popgreenL] (0,0) rectangle (3,1);
\fill[popgoldL] (3,1) rectangle (6,2);
\fill[poppinkL] (6,2) rectangle (9,3);
% Contour de l'escalier
\draw[popdark, thick] (0,0) -- (0,1) -- (3,1) -- (3,2) -- (6,2) -- (6,3) -- (9,3) -- (9,0) -- cycle;
% Textes dans les marches
\node[align=center, popdark] at (1.5,0.5) {\textbf{1. Glycogène}\\ (24--48 h)};
\node[align=center, popdark] at (4.5,1.5) {\textbf{2. Graisses}\\ (semaines)};
\node[align=center, popdark] at (7.5,2.5) {\textbf{3. Protéines}\\ \textbf{musculaires}\\ (dernier recours)};
% Axe temps
\draw[-{Stealth}, popdark, thick] (0,-0.5) -- (9.5,-0.5) node[below left, popdark] {durée du déficit énergétique};
% Flèches d'annotation
\draw[-{Stealth}, popmuted] (1.5,1.15) -- (1.5,1.6) node[above, popmuted] {réserve rapide};
\draw[-{Stealth}, popmuted] (4.5,2.15) -- (4.5,2.6) node[above, popmuted] {amaigrissement};
\draw[-{Stealth}, popmuted] (7.5,3.15) -- (7.5,3.6) node[above, popmuted] {fonte musculaire, danger vital};
\end{tikzpicture}
\end{center}
\legende{Schéma en escalier de l'ordre d'utilisation des réserves énergétiques lors d'un bilan énergétique négatif prolongé : le glycogène est mobilisé en premier, puis les graisses, enfin les protéines musculaires.}
\end{popfigure}

\begin{savaistu}[Pourquoi les protéines sont-elles utilisées en dernier recours ?]
Les protéines ne sont pas une réserve de stockage : ce sont des molécules \emph{fonctionnelles}. Elles constituent les muscles, les enzymes, les anticorps, l'hémoglobine. Les protéines musculaires ne sont utilisées qu'en dernier recours, car leur destruction compromet le fonctionnement d'organes vitaux : le \textbf{c\oe ur} (qui est un muscle) et les \textbf{muscles respiratoires} (diaphragme, muscles intercostaux). C'est pourquoi la dénutrition avancée peut entraîner la mort par défaillance cardiaque ou respiratoire.
\end{savaistu}

\courssub{La dénutrition : définition et mécanisme}

\begin{definition}[Dénutrition]
La \cle{dénutrition} est un état pathologique résultant d'apports énergétiques insuffisants par rapport aux besoins de l'organisme, entraînant l'épuisement progressif des réserves (glycogène, graisses, puis protéines) et une perte de masse corporelle.
\end{definition}

Le mécanisme est une simple conséquence du bilan énergétique négatif prolongé : chaque jour de déficit oblige l'organisme à « brûler » une partie de lui-même. Tant que ce sont les graisses qui fondent, l'amaigrissement reste réversible ; mais lorsque les protéines musculaires sont attaquées, les fonctions vitales elles-mêmes sont menacées : c'est l'\cle{épuisement}.

\begin{exemplebox}[Un déficit chiffré : le cas d'Amadou]
Reprenons notre portefaix du marché Mokolo : il dépense \SI{13000}{\kilo\joule/jour} mais n'absorbe que \SI{9000}{\kilo\joule/jour}, soit un déficit de \SI{4000}{\kilo\joule/jour}. Sachant qu'un gramme de graisse corporelle libère environ \SI{38}{\kilo\joule} (valeur arrondie à \SI{40}{\kilo\joule} pour les calculs simplifiés), ce déficit correspond à la mobilisation d'environ \textbf{105 g de graisse corporelle chaque jour} ($4000 / 38 \approx 105$). En un mois, cela représente environ 3,2 kg de masse perdue — un rythme dangereux s'il se prolonge.
\end{exemplebox}

\begin{attention}[Dénutrition et malnutrition : ne pas confondre]
La \textbf{dénutrition} désigne une insuffisance \emph{quantitative} d'énergie : l'organisme ne reçoit pas assez de « carburant ». La \textbf{malnutrition} est un terme plus large qui englobe tout déséquilibre nutritionnel, quantitatif \emph{ou qualitatif} : carence en vitamines, en sels minéraux ou en protéines, mais aussi excès alimentaire (l'obésité est une forme de malnutrition !). La dénutrition est donc un cas particulier de malnutrition. Au Probatoire, employer « malnutrition » à la place de « dénutrition » pour décrire un déficit énergétique est une imprécision sanctionnée.
\end{attention}

\courssub{Les formes sévères de dénutrition chez l'enfant}

Chez le jeune enfant, dont les besoins de croissance sont énormes, la dénutrition prend des formes dramatiques que les agents de santé des centres de nutrition du Cameroun connaissent bien (mention indicative, non exigible en détail) :

\begin{itemize}
\item le \cle{kwashiorkor} : dénutrition par \textbf{carence protéique} prédominante. L'enfant reçoit parfois assez de calories (bouillies de céréales) mais pas assez de protéines. Signes caractéristiques : ventre ballonné (ascite), \oe dèmes des membres, cheveux décolorés et fragiles, apathie ;
\item le \cle{marasme} : dénutrition par \textbf{carence énergétique globale} (manque de tout). L'enfant est extrêmement maigre, « peau sur os », visage de vieillard, mais sans \oe dèmes.
\end{itemize}

Ces deux formes sévères exigent une prise en charge nutritionnelle urgente (laits thérapeutiques, aliments prêts à l'emploi) dans des structures spécialisées.

\courssub{Les conséquences physiologiques du déséquilibre négatif}

Lorsque le bilan énergétique reste négatif pendant des semaines ou des mois, les conséquences s'enchaînent :

\begin{itemize}
\item \textbf{amaigrissement} : la fonte des graisses réduit la masse corporelle de façon visible ;
\item \textbf{fonte musculaire} : la dégradation des protéines musculaires entraîne une perte de force (Amadou ne soulève plus ses charges) ;
\item \textbf{fatigue chronique} : l'organisme, privé de carburant, réduit ses dépenses non essentielles ; la moindre activité devient épuisante ;
\item \textbf{baisse des défenses immunitaires} : la fabrication des anticorps (qui sont des protéines) et des cellules immunitaires ralentit ; les infections (rhumes, paludisme plus sévère, infections digestives) deviennent plus fréquentes et plus graves ;
\item \textbf{épuisement} : stade ultime où les réserves sont épuisées et les organes vitaux (c\oe ur, muscles respiratoires) sont atteints — le pronostic vital est engagé.
\end{itemize}

\begin{attention}[Erreur fréquente : « maigrir, c'est toujours bon pour la santé »]
Contrairement à une idée répandue, maigrir n'est \textbf{pas} toujours bénéfique. Une perte de masse \emph{volontaire, lente et encadrée} chez une personne en surpoids est profitable ; mais une perte de masse \textbf{rapide et involontaire} est un signal d'alarme. Ordre de grandeur à retenir : une perte de plus de \textbf{10 \% de la masse corporelle en 3 à 6 mois} est un signe d'alerte qui impose une consultation médicale. Elle peut révéler une dénutrition, mais aussi une maladie sous-jacente (tuberculose, diabète, cancer, infection chronique).
\end{attention}

\begin{popfigure}
\begin{center}
\begin{tikzpicture}
\begin{axis}[
  width=0.9\linewidth, height=6.5cm,
  xlabel={Durée du déficit énergétique (mois)},
  ylabel={Masse corporelle (kg)},
  xmin=0, xmax=12, ymin=40, ymax=75,
  xtick={0,2,4,6,8,10,12},
  ytick={40,45,50,55,60,65,70},
  grid=both, grid style={popline!40},
  axis line style={popdark},
  label style={popdark},
  tick label style={popdark, font=\footnotesize},
  legend style={font=\footnotesize, at={(0.02,0.02)}, anchor=south west}
]
% Courbe : perte rapide (graisses) puis ralentie (adaptation) puis chute (protéines)
\addplot[popcurve, domain=0:12, samples=200] {70 - 14*(1-exp(-x/3)) - 0.9*x};
\addlegendentry{masse corporelle}
% Seuil d'alerte : -10 % de 70 kg = 63 kg
\addplot[poporange, dashed, thick, domain=0:12] {63};
\addlegendentry{seuil d'alerte ($-10\,\%$)}
\node[poporange, font=\footnotesize, anchor=south west] at (axis cs:0.2,63.3) {63 kg};
% Annotations de phases
\node[popmuted, font=\footnotesize, align=center] at (axis cs:2,71.5) {phase 1 : glycogène\\ puis graisses};
\node[popmuted, font=\footnotesize, align=center] at (axis cs:10.5,47) {phase 3 : protéines\\ musculaires};
\end{axis}
\end{tikzpicture}
\end{center}
\legende{Évolution de la masse corporelle d'un adulte de 70 kg soumis à un bilan énergétique négatif prolongé. La perte est rapide au début (mobilisation du glycogène puis des graisses), franchit le seuil d'alerte de $-10\,\%$ (63 kg) vers le quatrième mois, puis s'accentue lorsque les protéines musculaires sont dégradées.}
\end{popfigure}

\courssub{Conséquences chez l'adolescent : un enjeu scolaire et sportif}

Chez l'adolescent — c'est-à-dire toi, élève de Première — les besoins énergétiques sont particulièrement élevés à cause de la \textbf{croissance} (pic de croissance pubertaire) et des activités scolaires et sportives. Un bilan énergétique négatif à cet âge a des conséquences spécifiques :

\begin{itemize}
\item \textbf{retards de croissance} : la taille et le développement pubertaire peuvent être freinés, parfois de façon irréversible si le déficit dure ;
\item \textbf{difficultés de concentration} : le cerveau consomme à lui seul environ 20 \% de l'énergie de repos ; en déficit, l'attention, la mémorisation et le raisonnement sont altérés ;
\item \textbf{baisse des performances scolaires} : fatigue en classe, somnolence, absentéisme lié aux infections ;
\item \textbf{baisse des performances sportives} : fonte musculaire, essoufflement rapide, temps de récupération allongé.
\end{itemize}

Sauter le petit-déjeuner avant six heures de cours, ou s'entraîner intensivement au football sans augmenter ses rations, sont des situations banales qui, répétées, installent un déficit énergétique chez l'adolescent.

\courssub{Les causes du déséquilibre énergétique négatif}

Un bilan énergétique négatif prolongé peut avoir plusieurs origines, souvent combinées :

\begin{enumerate}
\item \textbf{l'insuffisance alimentaire} : pauvreté, précarité, mauvaise répartition des repas dans la journée, saison de soudure dans les zones rurales (période entre l'épuisement des greniers et la nouvelle récolte) ;
\item \textbf{les jeûnes prolongés} : jeûnes thérapeutiques mal encadrés, régimes amaigrissants drastiques, troubles du comportement alimentaire (anorexie) ;
\item \textbf{l'activité physique intense non compensée} : travailleurs de force (portefaix, manoeuvres, agriculteurs en pleine saison), sportifs qui augmentent l'entraînement sans augmenter les apports ;
\item \textbf{certaines maladies} : infections chroniques (tuberculose, VIH/Sida non traité), parasitoses digestives, cancers, maladies qui empêchent l'absorption des nutriments ou augmentent les dépenses (fièvres prolongées).
\end{enumerate}

\begin{exoresolu}[Bilan énergétique et signe d'alerte]
\textbf{Énoncé.} Étudiant(e) de 18 ans, Manga pèse 60 kg en janvier. Pendant la préparation d'une compétition d'athlétisme, il s'entraîne deux heures par jour tout en sautant régulièrement le petit-déjeuner et le repas de midi. En avril, il pèse 53 kg, se plaint de fatigue permanente et ses notes chutent.
\begin{enumerate}
\item Son bilan énergétique est-il positif, nul ou négatif ? Justifier.
\item Calculer le pourcentage de masse perdue et dire si le seuil d'alerte est franchi.
\item Citer deux conséquences physiologiques observables chez Manga et deux mesures de rééquilibrage.
\end{enumerate}
\tcblower
\textbf{Solution détaillée.}
\begin{enumerate}
\item Le bilan énergétique de Manga est \textbf{négatif} : ses dépenses (métabolisme de base + entraînement quotidien intense) dépassent durablement ses apports, réduits par la suppression de deux repas. Son organisme puise donc dans ses réserves, ce qui explique la perte de masse.
\item Pourcentage de masse perdue :
\[ \frac{60 - 53}{60} \times 100 = \frac{7}{60} \times 100 \approx 11,7\,\% \]
La perte dépasse 10 \% de la masse corporelle en quelques mois : le \textbf{seuil d'alerte est franchi}. Une consultation médicale s'impose.
\item Conséquences observables : fatigue chronique et baisse des performances scolaires (chute des notes), probable fonte musculaire et baisse des défenses immunitaires. Mesures de rééquilibrage : rétablir trois repas quotidiens équilibrés (augmenter les apports) et réduire temporairement l'intensité de l'entraînement (diminuer les dépenses), jusqu'au retour à un bilan équilibré.
\end{enumerate}
\end{exoresolu}

\begin{exercice}[À toi de jouer]
Une cultivatrice de 55 kg dépense \SI{11500}{\kilo\joule/jour} pendant la saison des semis mais ne consomme que \SI{8500}{\kilo\joule/jour} pendant deux mois.
\begin{enumerate}
\item Calculer son déficit énergétique quotidien.
\item En admettant que ce déficit soit couvert uniquement par la graisse corporelle (\SI{38}{\kilo\joule/g}), estimer la masse de graisse mobilisée en 60 jours, puis le pourcentage de masse corporelle perdu. Le seuil d'alerte est-il atteint ?
\item Expliquer pourquoi, si le déficit se prolonge au-delà de l'épuisement des graisses, son état devient dangereux pour sa santé.
\end{enumerate}
\end{exercice}

\begin{corrige}[Exercice : la cultivatrice]
\begin{enumerate}
\item Déficit quotidien : $\num{11500} - \num{8500} = \SI{3000}{\kilo\joule/jour}$.
\item Masse de graisse mobilisée par jour : $\num{3000}/38 \approx 79$ g. Sur 60 jours : $79 \times 60 = \num{4740}$ g $\approx 4,7$ kg. Pourcentage perdu : $4,7/55 \times 100 \approx 8,5\,\%$. Le seuil d'alerte de 10 \% n'est pas encore atteint, mais il en est proche : la situation doit être corrigée rapidement.
\item Une fois les graisses épuisées, l'organisme dégrade ses protéines musculaires. Or les muscles comprennent le c\oe ur et les muscles respiratoires : leur destruction compromet les fonctions vitales. S'y ajoutent la baisse des défenses immunitaires (anticorps = protéines) et la fatigue chronique : c'est l'épuisement, qui engage le pronostic vital.
\end{enumerate}
\end{corrige}

\begin{aretenir}[L'essentiel de la section]
\begin{itemize}
\item Un bilan énergétique négatif prolongé oblige l'organisme à mobiliser ses réserves dans l'ordre : \textbf{glycogène} $\rightarrow$ \textbf{graisses} $\rightarrow$ \textbf{protéines musculaires}.
\item La \textbf{dénutrition} est l'état pathologique résultant d'apports énergétiques insuffisants ; chez l'enfant, ses formes sévères sont le kwashiorkor (carence protéique) et le marasme (carence énergétique globale).
\item Conséquences : amaigrissement, fonte musculaire, fatigue chronique, baisse des défenses immunitaires, épuisement ; chez l'adolescent, retards de croissance et baisse des performances scolaires et sportives.
\item Signe d'alerte : perte de plus de \textbf{10 \% de la masse corporelle en 3 à 6 mois}.
\end{itemize}
\end{aretenir}

Reste à découvrir, dans la section suivante, comment l'activité physique — loin d'être seulement une dépense — joue un rôle central dans la \emph{régulation} du bilan énergétique.

\coursec{Rôle de l'activité physique dans la régulation du bilan énergétique}

Dans la section précédente, nous avons vu qu'un \cle{bilan énergétique négatif} prolongé conduit à la dénutrition, voire à l'épuisement de l'organisme : les dépenses dépassent durablement les apports, et le corps puise dans ses réserves (glycogène, lipides, puis protéines). Mais une question se pose alors naturellement : si les dépenses énergétiques sont si dangereuses quand elles deviennent excessives, pourquoi les médecins recommandent-ils à tous de \emph{bouger davantage} ? La réponse tient en une idée capitale : parmi tous les postes de dépense, l'\cle{activité physique} est le seul que l'on peut \textbf{moduler volontairement}, et c'est un formidable outil de régulation du bilan énergétique — à condition de l'utiliser intelligemment.

\courssub{Une observation de la vie courante : pourquoi certains « mangent sans grossir » ?}

Dans un même quartier de Yaoundé, deux adolescents du même âge ont des habitudes alimentaires comparables. Le premier pratique le football trois fois par semaine et marche quotidiennement pour se rendre au lycée ; le second passe ses journées assis, devant les cours puis la télévision. À terme, le second prend du poids, pas le premier. Pourtant, ni le cœur, ni le cerveau, ni la température corporelle du second ne « consomment moins » : leur métabolisme de base est voisin. La différence se situe ailleurs.

\begin{itemize}
\item \textbf{Observation :} à alimentation comparable, l'adolescent physiquement actif maintient un poids stable, le sédentaire prend du poids.
\item \textbf{Hypothèse :} l'activité physique augmente la dépense énergétique totale et rétablit l'équilibre entre apports et dépenses.
\item \textbf{Données :} on mesure les dépenses énergétiques horaires de différentes activités (voir tableau et figure ci-dessous) : elles varient dans un rapport de 1 à 10 environ.
\item \textbf{Interprétation :} le métabolisme de base étant quasiment identique chez les deux adolescents, c'est le poste « activité physique » qui fait la différence sur la dépense totale.
\item \textbf{Conclusion :} l'activité physique est un \emph{levier de réglage} du bilan énergétique : elle permet d'augmenter les dépenses sans toucher aux fonctions vitales.
\end{itemize}

\courssub{L'activité physique : le poste de dépense le plus modulable}

Rappelons les trois grands postes de la dépense énergétique totale : le \cle{métabolisme de base} (fonctions vitales : battements cardiaques, respiration, maintien de la température à \SI{37}{\celsius}, activité cérébrale), l'\cle{effet thermique des aliments} (thermogenèse post-prandiale liée à la digestion, l'absorption et le stockage des nutriments), et l'\cle{activité physique}. Or :

\begin{itemize}
\item on ne peut pas réduire durablement son métabolisme de base sans mettre en danger les fonctions vitales ;
\item on ne peut pas supprimer l'effet thermique des aliments sans cesser de s'alimenter ;
\item en revanche, on peut faire varier la dépense liée à l'activité physique \textbf{de quelques centaines à plusieurs milliers de kilojoules par jour}, simplement en choisissant de bouger.
\end{itemize}

\begin{definition}[Activité physique]
L'\cle{activité physique} désigne l'ensemble des mouvements corporels produits par la contraction des muscles squelettiques, entraînant une dépense énergétique supérieure à celle du repos. Elle comprend le sport, mais aussi la marche, les travaux champêtres, les tâches domestiques et les déplacements.
\end{definition}

Le tableau suivant rassemble les dépenses horaires indicatives de quelques activités courantes, pour un adolescent d'environ \SI{55}{kg} :

\begin{center}
\begin{tabularx}{\linewidth}{|l|X|}
\hline
\rowcolor{popblueL} \textbf{Activité} & \textbf{Dépense énergétique indicative} \\
\hline
Sommeil & $\approx$ \SI{250}{kJ/h} \\
\hline
Travail assis (cours, bureau) & $\approx$ \SI{400}{kJ/h} \\
\hline
Marche rapide & $\approx$ \SI{1200}{kJ/h} \\
\hline
Travaux champêtres (houe, plantation) & $\approx$ \SI{1500}{kJ/h} \\
\hline
Football & $\approx$ \SI{2000}{kJ/h} \\
\hline
Course à pied & $\approx$ \SI{2500}{kJ/h} \\
\hline
\end{tabularx}
\end{center}

\begin{popfigure}
\begin{center}
\begin{tikzpicture}
\begin{axis}[
  xbar,
  width=0.95\linewidth,
  height=6.5cm,
  xmin=0, xmax=2900,
  xlabel={Dépense énergétique (kJ/h)},
  symbolic y coords={Sommeil, Travail assis, Marche rapide, Travaux champêtres, Football, Course},
  ytick=data,
  xtick={0,500,1000,1500,2000,2500},
  bar width=9pt,
  nodes near coords,
  every node near coord/.append style={font=\small},
  xmajorgrids=true,
  grid style={popline!40},
]
\addplot[fill=popblue, draw=popdark] coordinates {
  (250,Sommeil) (400,Travail assis) (1200,Marche rapide)
  (1500,Travaux champêtres) (2000,Football) (2500,Course)
};
\end{axis}
\end{tikzpicture}
\end{center}
\legende{Dépense énergétique horaire indicative de différentes activités chez un adolescent. La course dépense environ dix fois plus que le sommeil : l'activité physique est le poste de dépense le plus modulable.}
\end{popfigure}

\begin{attention}[Lecture des valeurs]
Ces valeurs sont des \emph{ordres de grandeur} : la dépense réelle dépend de la masse corporelle (une personne plus lourde dépense davantage pour le même mouvement), de l'intensité réelle de l'effort et de sa durée. Au Probatoire, on utilise les valeurs fournies par l'énoncé, sans chercher à les « corriger ».
\end{attention}

\courssub{Estimer la dépense d'une activité}

\begin{methode}[Calculer une dépense énergétique liée à une activité]
\begin{enumerate}
\item Relever la \textbf{dépense horaire} de l'activité (donnée de l'énoncé ou du tableau).
\item Convertir la \textbf{durée} de l'activité en heures (par exemple $1\,\text{h}\,30 = 1{,}5\,\text{h}$).
\item \textbf{Multiplier} la dépense horaire par la durée :
\[ \text{Dépense (kJ)} = \text{dépense horaire (kJ/h)} \times \text{durée (h)} \]
\item Comparer le résultat aux apports alimentaires de la journée pour raisonner sur le bilan.
\end{enumerate}
\end{methode}

\begin{attention}[Modèle additif simplifié utilisé en Première]
Dans les exercices de bilan énergétique, on utilise couramment un modèle simplifié : on additionne la dépense des activités physiques à la dépense « incompressible » sur 24 h (métabolisme de base + effet thermique des aliments), sans retrancher le coût de base des heures d'activité. Ce modèle, bien qu'il surestime légèrement la dépense totale, est la convention attendue au Probatoire pour sa simplicité et sa sécurité (il ne sous-estime jamais les besoins).
\end{attention}

\begin{exemplebox}[Un match de football « vaut » un repas]
Un élève joue au football pendant $1\,\text{h}\,30$. Avec une dépense indicative de \SI{2000}{kJ/h} :
\[ 1{,}5 \times 2000 = \SI{3000}{kJ} \]
soit approximativement l'équivalent énergétique d'\textbf{un repas complet} (par exemple une portion de foutou avec sauce et poisson). Une seule séance de sport peut donc « effacer » l'équivalent d'un repas entier : cela montre la puissance du levier « activité physique ».
\end{exemplebox}

\begin{exoresolu}[Bilan d'une journée d'adolescent]
Énoncé. Steve, élève de Première, a des apports alimentaires de \SI{11500}{kJ} dans la journée. Son métabolisme de base et son effet thermique des aliments représentent \SI{7500}{kJ}. Ce jour-là, il marche rapidement pendant 30 minutes ($\approx$ \SI{1200}{kJ/h}) et joue au football pendant 1 heure ($\approx$ \SI{2000}{kJ/h}). (1) Calculer sa dépense énergétique totale. (2) Son bilan est-il équilibré, positif ou négatif ? (3) Que se passerait-il à terme s'il supprimait toute activité physique sans modifier son alimentation ?
\tcblower
Solution détaillée.
(1) Dépense liée à l'activité physique :
\begin{align*}
\text{Marche} &: 0{,}5\,\text{h} \times 1200\,\text{kJ/h} = 600\,\text{kJ} \\
\text{Football} &: 1\,\text{h} \times 2000\,\text{kJ/h} = 2000\,\text{kJ}
\end{align*}
Dépense totale (modèle additif) :
\[ 7500 + 600 + 2000 = \SI{10100}{kJ} \]
(2) Bilan énergétique :
\[ \text{Bilan} = \text{apports} - \text{dépenses} = 11500 - 10100 = +\SI{1400}{kJ} \]
Le bilan est \textbf{légèrement positif} ce jour-là : l'excédent sera stocké sous forme de lipides dans le tissu adipeux.
(3) Sans activité physique, la dépense totale tomberait à environ \SI{7500}{kJ} ; le bilan deviendrait $11500 - 7500 = +\SI{4000}{kJ}$ par jour. Répété durablement, cet excédent conduirait au \textbf{surpoids}, puis à l'\textbf{obésité}. L'activité physique apparaît donc comme un régulateur essentiel du bilan.
\end{exoresolu}

\courssub{Les effets bénéfiques de l'exercice régulier}

L'activité physique ne se contente pas de « brûler » des kilojoules pendant l'effort. Pratiquée régulièrement, elle agit en profondeur sur l'organisme :

\begin{itemize}
\item \textbf{Maintien et développement de la masse musculaire.} Le muscle est un tissu qui consomme de l'énergie même au repos. En préservant la masse musculaire, l'exercice évite la fonte musculaire observée lors des régimes trop restrictifs.
\item \textbf{Amélioration du métabolisme de base.} Une masse musculaire plus importante élève légèrement la dépense de repos : l'organisme « brûle » davantage, même sans bouger. L'exercice agit donc deux fois : pendant l'effort \emph{et} au repos.
\item \textbf{Protection cardiovasculaire.} L'effort régulier renforce le muscle cardiaque, améliore la circulation sanguine, contribue à faire baisser la pression artérielle et le taux de lipides sanguins (cholestérol), réduisant le risque d'infarctus et d'accident vasculaire cérébral à l'âge adulte.
\item \textbf{Prévention du surpoids et aide au traitement de l'obésité.} Associée à une alimentation équilibrée, l'activité physique augmente les dépenses et aide à ramener le bilan énergétique à l'équilibre, voire à le rendre modérément négatif pour mobiliser les réserves adipeuses.
\end{itemize}

\begin{propriete}[Régulation du bilan énergétique par l'activité physique — admise]
L'activité physique régulière \textbf{augmente la dépense énergétique totale} et \textbf{favorise le maintien d'un bilan énergétique équilibré} : elle aide à prévenir le surpoids et participe, avec une alimentation équilibrée, au traitement de l'obésité. Les mécanismes hormonaux précis de cette régulation (rôle de l'insuline, de la leptine, de l'adrénaline) seront étudiés en classe de Terminale.
\end{propriete}

\begin{attention}[Erreur fréquente : « le sport seul fait maigrir »]
Penser que l'exercice seul suffit à perdre du poids \emph{sans surveiller l'alimentation} est une erreur classique. Une heure de course ($\approx$ \SI{2500}{kJ}) peut être « annulée » en quelques minutes par un excès alimentaire (boisson sucrée, beignets, friture abondante). Les deux volets — \textbf{alimentation équilibrée} et \textbf{activité physique régulière} — sont \emph{complémentaires} et indissociables. Rédige toujours cette complémentarité dans tes conclusions au Probatoire.
\end{attention}

\courssub{Le triangle de l'équilibre : alimentation, activité physique, repos}

Un bilan énergétique sain ne repose pas sur un seul pilier, mais sur trois, que l'on représente classiquement par un triangle :

\begin{popfigure}
\begin{center}
\begin{tikzpicture}[scale=1.05]
\coordinate (A) at (90:3);
\coordinate (B) at (210:3);
\coordinate (C) at (330:3);
\draw[very thick, popblue, fill=popblueL!50] (A) -- (B) -- (C) -- cycle;
\node[align=center, font=\small\bfseries, text=popdark] at (90:3.9) {Alimentation\\équilibrée};
\node[align=center, font=\small\bfseries, text=popdark] at (210:3.9) {Activité\\physique};
\node[align=center, font=\small\bfseries, text=popdark] at (330:3.9) {Repos et\\sommeil};
\node[align=center, font=\small, text=popdark] at (0,0) {Bilan\\énergétique\\équilibré};
\end{tikzpicture}
\end{center}
\legende{Le triangle de l'équilibre : la santé et un bilan énergétique stable reposent sur trois piliers indissociables — une alimentation équilibrée, une activité physique régulière et un repos suffisant. La faiblesse d'un seul pilier fragilise l'ensemble.}
\end{popfigure}

Le troisième pilier, le \cle{repos}, est souvent négligé : le sommeil permet la récupération musculaire et la régulation des sécrétions hormonales ; un adolescent a besoin d'environ 8 à 9 heures de sommeil par nuit.

\begin{savaistu}[La sédentarité, un poison silencieux]
Des études épidémiologiques internationales montrent que la \cle{sédentarité} — rester assis plus de 8 heures par jour — est un \textbf{facteur de risque indépendant} de maladies cardiovasculaires, de diabète de type 2 et de mortalité prématurée, \emph{même chez les personnes qui ne sont pas obèses} et même chez celles qui font du sport par ailleurs. Autrement dit : courir une heure le soir ne « répare » pas totalement dix heures passées assis dans la journée. D'où les recommandations de se lever et de marcher régulièrement, y compris pendant les journées d'étude. Au Cameroun, l'urbanisation rapide (transports motorisés, écrans, travail de bureau) fait progresser la sédentarité, et avec elle les maladies dites « de civilisation ».
\end{savaistu}

\courssub{Attention au revers de la médaille : l'exercice intense exige des apports adaptés}

Tout ce qui précède ne doit pas faire oublier la leçon de la section précédente : si l'activité physique augmente les dépenses, un exercice intense et répété \textbf{sans adaptation des apports alimentaires} fait basculer le bilan en territoire négatif. C'est le cas des sportifs de haut niveau, mais aussi des jeunes qui enchaînent entraînements et compétitions tout en sautant des repas : fatigue chronique, fonte musculaire, baisse des défenses immunitaires, troubles de la croissance chez l'adolescent. Chez le sportif, l'alimentation doit être \emph{augmentée} (glucides pour l'effort, protéines pour la réparation musculaire, eau en abondance) pour maintenir le bilan à l'équilibre.

\begin{exoresolu}[Le cas d'un jeune footballeur en sur-entraînement]
Énoncé. Junior, 16 ans, s'entraîne au football 2 heures par jour ($\approx$ \SI{2000}{kJ/h}) en plus de son métabolisme de base et de son effet thermique des aliments (\SI{7500}{kJ}). Il ne consomme que \SI{9500}{kJ} par jour. (1) Calculer son bilan énergétique journalier. (2) Quelles conséquences prévoir si cette situation dure plusieurs semaines ? (3) Proposer une mesure de rééquilibrage.
\tcblower
Solution détaillée.
(1) Dépense totale (modèle additif) : $7500 + 2 \times 2000 = \SI{11500}{kJ}$. Bilan :
\[ 9500 - 11500 = -\SI{2000}{kJ} \quad \text{par jour : bilan négatif.} \]
(2) Un déficit de \SI{2000}{kJ/j} répété conduit à l'épuisement des réserves de glycogène, à la fonte des lipides puis des \textbf{protéines musculaires} : fatigue persistante, perte de poids, baisse des performances, sensibilité accrue aux infections — c'est la \cle{dénutrition} d'effort.
(3) Rééquilibrage : augmenter les apports (repas riches en glucides complexes — riz, igname, manioc — et en protéines — poisson, haricots, arachides) pour couvrir les \SI{11500}{kJ} dépensés, et préserver un jour de repos hebdomadaire. L'objectif n'est pas de supprimer le sport, mais d'\emph{adapter l'assiette à l'effort}.
\end{exoresolu}

\courssub{Recommandations pratiques pour l'adolescent}

\begin{aretenir}[L'essentiel de la section]
\begin{itemize}
\item L'\cle{activité physique} est le poste de dépense énergétique \textbf{le plus modulable} : elle augmente les dépenses sans toucher aux fonctions vitales.
\item Ordres de grandeur : marche rapide $\approx$ \SI{1200}{kJ/h} ; travaux champêtres $\approx$ \SI{1500}{kJ/h} ; football $\approx$ \SI{2000}{kJ/h} ; course $\approx$ \SI{2500}{kJ/h}.
\item Dépense d'une activité = dépense horaire $\times$ durée.
\item L'exercice régulier maintient la masse musculaire, améliore le métabolisme de base et protège le système cardiovasculaire ; il prévient le surpoids et aide à traiter l'obésité, \textbf{en association avec une alimentation équilibrée}.
\item Un exercice intense exige des apports alimentaires adaptés, sous peine de bilan négatif (cas des sportifs).
\item Recommandation : \textbf{au moins 30 à 60 minutes d'activité physique par jour} pour un adolescent (marche, sport, travaux, jeux actifs), un sommeil suffisant, et des interruptions régulières de la position assise.
\end{itemize}
\end{aretenir}

\begin{exercice}[À toi de jouer]
Marie passe 8 h assise en cours, marche rapidement 45 minutes par jour ($\approx$ \SI{1200}{kJ/h}) et ses dépenses de base (métabolisme de base + effet thermique des aliments) s'élèvent à \SI{7200}{kJ}. Ses apports sont de \SI{9000}{kJ}. Calcule sa dépense totale, puis son bilan énergétique journalier. Ce bilan est-il favorable à long terme ? Propose deux mesures concrètes pour l'améliorer.
\end{exercice}

\begin{corrige}[Exercice « À toi de jouer »]
Dépense de marche : $0{,}75\,\text{h} \times 1200 = \SI{900}{kJ}$. Dépense totale (modèle additif) : $7200 + 900 = \SI{8100}{kJ}$. Bilan : $9000 - 8100 = +\SI{900}{kJ}$ par jour, donc \textbf{légèrement positif} : à long terme, risque de prise de poids progressive. Mesures possibles : augmenter la durée de marche ou ajouter une activité sportive hebdomadaire (par exemple 2 séances de football) pour accroître les dépenses ; réduire modérément les apports les plus denses (fritures, boissons sucrées) tout en gardant une alimentation équilibrée ; limiter le temps passé assis en dehors des cours. L'idéal est d'agir sur les deux volets, dépenses \emph{et} apports.
\end{corrige}

\coursec{Rééquilibrage du bilan énergétique et mesures préventives}

Dans la section précédente, tu as découvert que l'activité physique est un formidable régulateur du bilan énergétique : en augmentant les dépenses, elle permet de compenser des apports un peu trop généreux et de maintenir l'équilibre. Mais que faire lorsque le déséquilibre est déjà installé ? Comment aider concrètement une personne en surpoids, ou au contraire une personne dénutrie, à retrouver un bilan équilibré ? C'est toute la question du \cle{rééquilibrage}, que nous allons aborder maintenant. Tu verras qu'il ne s'agit jamais de solutions « miracles », mais d'une démarche logique, progressive et mesurable — exactement le type de raisonnement exigé au Probatoire.

\courssub{Le principe général : agir sur les deux volets du bilan}

\textbf{Une situation concrète pour commencer.} Au centre de santé de Ngaoundéré, une infirmière reçoit deux patients le même matin. Le premier, M.~Abba, 45 ans, présente un IMC de 31 : il est obèse. La seconde, Aminatou, 17 ans, pèse 41 kg pour 1,62 m (IMC $\approx$ 15,6) : elle est dénutrie. Deux situations opposées... mais une même logique de résolution. Dans les deux cas, l'infirmière commence par établir le \emph{bilan énergétique} du patient, puis agit sur ses deux composantes.

Rappelle-toi la formule centrale de la leçon :
\[
\text{Bilan énergétique} = \text{Apports (alimentation)} - \text{Dépenses (métabolisme de base + activité physique)}
\]

Ce bilan est une \emph{différence} : il possède donc deux leviers d'action. Pour le modifier, on peut :
\begin{itemize}
\item agir sur les \cle{apports}, c'est-à-dire sur l'alimentation (quantité et qualité des aliments) ;
\item agir sur les \cle{dépenses}, c'est-à-dire sur l'activité physique (et, dans une faible mesure, sur le métabolisme de base, qui augmente avec la masse musculaire).
\end{itemize}

\begin{aretenir}[Principe du rééquilibrage]
Tout rééquilibrage du bilan énergétique repose sur une action combinée sur les \textbf{deux volets} du bilan : les \cle{apports} (alimentation) et les \cle{dépenses} (activité physique). Agir sur un seul volet est possible, mais l'action combinée est plus efficace, plus douce pour l'organisme et plus durable.
\end{aretenir}

Pourquoi combiner les deux leviers plutôt que d'en utiliser un seul ? Prenons une analogie : pour vider une bassine qui déborde, tu peux fermer le robinet (réduire les apports) \emph{ou} ouvrir davantage l'évacuation (augmenter les dépenses). Mais si tu fais les deux à moitié, la bassine se vide deux fois plus vite, sans avoir à fermer complètement le robinet — ce qui serait insupportable (qui a envie de ne plus jamais manger de beignets ?). C'est exactement la stratégie du rééquilibrage intelligent.

\courssub{Rééquilibrer un bilan énergétique positif (surpoids, obésité)}

Un bilan \textbf{positif} signifie que les apports dépassent les dépenses : l'excédent est stocké, principalement sous forme de triglycérides dans le tissu adipeux. Pour inverser la tendance, il faut rendre le bilan légèrement \emph{négatif}, de manière contrôlée. Trois mesures complémentaires :

\begin{enumerate}
\item \textbf{Réduire les aliments très gras et très sucrés.} Ce sont les aliments les plus denses en énergie : 1 g de lipides apporte environ 38 kJ (9 kcal), contre 17 kJ pour 1 g de glucides ou de protéines. Les fritures (beignets, puff-puff, frites de plantain), les huiles en excès, les boissons sucrées et les friandises sont les premières cibles. Attention : il s'agit de \emph{réduire}, pas de supprimer totalement les lipides, indispensables (vitamines liposolubles, membranes cellulaires, précurseurs hormonaux).
\item \textbf{Augmenter l'activité physique.} Marche rapide, travaux champêtres, sport, ménage actif : toute dépense supplémentaire creuse le bilan du bon côté. On vise au moins 30 minutes d'activité modérée par jour.
\item \textbf{Fractionner raisonnablement les repas.} Répartir les apports en 3 repas équilibrés (éventuellement avec une collation saine comme un fruit) évite les grandes fringales qui poussent aux excès, et évite le grignotage permanent. Sauter le petit-déjeuner pour « compenser » est une fausse bonne idée : il conduit presque toujours à trop manger le soir.
\end{enumerate}

\begin{exemplebox}[Résorber un excès de +1 000 kJ/jour]
Une élève présente un excès énergétique estimé à $+1\,000$ kJ par jour. Plutôt qu'un régime sévère, on répartit l'effort sur les deux volets :
\begin{itemize}
\item \textbf{Côté apports :} supprimer une boisson sucrée de 33 cl ($\approx 600$ kJ) et remplacer les beignets du goûter par un fruit $\Rightarrow$ environ $-500$ kJ/jour ;
\item \textbf{Côté dépenses :} ajouter 25 minutes de marche rapide quotidienne $\Rightarrow$ environ $+500$ kJ/jour de dépense.
\end{itemize}
Total : le bilan passe de $+1\,000$ kJ/jour à environ $0$, puis devient légèrement négatif si l'effort est maintenu. Aucune privation extrême, aucune épreuve sportive insurmontable : juste deux petits ajustements quotidiens.
\end{exemplebox}

\begin{savaistu}[Les calories « cachées » des boissons sucrées]
Une bouteille de 50 cl de soda apporte environ \textbf{900 kJ} — l'équivalent énergétique d'un bol de riz — \emph{sans aucun effet de satiété} : le cerveau ne « comptabilise » pas les calories liquides comme il le fait pour les aliments solides. Résultat : on boit ces 900 kJ \emph{en plus} du repas, sans se sentir plus rassasié. C'est pourquoi les boissons sucrées sont l'une des premières sources d'apports « cachés » à supprimer lors d'un rééquilibrage. L'eau, elle, n'apporte aucune énergie et reste la seule boisson indispensable.
\end{savaistu}

\courssub{Rééquilibrer un bilan énergétique négatif (dénutrition, épuisement)}

Un bilan \textbf{négatif} signifie que les dépenses dépassent les apports : l'organisme puise dans ses réserves (glycogène hépatique et musculaire, puis graisses, puis \emph{protéines musculaires}), d'où l'amaigrissement, la fatigue chronique et la baisse des défenses immunitaires. Le rééquilibrage vise à rendre le bilan légèrement \emph{positif} :

\begin{enumerate}
\item \textbf{Augmenter les apports énergétiques \emph{et} protéiques.} Il ne suffit pas de « manger plus » : il faut manger \emph{mieux}. Les protéines sont indispensables pour reconstruire les muscles fondus pendant la dénutrition. On enrichit les plats : bouillie de maïs enrichie à l'arachide ou au soja, haricots niébé, poisson, œufs, lait caillé, huile de palme rouge en quantité raisonnable (énergie + vitamine A).
\item \textbf{Réduire temporairement les dépenses excessives.} Un adolescent dénutri qui fait 8 km à pied chaque jour et travaille aux champs doit temporairement alléger ces efforts, le temps de reconstituer ses réserves. L'activité physique n'est pas supprimée (elle entretient l'appétit et la masse musculaire), mais ramenée à un niveau modéré.
\item \textbf{Assurer un suivi médical si nécessaire.} Une dénutrition sévère (IMC très bas, œdèmes, infections répétées) est une \emph{maladie} : elle relève du centre de santé, qui peut prescrire des aliments thérapeutiques prêts à l'emploi (pâtes d'arachide enrichies, utilisées dans les programmes de prise en charge de la malnutrition aiguë au Cameroun) et rechercher une cause sous-jacente (parasitoses, tuberculose, VIH...).
\end{enumerate}

\begin{attention}[La symétrie n'est pas parfaite]
Chez la personne dénutrie, on ne « fait pas du sport pour compenser » : augmenter les dépenses aggraverait le bilan négatif ! Inversement, chez la personne obèse, on ne supprime pas totalement l'activité « pour se reposer ». Chaque mesure doit être adaptée au \emph{sens} du déséquilibre. C'est l'erreur de raisonnement la plus fréquente dans les copies.
\end{attention}

\courssub{Le socle commun : une alimentation variée et équilibrée}

Que le bilan soit positif ou négatif, le rééquilibrage ne consiste jamais à manger « n'importe quoi, en moins » ou « n'importe quoi, en plus ». La qualité des apports est aussi importante que leur quantité.

\begin{definition}[Alimentation équilibrée]
Une \cle{alimentation équilibrée} est une alimentation qui couvre les besoins énergétiques et nutritionnels de l'organisme — énergie, protéines, lipides, glucides, vitamines, sels minéraux, eau — \textbf{sans excès ni carence}, grâce à la variété des aliments consommés.
\end{definition}

Concrètement, au Cameroun, une alimentation équilibrée s'appuie sur les plats locaux qui combinent naturellement plusieurs groupes d'aliments :
\begin{itemize}
\item des \textbf{céréales ou tubercules} pour l'énergie (riz, maïs, mil, manioc, igname, plantain, taro) ;
\item des \textbf{légumineuses} pour les protéines végétales (niébé, haricots, arachides, soja) — éventuellement complétées par poisson, viande ou œufs ;
\item des \textbf{légumes et fruits} pour les vitamines, minéraux et fibres (feuilles de foléré, ndolé, gombo, éru, mangues, papayes, bananes, avocats).
\end{itemize}

Un plat comme le \emph{ndolé} (feuilles + arachides + poisson ou viande, accompagné de miondo ou de plantain) est un excellent exemple d'association équilibrée : glucides, protéines, lipides, vitamines et fibres dans un même repas. Le rééquilibrage ne consiste donc pas à abandonner la cuisine locale, mais à ajuster les \emph{proportions} et les \emph{modes de cuisson} (moins de friture, moins d'huile et de sucre ajoutés).

\courssub{Un rééquilibrage nécessairement progressif}

\textbf{Observation.} Deux personnes veulent perdre 10 kg. La première suit un « régime choc » : elle perd 6 kg en trois semaines... puis les reprend en deux mois, avec 2 kg de bonus (c'est le fameux « effet yo-yo »). La seconde perd 0,75 kg par semaine pendant trois mois et stabilise durablement. Pourquoi cette différence ?

\textbf{Interprétation.} Un amaigrissement brutal prive l'organisme de protéines : il puise dans ses muscles, son métabolisme de base \emph{chute} (moins de muscle = moins de dépense au repos), et les fringales deviennent irrésistibles. Dès l'arrêt du régime, les apports repartent à la hausse alors que les dépenses sont diminuées : le bilan redevient fortement positif. À l'inverse, une variation lente préserve les muscles, laisse le temps aux habitudes de s'installer et ne déclenche pas les mécanismes de « défense » de l'organisme.

\begin{propriete}[Règle de progressivité]
Tout rééquilibrage du bilan énergétique doit être \textbf{progressif} : on vise une variation de masse corporelle de \textbf{0,5 à 1 kg par semaine au maximum}, que ce soit à la baisse (surpoids) ou à la hausse (dénutrition). Au-delà, la perte (ou le gain) concerne surtout l'eau et les muscles, et le résultat n'est pas durable.
\end{propriete}

\begin{attention}[Erreur fréquente : les solutions extrêmes]
Dans les exercices comme dans la vie, il est \textbf{interdit} de proposer un régime drastique (suppression totale des lipides ou des glucides, jeûne prolongé, « un seul repas par jour ») ou l'arrêt total de l'activité physique. Ces solutions extrêmes sont :
\begin{itemize}
\item \textbf{inefficaces} à long terme (effet yo-yo, fonte musculaire, baisse du métabolisme de base) ;
\item \textbf{dangereuses} : carences en vitamines et minéraux, troubles du rythme cardiaque, dénutrition iatrogène, troubles du comportement alimentaire chez l'adolescent.
\end{itemize}
La bonne réponse, au Probatoire comme à l'hôpital, est toujours une mesure \emph{modérée, combinée et progressive}.
\end{attention}

\courssub{Éducation nutritionnelle et dépistage : prévenir plutôt que guérir}

Le meilleur rééquilibrage est celui qu'on n'a pas besoin de faire. Deux outils de prévention :

\begin{itemize}
\item \textbf{L'éducation nutritionnelle} : apprendre, dès l'école, à reconnaître les groupes d'aliments, à lire sa faim et sa satiété, à limiter les boissons sucrées et les fritures, à pratiquer une activité physique régulière. Les campagnes de sensibilisation (journées de la santé, émissions radio en langues locales) jouent un rôle majeur, car les habitudes alimentaires se construisent dans la famille et la communauté.
\item \textbf{Le dépistage} : la \textbf{mesure régulière de la masse et le calcul de l'IMC} permettent de détecter précocement une dérive du bilan, \emph{avant} que le surpoids ne devienne obésité ou que la maigreur ne devienne dénutrition. On rappelle la formule :
\[ \text{IMC} = \frac{\text{masse (kg)}}{\text{taille (m)}^2} \]
avec les repères chez l'adulte : IMC $< 18,5$ = maigreur ; $18,5 \leq \text{IMC} < 25$ = corpulence normale ; $25 \leq \text{IMC} < 30$ = surpoids ; $\text{IMC} \geq 30$ = obésité. Chez l'enfant et l'adolescent, on utilise les courbes de croissance de l'OMS, car les seuils varient avec l'âge et le sexe.
\end{itemize}

Une pesée trimestrielle à la maison ou au centre de santé, notée sur un carnet, suffit à repérer une tendance : c'est le même principe que le carnet de santé de la mère et de l'enfant utilisé dans les consultations prénatales et de suivi infantile.

\courssub{Synthèse : la démarche complète d'analyse d'une situation}

Voici la démarche en quatre étapes que tu dois savoir dérouler face à n'importe quelle situation-problème de bilan énergétique.

\begin{methode}[Proposer un rééquilibrage adapté à une situation]
\begin{enumerate}
\item \textbf{Établir le bilan chiffré} : estimer ou calculer les apports (alimentation) et les dépenses (métabolisme de base + activité physique), puis calculer $\text{Bilan} = \text{Apports} - \text{Dépenses}$.
\item \textbf{Déterminer le signe du bilan} et poser le diagnostic : bilan $\approx 0$ = équilibre ; bilan $> 0$ durable = risque de surpoids/obésité ; bilan $< 0$ durable = risque de dénutrition/épuisement. Confirmer éventuellement par l'IMC (courbes OMS pour les adolescents).
\item \textbf{Agir sur les apports ET/OU les dépenses} selon le signe : réduire les apports gras/sucrés et augmenter l'activité si bilan positif ; augmenter les apports énergétiques et protéiques et réduire temporairement les dépenses excessives si bilan négatif.
\item \textbf{Fixer un objectif progressif et contrôlable} : variation de 0,5 à 1 kg/semaine maximum, avec suivi régulier de la masse et de l'IMC, et avis médical si le déséquilibre est sévère.
\end{enumerate}
\end{methode}

\begin{popfigure}
\begin{center}
\begin{tikzpicture}[
  node distance=0.55cm,
  etape/.style={rectangle, rounded corners=3pt, draw=popblue, fill=popblueL, text width=4.6cm, align=center, minimum height=1cm, font=\small},
  decision/.style={diamond, draw=poporange, fill=poporangeL, text width=3.4cm, align=center, font=\small, aspect=2.2},
  mesure/.style={rectangle, rounded corners=3pt, draw=popgreen, fill=popgreenL, text width=4.4cm, align=center, font=\small},
  fleche/.style={-{Stealth[length=2.5mm]}, thick, popdark}
]
\node[etape, align=center] (bilan){\textbf{1. Calcul du bilan}\\ Bilan = Apports $-$ Dépenses};
\node[decision, below=of bilan] (signe) {\textbf{2. Signe du bilan ?}};
\node[mesure, below left=1.1cm and -1.2cm of signe, align=center] (positif){\textbf{3a. Bilan $>0$ (surpoids)}\\ $\downarrow$ gras/sucrés, $\uparrow$ activité, repas fractionnés};
\node[mesure, below right=1.1cm and -1.2cm of signe, align=center] (negatif){\textbf{3b. Bilan $<0$ (dénutrition)}\\ $\uparrow$ apports énergétiques et protéiques, $\downarrow$ dépenses excessives, avis médical};
\node[etape, below=2.6cm of signe, fill=poppurpleL, draw=poppurple, align=center] (suivi){\textbf{4. Suivi}\\ objectif : 0,5 à 1 kg/semaine max ; contrôle masse + IMC};
\draw[fleche] (bilan) -- (signe);
\draw[fleche] (signe.west) -| node[pos=0.25, above, font=\footnotesize]{positif} (positif.north);
\draw[fleche] (signe.east) -| node[pos=0.25, above, font=\footnotesize]{négatif} (negatif.north);
\draw[fleche] (positif.south) |- (suivi.west);
\draw[fleche] (negatif.south) |- (suivi.east);
\end{tikzpicture}
\end{center}
\legende{Organigramme de la démarche de rééquilibrage : calcul du bilan énergétique, identification du signe du déséquilibre, choix des mesures adaptées (agissant sur les apports et/ou les dépenses), puis suivi progressif et contrôlé.}
\end{popfigure}

\begin{center}
\begin{tabularx}{\linewidth}{|l|X|X|}
\hline
\rowcolor{popblueL} \textbf{Critère} & \textbf{Bilan positif (surpoids, obésité)} & \textbf{Bilan négatif (dénutrition, épuisement)} \\
\hline
Objectif & Rendre le bilan légèrement négatif & Rendre le bilan légèrement positif \\
\hline
Action sur les apports & Réduire les aliments très gras et très sucrés ; supprimer les boissons sucrées ; fractionner les repas & Augmenter les apports énergétiques \emph{et} protéiques ; enrichir les plats (arachide, soja, poisson, œufs) \\
\hline
Action sur les dépenses & Augmenter l'activité physique ($\geq$ 30 min/jour) & Réduire temporairement les dépenses excessives, sans arrêter toute activité \\
\hline
Encadrement & Suivi de la masse et de l'IMC ; éducation nutritionnelle & Suivi médical si dénutrition sévère ; recherche d'une cause (infection, parasitose) \\
\hline
Rythme visé & Perte de 0,5 à 1 kg/semaine maximum & Gain de 0,5 à 1 kg/semaine maximum \\
\hline
À proscrire & Régime drastique, jeûne, suppression d'un groupe d'aliments & Arrêt total de l'activité physique, alimentation déséquilibrée « pour grossir vite » \\
\hline
\end{tabularx}
\end{center}

\begin{exoresolu}[Une démarche complète, de A à Z]
\textbf{Énoncé.} Rodrigue, 19 ans, mesure 1,70 m et pèse 88 kg. Son médecin estime ses apports à 14\,000 kJ/jour et ses dépenses totales à 12\,000 kJ/jour. 1) Calcule son bilan énergétique et son IMC, puis pose le diagnostic. 2) Propose un plan de rééquilibrage chiffré et progressif.
\tcblower
\textbf{Solution.}
\textbf{1) Bilan et diagnostic.}
\[ \text{Bilan} = 14\,000 - 12\,000 = +2\,000 \text{ kJ/jour} > 0 \]
\[ \text{IMC} = \frac{88}{1,70^2} = \frac{88}{2,89} \approx 30,4 \]
Le bilan est positif et l'IMC dépasse 30 : Rodrigue est en \textbf{obésité modérée}. Il faut inverser le signe de son bilan.

\textbf{2) Plan de rééquilibrage.} On répartit l'effort sur les deux volets, en visant un bilan d'environ $-1\,000$ à $-2\,000$ kJ/jour :
\begin{itemize}
\item \textbf{Apports} : supprimer le soda quotidien de 50 cl ($-900$ kJ) et réduire les fritures ($-600$ kJ) $\Rightarrow$ apports ramenés à environ 12\,500 kJ/jour ;
\item \textbf{Dépenses} : 45 minutes de marche rapide ou de football par jour ($+1\,500$ kJ) $\Rightarrow$ dépenses portées à environ 13\,500 kJ/jour.
\end{itemize}
Nouveau bilan : $12\,500 - 13\,500 = -1\,000$ kJ/jour, soit un déficit modéré compatible avec une perte de $0,5$ à $1$ kg/semaine. On ajoute un suivi mensuel de la masse et de l'IMC, et une alimentation variée (céréales/tubercules + légumineuses + légumes et fruits) pour éviter toute carence. Objectif raisonnable : atteindre un IMC inférieur à 25 (environ 72 kg) en 4 à 8 mois, sans régime drastique.
\end{exoresolu}

\begin{exercice}[À toi de jouer]
Sandrine, 16 ans, mesure 1,60 m et pèse 40 kg. Ses apports sont estimés à 7\,500 kJ/jour et ses dépenses à 9\,500 kJ/jour (elle parcourt 10 km à pied chaque jour). 1) Calcule son bilan et son IMC ; pose le diagnostic. 2) Propose trois mesures de rééquilibrage adaptées, en précisant sur quel volet chacune agit. 3) Explique pourquoi il serait dangereux de lui conseiller de « faire du sport pour se muscler » comme mesure principale.
\end{exercice}

\begin{corrige}[Exercice]
\textbf{1)} Bilan $= 7\,500 - 9\,500 = -2\,000$ kJ/jour $< 0$. IMC $= 40 / 1,60^2 = 40/2,56 \approx 15,6$. Chez l'adulte, IMC $< 18,5$ indique une maigreur ; chez l'adolescente de 16 ans, on se réfère aux courbes de croissance OMS : une valeur de 15,6 est très en dessous du 3e percentile, confirmant une \textbf{dénutrition sévère (maigreur sévère)}.

\textbf{2)} Mesures : (a) augmenter les apports énergétiques et protéiques — bouillie enrichie à l'arachide, niébé, œufs, poisson (volet \emph{apports}) ; (b) réduire temporairement les dépenses excessives, par exemple le trajet à pied partiellement en transport (volet \emph{dépenses}) ; (c) consulter un centre de santé pour un suivi médical et la recherche d'une cause (parasitose, infection), avec pesées régulières et objectif de 0,5 à 1 kg/semaine (volet \emph{suivi}).

\textbf{3)} Le sport augmente les \emph{dépenses} énergétiques : chez une personne dont le bilan est déjà négatif, il aggraverait le déficit, accélérerait la fonte musculaire et l'épuisement. La priorité est de restaurer les apports ; l'activité physique modérée ne sera réintroduite progressivement qu'une fois les réserves reconstituées.
\end{corrige}

\begin{aretenir}[L'essentiel de la section]
\begin{itemize}
\item Le rééquilibrage agit sur les \textbf{deux volets} du bilan : apports (alimentation) et dépenses (activité physique).
\item Bilan \textbf{positif} : réduire les aliments très gras/sucrés, augmenter l'activité, fractionner les repas. Bilan \textbf{négatif} : augmenter les apports énergétiques et protéiques, réduire temporairement les dépenses excessives, suivi médical.
\item Le socle commun est une \textbf{alimentation variée et équilibrée}, appuyée sur les plats locaux (céréales/tubercules + légumineuses + légumes et fruits).
\item Progressivité obligatoire : \textbf{0,5 à 1 kg/semaine maximum} ; les régimes drastiques et l'arrêt total de l'activité sont inefficaces et dangereux.
\item Prévention : éducation nutritionnelle et dépistage par la mesure régulière de la masse et de l'IMC (courbes OMS pour les jeunes).
\end{itemize}
\end{aretenir}

\newpage

\coursec{Activité d'intégration}

\courssub{Objectifs de l'activité}

À l'issue de cette activité d'intégration, tu seras capable de :

\begin{itemize}
\item calculer les \cle{apports énergétiques} totaux d'une journée à partir d'un menu détaillé ;
\item calculer les \cle{dépenses énergétiques} totales en additionnant le métabolisme de base et les dépenses liées aux activités ;
\item établir un \cle{bilan énergétique} et déterminer son signe ;
\item calculer et interpréter l'\cle{indice de masse corporelle} (IMC) ;
\item proposer des mesures de \cle{rééquilibrage} chiffrées et vérifier leur efficacité par le calcul.
\end{itemize}

\courssub{Situation}

Ngono est une élève de Première C au lycée de Bafoussam. Elle a 16 ans, mesure \SI{1,62}{m} et pèse \SI{68}{kg}. Lors de la visite médicale scolaire, l'infirmier lui signale qu'elle est en début de surpoids et l'invite à examiner son alimentation et son mode de vie. Ngono, étonnée, affirme : « Je ne mange pourtant pas tant que ça ! » Le professeur de SVT propose alors à la classe d'établir le bilan énergétique complet d'une journée type de Ngono, afin de vérifier si ses apports et ses dépenses sont équilibrés.

\textbf{Document 1 : menu d'une journée type de Ngono}

\begin{center}
\begin{tabularx}{\linewidth}{|l|X|c|}
\hline
\rowcolor{popblueL} \textbf{Repas} & \textbf{Composition} & \textbf{Valeur énergétique (kJ)} \\
\hline
Petit-déjeuner & Bouillie de maïs sucrée + 3 beignets haricots & 2 000 \\
\hline
Déjeuner & Riz (grande portion) + sauce arachide + viande de bœuf + 1 banane mûre & 5 500 \\
\hline
Goûter & 1 bouteille de soda (50 cL) + 2 beignets & 1 300 \\
\hline
Dîner & Foutou maïs + sauce jaune + poisson fumé & 4 000 \\
\hline
\end{tabularx}
\end{center}

\textbf{Document 2 : emploi du temps et dépenses énergétiques}

Le \cle{métabolisme de base} de Ngono (dépense minimale au repos pour assurer les fonctions vitales) est de \SI{6500}{kJ} par jour. Le tableau suivant donne ses activités quotidiennes et la dépense énergétique horaire correspondante.

\begin{center}
\begin{tabularx}{\linewidth}{|X|c|c|}
\hline
\rowcolor{popblueL} \textbf{Activité} & \textbf{Durée} & \textbf{Dépense horaire (kJ/h)} \\
\hline
Sommeil & 8 h & 250 \\
\hline
Cours (position assise) & 7 h & 300 \\
\hline
Marche (trajet maison--lycée) & 30 min & 1 000 \\
\hline
Tâches ménagères & 1 h & 600 \\
\hline
Téléphone / télévision (assis ou allongé) & 4 h & 225 \\
\hline
\end{tabularx}
\end{center}

\textbf{Document 3 : rappels sur l'IMC}

L'indice de masse corporelle se calcule par : $IMC = \dfrac{m}{t^{2}}$, avec $m$ la masse en kg et $t$ la taille en m. Interprétation chez l'adolescent de 16 ans (repères simplifiés) : IMC $< 18{,}5$ : maigreur ; $18{,}5 \leq IMC < 25$ : corpulence normale ; $25 \leq IMC < 30$ : surpoids ; $IMC \geq 30$ : obésité.

\courssub{Consignes}

Travail en groupes de 4 élèves. Rédige toutes les étapes du calcul sur ta feuille.

\begin{enumerate}
\item Calcule les \textbf{apports énergétiques totaux} de la journée de Ngono à partir du document 1.
\item Calcule la \textbf{dépense énergétique} de chaque activité (document 2), puis la dépense totale de la journée en ajoutant le métabolisme de base.
\item Établis le \textbf{bilan énergétique} de la journée (bilan = apports $-$ dépenses) et précise son signe. Que peut-on en prévoir à long terme pour la masse de Ngono ?
\item Calcule l'\textbf{IMC} de Ngono et interprète le résultat à l'aide du document 3. Le diagnostic de l'infirmier est-il confirmé ?
\item Propose \textbf{trois mesures concrètes de rééquilibrage} (au moins une modification de l'alimentation et un ajout d'activité physique), chiffre chaque mesure, puis vérifie par le calcul que le nouveau bilan est proche de zéro.
\item Restitution : chaque groupe présente ses propositions ; la classe confronte les résultats et débat des mesures les plus réalistes dans le contexte de Ngono.
\end{enumerate}

\courssub{Ressources à mobiliser}

\begin{aretenir}[Formules utiles]
\begin{itemize}
\item Apports totaux = somme des valeurs énergétiques de tous les repas et boissons.
\item Dépenses totales = métabolisme de base $+$ somme des dépenses des activités (dépense d'une activité = durée $\times$ dépense horaire).
\item Bilan énergétique = apports $-$ dépenses. Bilan $> 0$ : stockage (prise de masse) ; bilan $< 0$ : déstockage (perte de masse) ; bilan $\approx 0$ : équilibre.
\item $IMC = \dfrac{m}{t^{2}}$ (kg/m$^2$).
\item Repère utile : \SI{1}{kg} de tissu adipeux correspond à environ \SI{38000}{kJ} stockés.
\end{itemize}
\end{aretenir}

\courssub{Démarche et corrigé commenté}

\begin{exoresolu}[Bilan énergétique de la journée de Ngono]
Énoncé : à partir des documents 1, 2 et 3, établis le bilan énergétique de Ngono, interprète son IMC et propose un rééquilibrage chiffré.

\tcblower

\textbf{Étape 1 : calcul des apports totaux.}

On additionne les valeurs énergétiques des quatre prises alimentaires :
\begin{align*}
A &= 2\,000 + 5\,500 + 1\,300 + 4\,000 \\
A &= 12\,800 \ \text{kJ}
\end{align*}
Les apports de la journée sont de \SI{12800}{kJ}.

\textbf{Étape 2 : calcul des dépenses totales.}

Pour chaque activité : dépense = durée $\times$ dépense horaire.
\begin{align*}
\text{Sommeil} &: 8 \times 250 = 2\,000 \ \text{kJ} \\
\text{Cours} &: 7 \times 300 = 2\,100 \ \text{kJ} \\
\text{Marche} &: 0{,}5 \times 1\,000 = 500 \ \text{kJ} \\
\text{Tâches ménagères} &: 1 \times 600 = 600 \ \text{kJ} \\
\text{Téléphone / TV} &: 4 \times 225 = 900 \ \text{kJ}
\end{align*}
Somme des dépenses d'activités : $2\,000 + 2\,100 + 500 + 600 + 900 = 6\,100$ kJ.

Dépense totale de la journée :
\[ D = 6\,500 + 6\,100 = 12\,600 \ \text{kJ} \]

\textbf{Étape 3 : bilan énergétique.}

\[ B = A - D = 12\,800 - 12\,600 = +200 \ \text{kJ} \]

Le bilan est \textbf{positif} : Ngono consomme chaque jour environ \SI{200}{kJ} de plus qu'elle n'en dépense. Cet excédent, faible en apparence, est stocké sous forme de graisse. À long terme :
\[ 200 \times 365 = 73\,000 \ \text{kJ/an} \approx \frac{73\,000}{38\,000} \approx 1{,}9 \ \text{kg de graisse par an} \]
Cela confirme qu'un petit excédent quotidien suffit à provoquer une prise de masse progressive : le surpoids s'installe lentement, souvent sans qu'on s'en aperçoive.

\textbf{Étape 4 : calcul et interprétation de l'IMC.}

\[ IMC = \frac{m}{t^{2}} = \frac{68}{1{,}62^{2}} = \frac{68}{2{,}6244} \approx 25{,}9 \ \text{kg/m}^{2} \]

Comme $25 \leq 25{,}9 < 30$, Ngono est en situation de \textbf{surpoids}. Le diagnostic de l'infirmier est confirmé. Le bilan positif calculé à l'étape 3 en donne l'explication énergétique.

\textbf{Étape 5 : propositions de rééquilibrage chiffrées.}

\begin{itemize}
\item \textbf{Mesure 1 (alimentation)} : remplacer le soda quotidien par de l'eau. Économie : environ \SI{600}{kJ/jour}. Nouveaux apports : $12\,800 - 600 = 12\,200$ kJ.
\item \textbf{Mesure 2 (activité physique)} : ajouter 30 minutes de marche rapide chaque jour (par exemple en rentrant du lycée à pied). Dépense supplémentaire : $0{,}5 \times 1\,000 = 500$ kJ. Nouvelles dépenses : $12\,600 + 500 = 13\,100$ kJ.
\item \textbf{Mesure 3 (hygiène de vie)} : réduire le temps passé devant le téléphone/la télévision de 4 h à 2 h et consacrer le temps libéré à des tâches actives ou à du sport.
\end{itemize}

Vérification avec les mesures 1 et 2 :
\[ B' = 12\,200 - 13\,100 = -900 \ \text{kJ} \]

Le nouveau bilan est \textbf{proche de zéro et légèrement négatif}, ce qui est exactement ce qu'il faut : un déficit modéré permet de résorber progressivement le surpoids (environ \SI{1}{kg} de graisse en 6 semaines), sans régime drastique dangereux pour une adolescente en pleine croissance. Une fois la corpulence normale atteinte, Ngono pourra remonter légèrement ses apports pour stabiliser le bilan autour de zéro.
\end{exoresolu}

\begin{attention}[Pièges fréquents]
\begin{itemize}
\item Ne pas oublier le \cle{métabolisme de base} dans les dépenses : il représente ici plus de la moitié de la dépense totale.
\item Convertir les durées en heures : 30 min $= 0{,}5$ h avant de multiplier par la dépense horaire.
\item Ne pas conclure « maigreur » ou « obésité » trop vite : vérifie les bornes de l'IMC. Ici $IMC \approx 25{,}9$, c'est du \textbf{surpoids}, pas de l'obésité.
\item Un bilan légèrement positif chaque jour n'est pas anodin : c'est l'\textbf{accumulation} sur des mois qui produit la prise de masse.
\end{itemize}
\end{attention}

\courssub{Bilan de l'activité}

Cette activité t'a permis de mettre en évidence plusieurs idées essentielles :

\begin{itemize}
\item Le \cle{bilan énergétique} se calcule concrètement : apports (alimentation) d'un côté, dépenses (métabolisme de base + activités) de l'autre.
\item Un déséquilibre même faible ($+200$ kJ/jour, soit l'équivalent d'un petit beignet) provoque, par accumulation, une prise de masse d'environ \SI{2}{kg} par an : le surpoids et l'obésité résultent souvent de petits excès quotidiens prolongés.
\item L'\cle{IMC} est un outil simple de diagnostic du surpoids, qui confirme ici l'alerte médicale.
\item Le rééquilibrage repose sur \textbf{deux leviers complémentaires} : réduire les apports superflus (boissons sucrées, fritures) et augmenter les dépenses par l'\cle{activité physique}. La vérification par le calcul montre qu'un bilan proche de zéro, légèrement négatif, permet un retour progressif et sain à une corpulence normale.
\item Enfin, tu as mobilisé les trois savoir-faire de la leçon : établir un bilan énergétique, interpréter ses conséquences physiologiques et proposer des mesures de rééquilibrage adaptées à une situation réelle.
\end{itemize}

\newpage

\coursec{Méthodes et exercices résolus}

\coursec{Déséquilibres énergétiques des organismes et leurs conséquences}

\courssub{Situation déclenchante : deux visages du déséquilibre au Cameroun}
\begin{itemize}
\item \textbf{Cas 1 (Yaoundé, Douala) :} Une enqu\^ete de 2022 montre que pr\`es de 30\% des adultes en milieu urbain sont en surpoids ou ob\`esit\'e (IMC $\ge$ 25). L'alimentation riche en huiles, fritures (beignets, poisson brais\'e), boissons sucr\'ees et la s\'edentarit\'e (voiture, \'ecrans) cr\'eent un exc\'edent \'energ\'etique chronique.
\item \textbf{Cas 2 (Extr\^eme-Nord, r\'egions en crise) :} Pr\`es de 30\% des enfants de moins de 5 ans souffrent de malnutrition aigu\"e ou chronique (rapport SMART 2023). L'ins\'ecurit\'e alimentaire, les r\'ecoltes insuffisantes et les d\'eplacements de populations entra\^inent un d\'eficit \'energ\'etique s\'ev\`ere.
\end{itemize}
\textbf{Probl\'ematique :} Comment l'organisme g\`ere-t-il ses apports et ses d\'epenses \'energ\'etiques ? Quelles sont les cons\'equences physiologiques d'un d\'esequilibre (positif ou n\'egatif) et comment y rem\'edier ?

\courssub{Les besoins \'energ\'etiques de l'organisme}

\begin{definition}[Besoins \'energ\'etiques journaliers (BEJ)]
Quantit\'e d'\'energie (en kJ ou kcal) que l'organisme doit recevoir par l'alimentation pour couvrir l'ensemble de ses d\'epenses \'energ\'etiques et maintenir une masse corporelle stable, une composition corporelle saine et un niveau d'activit\'e compatible avec la sant\'e. Ils varient selon l'\^age, le sexe, la masse corporelle, l'\'etat physiologique (grossesse, allaitement, croissance) et surtout le \textbf{niveau d'activit\'e physique (NAP)}.
\end{definition}

\begin{definition}[M\'etabolisme de base (MB) ou D\'epense \'energ\'etique de base (DEB)]
\'Energie minimale indispensable \`a la vie v\'eg\'etative au repos complet, \`a jeun, en neutralit\'e thermique (18--20 \textdegree C), pour maintenir les fonctions vitales (ventilation, circulation sanguine, activit\'e c\'er\'ebrale, turn-over prot\'eique, pompes ioniques, thermor\'egulation). Il repr\'esente \textbf{60 à 70 \%} de la d\'epense totale journali\`ere.
\end{definition}

\begin{propriete}[Facteurs faisant varier le m\'etabolisme de base]
\begin{center}
\begin{tabularx}{\linewidth}{|l|X|}\hline
\rowcolor{popblueL} \textbf{Facteur} & \textbf{Effet sur le MB} \\ \hline
\cle{Masse musculaire} & Principal d\'eterminant : le muscle est tr\`es m\'etaboliquement actif. Plus la masse musculaire est \'elev\'ee, plus le MB est haut. \\ \hline
\cle{Sexe} & Les hommes ont g\'en\'eralement un MB plus \'elev\'e que les femmes \`a masse \'egale (plus de muscle, moins de graisse). \\ \hline
\cle{\^Age} & \'Elev\'e chez l'enfant (croissance) et l'adolescent ; d\'ecro\^it avec l'\^age (perte musculaire, s\'edentarit\'e). \\ \hline
\cle{\'Etat physiologique} & Grossesse (+10 à 15 \%), allaitement (+25 \%), f\'evre (+13 \% par \textdegree C au-dessus de 37 \textdegree C), hyperthyro\"idie. \\ \hline
\cle{Climat / Temp\'erature} & Augmentation par le froid (thermog\'en\`ese de frisson) ou chaleur extr\^eme (ventilation). \\ \hline
\end{tabularx}
\end{center}
\end{propriete}

\begin{aretenir}[Composantes de la d\'epense \'energ\'etique totale (DET)]
La DET = MB + \textbf{Thermor\'egulation} (maintien 37 \textdegree C) + \textbf{Action dynamique sp\'ecifique (ADS)} (co\^ut \'energ\'etique de la digestion, absorption, stockage : $\approx$ 10 \% des apports) + \textbf{Activit\'e physique} (variable de 15 \% \`a 50 \% de la DET selon le NAP).
\end{aretenir}

\begin{experience}[Estimation du m\'etabolisme de base (Formules de Black et al. / OMS 1985, utilis\'ees au Probatoire)]
\begin{itemize}
\item \textbf{Mat\'eriel :} Balance, toise, tableau de formules.
\item \textbf{Protocole :} Mesurer la masse $m$ (kg) et la taille $T$ (m). Appliquer la formule selon le sexe et la tranche d'\^age.
\item \textbf{Formules simplifi\'ees (kJ/jour) :}
\begin{align*}
\text{Homme 18--30 ans} &: MB \approx 63 \times m + 2896 \\
\text{Femme 18--30 ans} &: MB \approx 62 \times m + 2036
\end{align*}
\item \textbf{Interpr\'etation :} Ces formules donnent une estimation du MB. La DET = MB $\times$ \textbf{Coefficient d'activit\'e} (1,3 s\'edentaire ; 1,5 actif mod\'er\'e ; 1,7 tr\`es actif ; 2,0 athl\`ete).
\end{itemize}
\legende{Estimation du m\'etabolisme de base et de la d\'epense totale selon l'activit\'e physique.}
\end{experience}

\courssub{Le bilan \'energ\'etique : apports et d\'epenses}

\begin{definition}[Apports \'energ\'etiques (AE)]
\'Energie fournie par les macronutriments de l'alimentation, lib\'er\'ee par la respiration cellulaire (oxydation compl\`ete). On utilise les \textbf{\'equivalences \'energ\'etiques d'Atwater} (valeurs nettes m\'etabolisables) :
\begin{center}
\begin{tabular}{|c|c|c|}\hline
\rowcolor{popblueL} \textbf{Nutriment} & \textbf{\'Energie (kJ/g)} & \textbf{\'Energie (kcal/g)} \\ \hline
Glucides & 17 & 4 \\ \hline
Lipides & 38 & 9 \\ \hline
Prot\'eines & 17 & 4 \\ \hline
Alcool (non nutritif) & 29 & 7 \\ \hline
\end{tabular}
\end{center}
\end{definition}

\begin{definition}[Bilan \'energ\'etique (B)]
Diff\'erence entre les apports \'energ\'etiques (AE) et la d\'epense \'energ\'etique totale (DET) sur une p\'eriode donn\'ee (g\'en\'eralement la journ\'ee) :
\[
B = AE - DET \quad (\text{en kJ ou kcal})
\]
\end{definition}

\begin{propriete}[Interpr\'etation du signe du bilan]
\begin{center}
\begin{tabularx}{\linewidth}{|c|X|X|}\hline
\rowcolor{popblueL} \textbf{Bilan} & \textbf{Signification physiologique} & \textbf{Cons\'equence \`a long terme} \\ \hline
$B \approx 0$ & \cle{\'Equilibre \'energ\'etique} : Apports = D\'epenses. & Masse corporelle stable. Homeostasie \'energ\'etique. \\ \hline
$B > 0$ & \cle{Bilan positif} (Exc\'edent) : Apports $>$ D\'epenses. & Stockage de l'exc\'edent (lipogen\`ese) $\rightarrow$ Prise de masse grasse $\rightarrow$ Surpoids / Ob\'esit\'e. \\ \hline
$B < 0$ & \cle{Bilan n\'egatif} (D\'eficit) : Apports $<$ D\'epenses. & Mobilisation des r\'eserves (glycog\'enolyse, lipolyse, n\'eoglucogen\`ese) $\rightarrow$ Perte de masse $\rightarrow$ D\'enutrition / \'Epuisement. \\ \hline
\end{tabularx}
\end{center}
\end{propriete}

\begin{center}
\begin{popfigure}
\begin{tikzpicture}[node distance=1.2cm, >=Stealth, thick]
\node[draw, rounded corners, fill=popgreenL, minimum width=3cm, minimum height=1.2cm, align=center] (apports){APORTS \\ (Alimentation \\ Glucides, Lipides, Prot\'eines)};
\node[draw, rounded corners, fill=poporangeL, minimum width=4cm, minimum height=2.5cm, right=of apports, align=center] (organisme){ORGANISME \\ \textbf{M\'etabolisme de base} \\ Thermor\'egulation \\ Action Dynamique Sp\'ecifique};
\node[draw, rounded corners, fill=popblueL, minimum width=3cm, minimum height=1.2cm, below=of organisme, align=center] (activite){ACTIVIT\'E PHYSIQUE \\ (Marche, Sport, Travail)};
\node[draw, rounded corners, fill=poppinkL, minimum width=3cm, minimum height=1.2cm, left=of organisme, align=center] (stockage){STOCKAGE / \\ MOBILISATION \\ (Tissu adipeux, Glycog\`ene, Muscle)};
\draw[->, popgreen] (apports) -- node[above, sloped] {\textbf{Entr\'ee \'energ\'etique}} (organisme);
\draw[->, poporange] (organisme) -- node[right] {\textbf{D\'epense de base}} (activite);
\draw[->, popblue] (activite) -- node[right] {\textbf{D\'epense activit\'e}} (organisme);
\draw[<->, poppink, dashed] (organisme) -- node[above, sloped] {\textbf{Bilan $B$}} (stockage);
\end{tikzpicture}
\legende{Schéma du bilan énergétique : les apports alimentaires couvrent les dépenses (base + activité) ; l'excédent ou le déficit est géré par le stockage ou la mobilisation des réserves.}
\end{popfigure}
\end{center}

\courssub{Établir un bilan énergétique simple à partir de données}

\begin{methode}[Calculer un bilan \'energ\'etique]
\begin{enumerate}
\item \textbf{Identifier les apports} : additionner l'\'energie fournie par chaque aliment consommé dans la journ\'ee (en kJ ou en kcal). Si l'\'enonc\'e donne les masses de glucides, lipides et prot\'eines, utiliser les \'equivalences \'energ\'etiques : 1 g de glucides $\approx$ 17 kJ ; 1 g de lipides $\approx$ 38 kJ ; 1 g de prot\'eines $\approx$ 17 kJ.
\item \textbf{Identifier les d\'epenses} : additionner la d\'epense de base (m\'etabolisme de base + thermor\'egulation + digestion/ADS) et les d\'epenses li\'ees aux activit\'es physiques de la journ\'ee.
\item \textbf{Calculer le bilan} : $B = \text{Apports} - \text{D\'epenses}$ (en kJ).
\item \textbf{Interpr\'eter le signe} : $B > 0$ : bilan positif (stockage, prise de masse) ; $B < 0$ : bilan n\'egatif (d\'estockage, perte de masse) ; $B \approx 0$ : \'equilibre (masse stable).
\item \textbf{Conclure} par une phrase compl\`ete reliant le r\'esultat num\'erique \`a la situation de l'individu.
\end{enumerate}
\end{methode}

\begin{exoresolu}[Bilan \'energ\'etique d'un lyc\'een de Yaound\'e]
Énoncé. Marc, élève de Première C, consomme en une journ\'ee : 400 g de glucides, 80 g de lipides et 70 g de prot\'eines. Sa d\'epense \'energ\'etique totale de la journ\'ee est estim\'ee à \SI{11000}{kJ}. On donne : 1 g de glucides $=$ 17 kJ ; 1 g de lipides $=$ 38 kJ ; 1 g de prot\'eines $=$ 17 kJ.
\begin{enumerate}
\item Calculer les apports \'energ\'etiques de Marc.
\item Établir son bilan \'energ\'etique et l'interpr\'eter.
\end{enumerate}
\tcblower
Solution détaillée.
\begin{enumerate}
\item \textbf{Calcul des apports.} On calcule l'\'energie fournie par chaque cat\'egorie de nutriments :
\begin{align*}
E_{\text{glucides}} &= 400 \times 17 = \SI{6800}{kJ}\\
E_{\text{lipides}} &= 80 \times 38 = \SI{3040}{kJ}\\
E_{\text{prot\'eines}} &= 70 \times 17 = \SI{1190}{kJ}
\end{align*}
Apports totaux : $6800 + 3040 + 1190 = \SI{11030}{kJ}$.
\item \textbf{Bilan \'energ\'etique.} $B = \text{Apports} - \text{D\'epenses} = 11030 - 11000 = +\SI{30}{kJ}$.
Le bilan est l\'eg\`erement positif : les apports d\'epassent tr\`es faiblement les d\'epenses. L'exc\'edent de \SI{30}{kJ} sera stock\'e sous forme de graisse dans le tissu adipeux. Sur une journ\'ee, c'est n\'egligeable ; mais si ce l\'eger exc\'edent se r\'ep\`ete chaque jour, il entra\^inera \`a long terme une prise de masse progressive.
\end{enumerate}
\textbf{Conclusion :} les apports de Marc s'\'el\`evant à \SI{11030}{kJ} et son bilan \'energ\'etique est de $+\SI{30}{kJ}$ : il est pratiquement à l'équilibre, avec une très légère tendance au stockage.
\end{exoresolu}

\courssub{Prévoir la variation de masse corporelle à partir d'un déséquilibre}

\begin{methode}[Convertir un excédent ou un déficit énergétique en variation de masse]
\begin{enumerate}
\item Calculer le bilan journalier $B$ (voir méthode précédente).
\item Multiplier par le nombre de jours pour obtenir l'excédent (ou le déficit) cumulé : $E = B \times n$.
\item Utiliser la conversion fournie par l'énoncé (valeur de référence : 1 kg de tissu adipeux $\approx$ \SI{30000}{kJ}, soit environ \SI{7000}{kcal}).
\item Diviser : variation de masse $= \dfrac{E}{\SI{30000}{kJ/kg}}$.
\item Vérifier la cohérence du signe : excédent $\Rightarrow$ prise de masse ; déficit $\Rightarrow$ perte de masse.
\end{enumerate}
\end{methode}

\begin{exoresolu}[Conséquence d'un excédent alimentaire chronique]
Énoncé. Une employée de bureau de Douala consomme en moyenne \SI{10500}{kJ} par jour alors que sa dépense totale est de \SI{10000}{kJ}. On admet que 1 kg de tissu adipeux correspond à environ \SI{30000}{kJ}.
\begin{enumerate}
\item Calculer son excédent énergétique journalier.
\item Calculer la prise de masse prévisible au bout de 60 jours si les habitudes restent identiques.
\end{enumerate}
\tcblower
Solution détaillée.
\begin{enumerate}
\item \textbf{Excédent journalier.} $B = 10500 - 10000 = +\SI{500}{kJ}$ par jour. Le bilan est positif : chaque jour, \SI{500}{kJ} ne sont pas dépensés et sont stockés.
\item \textbf{Prise de masse sur 60 jours.}
\begin{align*}
E_{\text{cumulé}} &= 500 \times 60 = \SI{30000}{kJ}\\
\text{Prise de masse} &= \frac{30000}{30000} = \SI{1}{kg}
\end{align*}
\end{enumerate}
\textbf{Conclusion :} avec un excédent de seulement \SI{500}{kJ} par jour (l'équivalent d'un beignet-haricots copieux), cette employée prendrait environ \SI{1}{kg} en deux mois, soit près de \SI{6}{kg} par an. Cela illustre comment un petit déséquilibre chronique conduit progressivement au surpoids puis à l'obésité.
\end{exoresolu}

\courssub{Le déséquilibre énergétique positif : surpoids et obésité}

\begin{definition}[Surpoids et Obésité]
Accumulation anormale ou excessive de graisse corporelle (tissu adipeux) présentant un risque pour la santé. Le diagnostic clinique utilise l'\textbf{Indice de Masse Corporelle (IMC)}.
\end{definition}

\begin{propriete}[Mécanismes physiologiques du bilan positif]
Un apport excédentaire en glucides et lipides entraîne :
\begin{enumerate}
\item \textbf{Lipogenèse hépatique :} Excès de glucose $\rightarrow$ Acétyl-CoA $\rightarrow$ Acides gras $\rightarrow$ Triglycérides (VLDL) libérés dans le sang.
\item \textbf{Stockage adipocytaire :} Les triglycérides circulants sont hydrolysés par la \textbf{lipoprotéine lipase (LPL)} au niveau des capillaires du tissu adipeux ; les acides gras libres entrent dans l'adipocyte et sont ré-estérifiés (lipogenèse).
\item \textbf{Hypertrophie puis hyperplasie des adipocytes :} Les adipocytes grossissent (hypertrophie) ; au-delà d'une taille critique, ils se divisent (hyperplasie), augmentant le nombre de cellules graisseuses (irréversible).
\end{enumerate}
\end{propriete}

\begin{aretenir}[Conséquences pathologiques de l'obésité (adiposité viscérale surtout)]
L'obésité n'est pas qu'esthétique : c'est une \textbf{maladie chronique inflammatoire}. Le tissu adipeux viscéral sécrète des adipokines pro-inflammatoires (TNF-$\alpha$, IL-6, leptine) et moins d'adiponectine (protectrice).
\begin{itemize}
\item \textbf{Métaboliques :} Résistance à l'insuline $\rightarrow$ \cle{Diabète de type 2}, dyslipidémie (hypertriglycéridémie, HDL bas), stéatose hépatique non alcoolique (NASH).
\item \textbf{Cardiovasculaires :} \cle{Hypertension artérielle}, athérosclérose accélérée, infarctus du myocarde, AVC, insuffisance cardiaque.
\item \textbf{Mécaniques :} \cle{Arthrose} (genoux, hanches), apnée du sommeil, hernie hiatale.
\item \textbf{Autres :} Certains cancers (sein, côlon, endomètre), infertilité (SOPK), dépression, stigmatisation sociale.
\end{itemize}
\end{aretenir}

\courssub{Exploiter l'IMC pour qualifier un état pondéral}

\begin{methode}[Calculer et interpréter l'Indice de Masse Corporelle]
\begin{enumerate}
\item Relever la masse $m$ (en kg) et la taille $T$ (en m).
\item Appliquer la formule : $\mathrm{IMC} = \dfrac{m}{T^2}$ (en kg/m$^2$). Attention : la taille doit être en \textbf{mètres} et elle est \textbf{au carré}.
\item Comparer le résultat aux seuils de référence (adulte) :
\begin{center}
\begin{tabularx}{\linewidth}{|l|X|}\hline
\rowcolor{popblueL} \textbf{IMC (kg/m$^2$)} & \textbf{État pondéral} \\ \hline
$< 18,5$ & Maigreur (dénutrition si très bas) \\ \hline
18,5 à 24,9 & Corpulence normale \\ \hline
25 à 29,9 & Surpoids \\ \hline
30 à 34,9 & Obésité modérée (Classe I) \\ \hline
35 à 39,9 & Obésité sévère (Classe II) \\ \hline
$\geq 40$ & Obésité massive / morbide (Classe III) \\ \hline
\end{tabularx}
\end{center}
\item Conclure en nommant l'état pondéral et en rappelant le risque associé.
\end{enumerate}
\end{methode}

\begin{exoresolu}[Diagnostic pondéral par l'IMC]
Énoncé. Mme Ngo Bell mesure \SI{1,60}{m} et pèse \SI{82}{kg}.
\begin{enumerate}
\item Calculer son IMC.
\item Qualifier son état pondéral et citer deux risques pour sa santé.
\end{enumerate}
\tcblower
Solution détaillée.
\begin{enumerate}
\item \textbf{Calcul de l'IMC.}
\[
\mathrm{IMC} = \frac{m}{T^2} = \frac{82}{(1,60)^2} = \frac{82}{2,56} \approx \SI{32,0}{kg/m^2}
\]
\item \textbf{Interprétation.} Un IMC de 32,0 kg/m$^2$ se situe dans l'intervalle [30 ; 34,9] : Mme Ngo Bell est en \cle{obésité modérée (Classe I)}. Cet état résulte d'un bilan énergétique positif prolongé (apports supérieurs aux dépenses, stockage de graisse). Elle présente un risque élevé de :
\begin{itemize}
\item diabète de type 2 (résistance des cellules à l'insuline) ;
\item hypertension artérielle et maladies cardiovasculaires (infarctus, AVC).
\end{itemize}
\end{enumerate}
\textbf{Conclusion :} avec un IMC $\approx$ 32 kg/m$^2$, Mme Ngo Bell est en obésité modérée ; un rééquilibrage alimentaire associé à une activité physique régulière est nécessaire pour prévenir le diabète et les complications cardiovasculaires.
\end{exoresolu}

\courssub{Le déséquilibre énergétique négatif : dénutrition et épuisement}

\begin{definition}[Dénutrition (Malnutrition énergétique)]
État pathologique résultant d'un bilan énergétique \textbf{chroniquement négatif} (apports $<$ dépenses), entraînant une perte de masse corporelle, une altération de la composition corporelle (fonte musculaire, perte de graisse) et une détérioration des fonctions physiologiques (immunité, cicatrisation, cognition, croissance).
\end{definition}

\begin{propriete}[Mécanismes d'adaptation au jeûne / bilan négatif (Mobilisation des réserves)]
L'organisme hiérarchise l'utilisation des substrats pour préserver la glycémie (carburant obligatoire du cerveau et globules rouges) :
\begin{enumerate}
\item \textbf{Phase 1 (0--24 h) : Glycogène.} Glycogénolyse hépatique (maintien glycémie). Réserves limitées ($\approx$ 400--500 kJ).
\item \textbf{Phase 2 (1--3 jours) : Lipides + Protéines.} Lipolyse adipose intense $\rightarrow$ Acides gras (carburant muscle, foie, cœur) + Corps cétoniques (carburant cerveau). Protéolyse musculaire modérée pour néoglucogenèse (alanine, glutamine).
\item \textbf{Phase 3 (Jeûne prolongé / Dénutrition) : Épargne protéique.} Cerveau utilise massivement les corps cétoniques ($\beta$-hydroxybutyrate) $\rightarrow$ économie de glucose $\rightarrow$ diminution de la protéolyse musculaire. Mais si apports restent nuls/insuffisants : \textbf{fonte musculaire majeure} (perte force, respiration, immunité), atrophie viscérale (intestin, cœur), œdèmes (hypoalbuminémie).
\end{enumerate}
\end{propriete}

\begin{aretenir}[Formes cliniques de dénutrition sévère (enfant) - Contexte Cameroun]
\begin{itemize}
\item \textbf{Marasme (Dénutrition globale) :} Déficit énergétique + protéique pur. Enfant "vieux", peau fripée, masse $<$ 60\% norme, pas d'œdème. Fonte musculaire extrême. Risque vital immédiat (hypoglycémie, hypothermie, infection).
\item \textbf{Kwashiorkor (Dénutrition protéique prédominante) :} Apports énergétiques juste suffisants (tubercules) mais protéines quasi absentes. \textbf{Œdèmes} (hypoalbuminémie), dermatose (peau qui pèle), hépatomégalie (stéatose), cheveux décolorés ("cheveux en drapeau"), anorexie sévère.
\item \textbf{Marasme-Kwashiorkor :} Forme mixte (œdèmes + maigreur extrême).
\end{itemize}
\end{aretenir}

\courssub{Interpréter les conséquences physiologiques d'un déséquilibre énergétique}

\begin{methode}[Raisonner sur les conséquences d'un déséquilibre]
\begin{enumerate}
\item \textbf{Déterminer le sens du déséquilibre} (positif ou négatif) à partir des données ou des signes cliniques décrits.
\item \textbf{Identifier le mécanisme physiologique} :
\begin{itemize}
\item Bilan positif $\rightarrow$ stockage des excédents sous forme de triglycérides dans les adipocytes $\rightarrow$ hypertrophie du tissu adipeux $\rightarrow$ surpoids puis obésité $\rightarrow$ complications (diabète de type 2, hypertension artérielle, maladies cardiovasculaires, arthrose).
\item Bilan négatif $\rightarrow$ mobilisation des réserves (glycogène puis lipides puis protéines musculaires) $\rightarrow$ amaigrissement, fonte musculaire, fatigue, baisse des défenses immunitaires $\rightarrow$ dénutrition, épuisement.
\end{itemize}
\item \textbf{Relier chaque signe clinique au mécanisme} : par exemple, la fatigue chronique s'explique par l'insuffisance d'ATP fournie aux cellules ; les œdèmes par l'hypoalbuminémie (foie ne synthétise plus assez d'albumine faute d'acides aminés).
\item \textbf{Conclure} en nommant précisément l'état (surpoids, obésité, dénutrition, épuisement) et en citant au moins une complication.
\end{enumerate}
\end{methode}

\begin{exoresolu}[Interprétation d'un cas clinique de dénutrition]
Énoncé. Dans un centre de santé de la région de l'Extrême-Nord, un enfant de 10 ans présente : une masse de \SI{20}{kg} (norme : \SI{30}{kg}), une fatigue permanente, des infections respiratoires à répétition et des muscles très peu développés. Son alimentation, pauvre et peu variée, couvre environ 60 \% de ses besoins énergétiques.
\begin{enumerate}
\item Nommer le trouble dont souffre cet enfant et justifier.
\item Expliquer l'origine de deux de ses symptômes.
\end{enumerate}
\tcblower
Solution détaillée.
\begin{enumerate}
\item \textbf{Diagnostic.} L'enfant souffre de \cle{dénutrition} (malnutrition par carence énergétique et protéique, type marasme). Justification : ses apports ne couvrent que 60 \% de ses dépenses, donc son bilan énergétique est fortement négatif ($B < 0$), ce qui se traduit par un déficit pondéral majeur (\SI{20}{kg} au lieu de \SI{30}{kg}, soit un déficit de 33 \%) et une fonte musculaire.
\item \textbf{Explication des symptômes.}
\begin{itemize}
\item \emph{Fatigue permanente et fonte musculaire :} les apports étant insuffisants, l'organisme puise d'abord dans ses réserves de glycogène et de lipides ; celles-ci épuisées, il dégrade les protéines de ses propres muscles pour produire de l'énergie (néoglucogenèse à partir des acides aminés). D'où l'atrophie musculaire et l'asthénie.
\item \emph{Infections à répétition :} la carence en protéines et en énergie empêche la fabrication normale des anticorps (immunoglobulines) et des cellules immunitaires (lymphocytes) : les défenses de l'organisme s'effondrent (immunodépression), ce qui favorise les infections respiratoires et digestives.
\end{itemize}
\end{enumerate}
\textbf{Conclusion :} le bilan énergétique chroniquement négatif de cet enfant explique l'ensemble du tableau clinique : amaigrissement, fonte musculaire, épuisement et immunodépression.
\end{exoresolu}

\courssub{Rôle de l'activité physique dans la régulation du bilan énergétique}

\begin{propriete}[Effets de l'activité physique sur le bilan et la santé]
L'activité physique est le \textbf{seul levier volontaire majeur} sur les dépenses (le MB est imposé).
\begin{itemize}
\item \textbf{Augmentation de la DET :} Dépense immédiate pendant l'effort (kJ/min selon intensité) + \textbf{EPOC} (Excès de consommation d'O$_2$ post-effort, "afterburn") + augmentation du MB à long terme par gain de masse musculaire.
\item \textbf{Régulation de l'appétit :} Amélioration de la sensibilité à la leptine et à l'insuline ; régulation court terme (PYY, GLP-1, ghréline).
\item \textbf{Partitionnement des nutriments :} Favorise l'oxydation des lipides et le stockage du glycogène musculaire plutôt que la lipogenèse hépatique.
\item \textbf{Prévention des complications :} Baisse PA, amélioration profil lipidique, sensibilité insuline, densité osseuse, santé mentale.
\end{itemize}
\end{propriete}

\begin{methode}[Quantifier l'effet d'une activité physique sur le bilan]
\begin{enumerate}
\item Relever la dépense énergétique de l'activité (donnée en kJ/min ou kJ/heure) et sa durée.
\item Calculer la dépense supplémentaire : $E_{\text{activité}} = \text{dépense unitaire} \times \text{durée}$. Vérifier la cohérence des unités (convertir les heures en minutes si nécessaire).
\item Recalculer le bilan en ajoutant cette dépense aux dépenses totales : $B' = \text{Apports} - (\text{Dépenses} + E_{\text{activité}})$.
\item Comparer $B$ et $B'$ et conclure sur le rôle régulateur de l'activité physique : elle augmente les dépenses et peut ramener le bilan à l'équilibre sans réduire excessivement les apports.
\end{enumerate}
\end{methode}

\begin{exoresolu}[Effet de la marche quotidienne sur le bilan énergétique]
Énoncé. Reprenons le cas de l'employée de Douala : apports $=$ \SI{10500}{kJ}/jour ; dépenses $=$ \SI{10000}{kJ}/jour. Elle décide de marcher à pas rapides 30 minutes par jour. La marche rapide dépense environ \SI{20}{kJ/min}.
\begin{enumerate}
\item Calculer la dépense supplémentaire quotidienne due à la marche.
\item Établir le nouveau bilan énergétique et conclure.
\end{enumerate}
\tcblower
Solution détaillée.
\begin{enumerate}
\item \textbf{Dépense supplémentaire.} $E_{\text{marche}} = 20 \times 30 = \SI{600}{kJ}$ par jour.
\item \textbf{Nouveau bilan.}
\[
B' = 10500 - (10000 + 600) = 10500 - 10600 = -\SI{100}{kJ}
\]
Le bilan passe de $+\SI{500}{kJ}$ à $-\SI{100}{kJ}$ : il devient légèrement négatif. Au lieu de prendre du poids, l'employée en perdra très lentement (environ \SI{1}{kg} en 300 jours), et surtout elle stoppe la prise de masse.
\end{enumerate}
\textbf{Conclusion :} 30 minutes de marche rapide quotidienne suffisent à inverser le bilan énergétique de cette employée. L'activité physique est donc un régulateur puissant du bilan énergétique : elle augmente les dépenses et permet de prévenir le surpoids sans régime drastique.
\end{exoresolu}

\begin{savaistu}[Recommandations OMS et contexte camerounais]
L'OMS recommande aux adultes (18--64 ans) : \textbf{150 à 300 min/semaine d'activité d'endurance modérée} (marche rapide, vélo, danse, travaux champs) \textbf{OU 75--150 min/semaine d'intensité soutenue} + \textbf{renforcement musculaire 2 fois/semaine}. Au Cameroun, la marche pour aller au travail/marché, l'agriculture, le "sport du dimanche" (football, course) sont des leviers accessibles. La s\'edentarit\'e (écrans, moto-taxi, bureau) est le nouveau facteur de risque majeur en ville.
\end{savaistu}

\courssub{Rééquilibrage du bilan énergétique et mesures préventives}

\begin{methode}[Construire un programme de rééquilibrage]
\begin{enumerate}
\item \textbf{Analyser la situation} : sens et ampleur du déséquilibre, causes (alimentation excessive ou insuffisante, sédentarité, pauvreté, maladie).
\item \textbf{Fixer l'objectif chiffré} : ramener le bilan à l'équilibre ($B \approx 0$) ou créer un déficit/excédent modéré et progressif (jamais brutal : max 0,5--1 kg/semaine perte, 0,5 kg/semaine prise).
\item \textbf{Agir sur les deux leviers} :
\begin{itemize}
\item les \emph{apports} : réduire les aliments très énergétiques (huiles, fritures, boissons sucrées, alcool) en cas de bilan positif ; enrichir et diversifier l'alimentation (légumineuses : haricots, niébé ; poissons : capitaine, carpe ; fruits : papaye, mangue ; tubercules : manioc, igname, patate douce) en cas de bilan négatif ;
\item les \emph{dépenses} : augmenter l'activité physique (marche, sport, travaux) en cas de bilan positif ; limiter les dépenses excessives et se reposer en cas d'épuisement.
\end{itemize}
\item \textbf{Vérifier par le calcul} que les mesures proposées ramènent bien le bilan vers l'équilibre.
\item \textbf{Préciser le caractère progressif et durable} des mesures, et l'accompagnement médical si nécessaire (nutritionniste, centre de santé, PEC SAM/MAM pour enfants).
\end{enumerate}
\end{methode}

\begin{exoresolu}[Rééquilibrage du cas de Mme Ngo Bell]
Énoncé. Mme Ngo Bell (IMC $=$ 32 kg/m$^2$) a des apports quotidiens de \SI{11000}{kJ} pour des dépenses de \SI{9800}{kJ}. Son médecin lui conseille de supprimer les boissons sucrées ($-\SI{800}{kJ}$/jour) et de pratiquer 45 min de marche rapide par jour (\SI{20}{kJ/min}).
\begin{enumerate}
\item Calculer son bilan énergétique actuel.
\item Vérifier que les mesures proposées rééquilibrent son bilan.
\item Citer une précaution à respecter dans ce rééquilibrage.
\end{enumerate}
\tcblower
Solution détaillée.
\begin{enumerate}
\item \textbf{Bilan actuel.} $B = 11000 - 9800 = +\SI{1200}{kJ}$/jour. Ce fort excédent quotidien explique son obésité : il correspond à une prise théorique de $\frac{1200 \times 30}{30000} = \SI{1,2}{kg}$ par mois.
\item \textbf{Vérification des mesures.}
\begin{itemize}
\item Nouveaux apports : $11000 - 800 = \SI{10200}{kJ}$.
\item Nouvelles dépenses : $9800 + (20 \times 45) = 9800 + 900 = \SI{10700}{kJ}$.
\item Nouveau bilan : $B' = 10200 - 10700 = -\SI{500}{kJ}$/jour.
\end{itemize}
Le bilan devient modérément négatif : Mme Ngo Bell perdra environ $\frac{500 \times 60}{30000} = \SI{1}{kg}$ en 60 jours, soit un amaigrissement lent et régulier (\SI{0,5}{kg}/mois).
\item \textbf{Précaution.} Le rééquilibrage doit être \emph{progressif} : un déficit trop brutal (régime drastique $<$ 800 kcal/j) provoquerait fonte musculaire, carences en vitamines/sels minéraux, calculs biliaires et reprise de poids (effet yoyo). L'alimentation doit rester variée et équilibrée (protéines, vitamines, sels minéraux, fibres), et le suivi médical doit être maintenu.
\end{enumerate}
\textbf{Conclusion :} en agissant simultanément sur les apports ($-\SI{800}{kJ}$) et les dépenses ($+\SI{900}{kJ}$), les mesures proposées font passer le bilan de $+\SI{1200}{kJ}$ à $-\SI{500}{kJ}$ par jour : elles permettent un amaigrissement progressif d'environ \SI{0,5}{kg} par mois, compatible avec la santé.
\end{exoresolu}

\begin{attention}[Les 5 erreurs qui coûtent des points au Probatoire]
\begin{enumerate}
\item \textbf{Oublier d'élever la taille au carré dans l'IMC.} $\mathrm{IMC} = \dfrac{m}{T^2}$, et non $\dfrac{m}{T}$. De plus, la taille doit être convertie en \textbf{mètres} : pour \SI{170}{cm}, on utilise $T = 1,70$ m, donc $T^2 = 2,89$ m$^2$.
\item \textbf{Inverser la formule du bilan.} Le bilan se calcule $B = \text{Apports} - \text{Dépenses}$, jamais l'inverse. Un bilan positif signifie \emph{stockage} (excédent), un bilan négatif signifie \emph{puisement dans les réserves}. Une inversion de signe fausse toute l'interprétation.
\item \textbf{Additionner des masses au lieu d'énergies.} On ne peut pas additionner des grammes de glucides, lipides et protéines : il faut d'abord convertir chaque masse en énergie avec les coefficients (17, 38, 17 kJ/g), puis additionner les énergies. Ne pas oublier les unités (kJ) à chaque étape.
\item \textbf{Confondre dénutrition et sous-alimentation quantitative seule.} La dénutrition résulte d'un bilan énergétique chroniquement négatif ; elle entraîne non seulement l'amaigrissement, mais aussi la fonte musculaire (destruction des protéines corporelles) et la baisse des défenses immunitaires. Une rédaction qui ne cite que la « maigreur » est incomplète.
\item \textbf{Proposer un rééquilibrage non chiffré ou brutal.} Au Probatoire, toute mesure proposée doit être \emph{vérifiée par le calcul} (montrer que le nouveau bilan tend vers l'équilibre) et qualifiée de \emph{progressive}. Écrire « il faut manger moins » sans calcul ni nuance (diversification alimentaire, activité physique, suivi médical) fait perdre des points de savoir-être et de rigueur.
\end{enumerate}
\end{attention}

\newpage

\coursec{Exercices}

\courssub{Je vérifie mes connaissances}

\begin{exercice}[QCM : les bases du bilan énergétique]
Pour chaque question, choisis la (ou les) bonne(s) réponse(s).
\begin{enumerate}
\item Le bilan énergétique d'un organisme compare :
\begin{enumerate}
\item la masse d'aliments ingérée et la masse de déchets éliminée ;
\item les apports énergétiques alimentaires et les dépenses énergétiques totales ;
\item la quantité de dioxygène consommée et celle de dioxyde de carbone rejetée.
\end{enumerate}
\item La valeur énergétique des lipides est d'environ :
\begin{enumerate}
\item \SI{17}{kJ/g} ; \item \SI{4}{kJ/g} ; \item \SI{38}{kJ/g}.
\end{enumerate}
\item Le métabolisme de base correspond :
\begin{enumerate}
\item à l'énergie dépensée pendant un match de football ;
\item à l'énergie minimale nécessaire au fonctionnement de l'organisme au repos, à jeun, à température neutre ;
\item à l'énergie contenue dans un repas équilibré.
\end{enumerate}
\item Un bilan énergétique durablement positif entraîne :
\begin{enumerate}
\item une dénutrition ; \item un stockage de graisses et une prise de masse ; \item un épuisement des réserves.
\end{enumerate}
\end{enumerate}
\end{exercice}

\begin{exercice}[Vrai ou faux, en justifiant]
Indique si chaque affirmation est vraie ou fausse, puis justifie ta réponse en une ou deux phrases.
\begin{enumerate}
\item Le métabolisme de base représente la plus grande partie des dépenses énergétiques d'une personne sédentaire.
\item Un adolescent en pleine croissance a les mêmes besoins énergétiques qu'un adulte sédentaire.
\item Lorsque les dépenses dépassent les apports, l'organisme puise d'abord dans ses réserves de glycogène puis dans ses graisses.
\item L'activité physique n'a aucune influence sur le bilan énergétique car elle ne modifie que la masse musculaire.
\end{enumerate}
\end{exercice}

\begin{exercice}[Définitions et postes de dépenses]
\begin{enumerate}
\item Définis les termes suivants : \emph{bilan énergétique}, \emph{métabolisme de base}, \emph{dénutrition}.
\item Cite les trois postes de dépenses énergétiques de l'organisme et précise, pour chacun, un facteur qui le fait varier.
\item Rappelle la valeur énergétique d'un gramme de glucides et d'un gramme de protéines.
\end{enumerate}
\end{exercice}

\begin{exercice}[Calculs directs]
\begin{enumerate}
\item Un en-cas contient \SI{20}{g} de glucides, \SI{8}{g} de lipides et \SI{5}{g} de protéines. Calcule son apport énergétique total en kJ (données : glucides et protéines : \SI{17}{kJ/g} ; lipides : \SI{38}{kJ/g}).
\item Une élève dépense \SI{9800}{kJ} par jour et absorbe \SI{10200}{kJ} par jour. Calcule son bilan énergétique journalier et précise s'il est positif, négatif ou nul.
\item Un cultivateur de la Menoua dépense \SI{14500}{kJ/j} et consomme des aliments apportant \SI{13000}{kJ/j}. Calcule son déficit énergétique hebdomadaire.
\end{enumerate}
\end{exercice}

\courssub{Je m'entraîne}

\begin{exercice}[La journée de Chantal]
Chantal, élève de Première C à Yaoundé, consomme en une journée : \SI{280}{g} de glucides, \SI{70}{g} de lipides et \SI{60}{g} de protéines. Ses dépenses sont : métabolisme de base \SI{5800}{kJ}, activités diverses (marche, cours, tâches domestiques) \SI{3200}{kJ}, thermorégulation et digestion \SI{900}{kJ}.
\begin{enumerate}
\item Calcule ses apports énergétiques totaux (glucides et protéines : \SI{17}{kJ/g} ; lipides : \SI{38}{kJ/g}).
\item Calcule ses dépenses énergétiques totales.
\item Établis son bilan énergétique et conclus sur l'évolution prévisible de sa masse si ce régime se poursuit.
\end{enumerate}
\end{exercice}

\begin{exercice}[Trois IMC à interpréter]
On donne les mensurations de trois adultes : M.~Abena : masse $m = \SI{92}{kg}$, taille $t = \SI{1,70}{m}$ ; Mme~Etoundi : $m = \SI{47}{kg}$, $t = \SI{1,65}{m}$ ; M.~Fotso : $m = \SI{68}{kg}$, $t = \SI{1,75}{m}$.
\begin{enumerate}
\item Calcule l'IMC de chaque personne ($\mathrm{IMC} = m/t^{2}$, avec $m$ en kg et $t$ en m).
\item Classe l'état nutritionnel de chacun à l'aide des repères suivants : IMC $< 18,5$ : dénutrition ; $18,5 \leq \mathrm{IMC} < 25$ : corpulence normale ; $25 \leq \mathrm{IMC} < 30$ : surpoids ; $\mathrm{IMC} \geq 30$ : obésité.
\item Propose pour chacune des deux personnes en déséquilibre une mesure de rééquilibrage adaptée, en la justifiant.
\end{enumerate}
\end{exercice}

\begin{exercice}[L'athlète en déficit]
Une sprinteuse de l'équipe nationale, en stage de préparation, dépense \SI{15000}{kJ/jour} mais n'absorbe que \SI{11000}{kJ/jour} par son alimentation.
\begin{enumerate}
\item Calcule son déficit énergétique journalier, puis sur un mois de 30 jours.
\item Explique, en nommant les réserves mobilisées, les conséquences physiologiques de ce déficit s'il se prolonge (sur la masse corporelle, les muscles, les performances).
\item Sachant qu'un gramme de lipides corporels libère environ \SI{38}{kJ}, estime la perte de masse graisseuse théorique sur 30 jours.
\item Propose un ajustement alimentaire chiffré pour rétablir l'équilibre (composition en glucides, lipides, protéines de l'apport supplémentaire).
\end{enumerate}
\end{exercice}

\begin{exercice}[Rééquilibrer un menu]
Un employé de bureau à Douala, sédentaire, dépense \SI{9200}{kJ/jour}. Son alimentation quotidienne apporte \SI{11800}{kJ}, dont \SI{110}{g} de lipides.
\begin{enumerate}
\item Calcule son excédent énergétique journalier.
\item Explique pourquoi cet excédent favorise le surpoids, en précisant la forme de stockage de l'énergie excédentaire.
\item Propose deux modifications concrètes et chiffrées (une alimentaire, une liée à l'activité physique) pour ramener son bilan à l'équilibre, sachant qu'une heure de marche rapide dépense environ \SI{1100}{kJ}.
\end{enumerate}
\end{exercice}

\courssub{Je me prépare au Probatoire}

\begin{exercice}[Exploitation de données : la journée d'Arnaud]
Arnaud, élève de Première TI au lycée de Bafoussam, a noté ses repas et ses activités d'une journée type.

\textbf{Document 1 : apports alimentaires de la journée}
\begin{center}
\begin{tabularx}{\linewidth}{|l|X|c|c|c|}
\hline
\rowcolor{popblueL} \textbf{Repas} & \textbf{Composition} & \textbf{Glucides (g)} & \textbf{Lipides (g)} & \textbf{Protéines (g)} \\
\hline
Matin & beignets-haricots, bouillie & 90 & 25 & 15 \\
\hline
Midi & eru, water fufu & 110 & 35 & 25 \\
\hline
Soir & riz sauté, fruit & 80 & 15 & 10 \\
\hline
\end{tabularx}
\end{center}

\textbf{Document 2 : dépenses de la journée} : métabolisme de base : \SI{6500}{kJ} ; trajets à pied et jeux : \SI{2600}{kJ} ; thermorégulation et digestion : \SI{1000}{kJ}.

Données : glucides et protéines : \SI{17}{kJ/g} ; lipides : \SI{38}{kJ/g}.
\begin{enumerate}
\item Calcule les apports énergétiques totaux d'Arnaud pour cette journée, en détaillant la contribution de chaque type de nutriment.
\item Calcule ses dépenses énergétiques totales.
\item Établis le bilan énergétique de la journée et précise sa nature (positif, négatif ou équilibré).
\item Déduis-en le risque encouru par Arnaud si cette journée est représentative de son mode de vie, en nommant précisément le trouble possible et son mécanisme.
\item Propose deux mesures de rééquilibrage adaptées à sa situation d'élève, en les justifiant.
\end{enumerate}
\end{exercice}

\begin{exercice}[Étude documentaire : l'obésité en milieu urbain camerounais]
\textbf{Document.} \emph{« Dans les grandes villes du Cameroun comme Douala et Yaoundé, les enquêtes nutritionnelles révèlent une augmentation régulière du surpoids et de l'obésité, y compris chez les adolescents. Plusieurs facteurs sont mis en cause : la consommation croissante de plats de rue très gras et de boissons sucrées, la réduction des déplacements à pied au profit des motos-taxis et des transports motorisés, la diminution des espaces de jeu et la sédentarité liée aux écrans. Dans le même temps, dans certaines zones rurales en insécurité alimentaire, la dénutrition touche encore des enfants dont les apports ne couvrent pas les dépenses. »}
\begin{enumerate}
\item À partir du document, relève deux causes de déséquilibre énergétique \emph{positif} et précise pour chacune si elle agit sur les apports ou sur les dépenses.
\item Explique, en t'appuyant sur la notion de bilan énergétique, comment ces causes conduisent au surpoids puis à l'obésité (nomme la forme de stockage de l'énergie excédentaire).
\item Explique pourquoi, dans les zones rurales décrites, le déséquilibre est de nature opposée, et cite deux conséquences physiologiques de la dénutrition chez l'enfant.
\item Rédige un court paragraphe argumenté (8 à 12 lignes) proposant une stratégie de prévention de l'obésité adaptée aux élèves d'un lycée urbain camerounais, en mobilisant au moins trois savoirs de la leçon.
\end{enumerate}
\end{exercice}

\begin{exercice}[Problème de synthèse : deux situations contrastées]
\textbf{Situation A.} Mme~Ngo Bell, commerçante au marché Mokolo, marche et porte des charges toute la journée : elle dépense \SI{12500}{kJ/jour}. Ses repas, souvent pris à la hâte, ne lui apportent que \SI{9800}{kJ/jour}.

\textbf{Situation B.} Son fils Jean-Claude, employé de banque, se déplace en voiture et passe ses soirées devant la télévision : il dépense \SI{9000}{kJ/jour} et absorbe \SI{12100}{kJ/jour}. Il mesure \SI{1,72}{m} et pèse \SI{95}{kg}.

Données : glucides et protéines : \SI{17}{kJ/g} ; lipides : \SI{38}{kJ/g} ; une heure de marche rapide dépense environ \SI{1100}{kJ}.
\begin{enumerate}
\item Calcule le bilan énergétique journalier de Mme~Ngo Bell et celui de Jean-Claude.
\item Calcule l'IMC de Jean-Claude et qualifie son état nutritionnel à l'aide des repères : IMC $< 18,5$ : dénutrition ; $18,5 \leq \mathrm{IMC} < 25$ : corpulence normale ; $25 \leq \mathrm{IMC} < 30$ : surpoids ; $\mathrm{IMC} \geq 30$ : obésité.
\item Décris, pour chacune des deux personnes, les conséquences physiologiques prévisibles à terme de son déséquilibre, en nommant les réserves concernées et les troubles possibles.
\item Élabore pour Jean-Claude un programme de rééquilibrage chiffré sur une journée : nouvel apport alimentaire visé (en kJ) et durée quotidienne de marche rapide nécessaire, en supposant qu'il réduit ses apports de \SI{1600}{kJ} et compense le reste par l'activité physique.
\item Explique en quelques lignes pourquoi l'activité physique est un régulateur privilégié du bilan énergétique, au-delà de la simple dépense supplémentaire.
\end{enumerate}
\end{exercice}

\newpage

\coursec{Corrigés des exercices}

\begin{corrige}[Exercice 1 — QCM : les bases du bilan énergétique]
\tcblower
\begin{enumerate}
\item \textbf{Bonne réponse : b.} Le bilan énergétique compare les \cle{apports énergétiques} fournis par l'alimentation et les \cle{dépenses énergétiques} totales de l'organisme. La réponse a compare des masses (et non des énergies) ; la réponse c décrit les échanges gazeux respiratoires, qui ne sont qu'un reflet indirect de la dépense.
\item \textbf{Bonne réponse : c.} Un gramme de lipides libère environ \SI{38}{kJ} (soit \SI{9}{kcal/g}), contre \SI{17}{kJ/g} pour les glucides et les protéines. Les lipides sont donc les nutriments les plus énergétiques.
\item \textbf{Bonne réponse : b.} Le \cle{métabolisme de base} est l'énergie minimale dépensée par l'organisme \emph{au repos, à jeun et à température neutre}, pour assurer les fonctions vitales (battements cardiaques, respiration, maintien de la température corporelle, activité cérébrale).
\item \textbf{Bonne réponse : b.} Un bilan durablement positif signifie que les apports excèdent les dépenses : l'énergie excédentaire est stockée sous forme de \cle{triglycérides} dans le tissu adipeux, d'où une prise de masse (surpoids puis obésité). Les réponses a et c décrivent au contraire les conséquences d'un bilan négatif.
\end{enumerate}
\end{corrige}

\begin{corrige}[Exercice 2 — Vrai ou faux, en justifiant]
\tcblower
\begin{enumerate}
\item \textbf{Vrai.} Chez une personne sédentaire, le métabolisme de base représente environ 60 à 70 \% des dépenses énergétiques totales, devant l'activité physique et la thermorégulation/digestion.
\item \textbf{Faux.} L'adolescent en croissance doit fabriquer de la matière vivante (os, muscles, organes) : ses besoins énergétiques sont \emph{supérieurs} à ceux d'un adulte sédentaire, d'autant plus qu'il est souvent plus actif.
\item \textbf{Vrai.} L'organisme mobilise d'abord le \cle{glycogène} (réserve glucidique du foie et des muscles), puis, si le déficit persiste, les \cle{triglycérides} du tissu adipeux, et enfin, en cas de jeûne prolongé, les protéines musculaires.
\item \textbf{Faux.} L'activité physique est un poste majeur et \emph{modifiable} des dépenses énergétiques : elle augmente directement la dépense et élève aussi le métabolisme de base à long terme (la masse musculaire consomme de l'énergie même au repos). Elle influence donc fortement le bilan énergétique.
\end{enumerate}
\end{corrige}

\begin{corrige}[Exercice 3 — Définitions et postes de dépenses]
\tcblower
\begin{enumerate}
\item \textbf{Définitions.}
\begin{itemize}
\item \cle{Bilan énergétique} : comparaison, sur une période donnée, entre les apports énergétiques alimentaires et les dépenses énergétiques totales de l'organisme. Il est positif (apports $>$ dépenses), négatif (apports $<$ dépenses) ou nul (équilibre).
\item \cle{Métabolisme de base} : dépense énergétique minimale de l'organisme au repos, à jeun, à température neutre, indispensable aux fonctions vitales.
\item \cle{Dénutrition} : état pathologique résultant d'un bilan énergétique durablement négatif, caractérisé par un épuisement des réserves (glycogène, graisses puis protéines), un amaigrissement et une altération des fonctions de l'organisme.
\end{itemize}
\item \textbf{Les trois postes de dépenses} :
\begin{itemize}
\item le \textbf{métabolisme de base} — varie avec l'âge, le sexe ou la masse musculaire (il est plus élevé chez l'adolescent et chez l'homme) ;
\item l'\textbf{activité physique} — varie avec l'intensité et la durée de l'effort (un cultivateur dépense bien plus qu'un employé de bureau) ;
\item la \textbf{thermorégulation et la digestion} (thermogénèse) — varient avec la température extérieure et la nature des repas.
\end{itemize}
\item Glucides : \SI{17}{kJ/g} ; protéines : \SI{17}{kJ/g} (soit environ \SI{4}{kcal/g} chacun).
\end{enumerate}
\end{corrige}

\begin{corrige}[Exercice 4 — Calculs directs]
\tcblower
\begin{enumerate}
\item Apport de l'en-cas :
\begin{align*}
E &= (20 \times 17) + (8 \times 38) + (5 \times 17)\\
&= 340 + 304 + 85 = \SI{729}{kJ}.
\end{align*}
\item Bilan journalier : $B = \text{apports} - \text{dépenses} = 10200 - 9800 = +\SI{400}{kJ}$. Le bilan est \textbf{positif} : à terme, risque de prise de masse.
\item Déficit journalier : $14500 - 13000 = \SI{1500}{kJ/j}$. Déficit hebdomadaire : $1500 \times 7 = \SI{10500}{kJ}$.
\end{enumerate}
\end{corrige}

\begin{corrige}[Exercice 5 — La journée de Chantal]
\tcblower
\begin{enumerate}
\item \textbf{Apports énergétiques totaux :}
\begin{align*}
E_{\text{glucides}} &= 280 \times 17 = \SI{4760}{kJ}\\
E_{\text{lipides}} &= 70 \times 38 = \SI{2660}{kJ}\\
E_{\text{protéines}} &= 60 \times 17 = \SI{1020}{kJ}\\
E_{\text{total}} &= 4760 + 2660 + 1020 = \SI{8440}{kJ}.
\end{align*}
\item \textbf{Dépenses totales :} $5800 + 3200 + 900 = \SI{9900}{kJ}$.
\item \textbf{Bilan :} $B = 8440 - 9900 = -\SI{1460}{kJ}$ : bilan \textbf{négatif}. Les apports ne couvrent pas les dépenses ; si ce régime se poursuit, Chantal puisera dans ses réserves (glycogène puis graisses) et sa masse corporelle \textbf{diminuera}, avec risque de fatigue et de baisse de concentration en classe.
\end{enumerate}
\end{corrige}

\begin{corrige}[Exercice 6 — Trois IMC à interpréter]
\tcblower
\begin{enumerate}
\item \textbf{Calculs des IMC} ($\mathrm{IMC} = m/t^{2}$) :
\begin{itemize}
\item M. Abena : $\mathrm{IMC} = \dfrac{92}{1,70^{2}} = \dfrac{92}{2,89} \approx 31,8$ ;
\item Mme Etoundi : $\mathrm{IMC} = \dfrac{47}{1,65^{2}} = \dfrac{47}{2,7225} \approx 17,3$ ;
\item M. Fotso : $\mathrm{IMC} = \dfrac{68}{1,75^{2}} = \dfrac{68}{3,0625} \approx 22,2$.
\end{itemize}
\item \textbf{Classification :}
\begin{itemize}
\item M. Abena : $\mathrm{IMC} \approx 31,8 \geq 30$ : \textbf{obésité} ;
\item Mme Etoundi : $\mathrm{IMC} \approx 17,3 < 18,5$ : \textbf{dénutrition} ;
\item M. Fotso : $18,5 \leq \mathrm{IMC} \approx 22,2 < 25$ : \textbf{corpulence normale}.
\end{itemize}
\item \textbf{Mesures de rééquilibrage :}
\begin{itemize}
\item \emph{M. Abena (bilan positif chronique)} : réduire les apports (limiter fritures et boissons sucrées) et augmenter les dépenses par une activité physique régulière (marche rapide quotidienne), afin de rendre le bilan négatif et de mobiliser les graisses stockées.
\item \emph{Mme Etoundi (bilan négatif chronique)} : augmenter les apports énergétiques et protéiques (repas plus copieux et plus fréquents : haricots, arachides, poisson, bouillies enrichies) pour rendre le bilan positif et reconstituer les réserves, avec suivi médical.
\end{itemize}
\end{enumerate}
\end{corrige}

\begin{corrige}[Exercice 7 — L'athlète en déficit]
\tcblower
\begin{enumerate}
\item Déficit journalier : $15000 - 11000 = \SI{4000}{kJ/j}$. Sur 30 jours : $4000 \times 30 = \SI{120000}{kJ}$.
\item \textbf{Conséquences physiologiques :} l'organisme comble le déficit en mobilisant d'abord le \cle{glycogène} hépatique et musculaire (réserve glucidique rapidement épuisable), puis les \cle{triglycérides} du tissu adipeux (lipolyse). Si le déficit se prolonge, il dégrade aussi les \cle{protéines musculaires} (protéolyse) : fonte musculaire, perte de masse corporelle, fatigue chronique, baisse des performances (puissance et endurance réduites), risque de blessures et de baisse d'immunité.
\item Perte de masse graisseuse théorique :
\[ m = \dfrac{120000}{38} \approx \SI{3158}{g} \approx \SI{3,2}{kg} \text{ de graisse en 30 jours.} \]
\item \textbf{Ajustement alimentaire :} il faut combler \SI{4000}{kJ/j} en privilégiant les glucides (carburant du sprint). Proposition : $180$ g de glucides ($180 \times 17 = \SI{3060}{kJ}$) + $20$ g de lipides ($20 \times 38 = \SI{760}{kJ}$) + $10$ g de protéines ($10 \times 17 = \SI{170}{kJ}$), soit $3060 + 760 + 170 = \SI{3990}{kJ} \approx \SI{4000}{kJ}$. L'apport total passe à environ \SI{15000}{kJ/j}, ce qui rétablit l'équilibre.
\end{enumerate}
\end{corrige}

\begin{corrige}[Exercice 8 — Rééquilibrer un menu]
\tcblower
\begin{enumerate}
\item Excédent journalier : $B = 11800 - 9200 = +\SI{2600}{kJ}$.
\item Les apports dépassent durablement les dépenses : le glucose et les acides gras non utilisés sont transformés en \cle{triglycérides}, stockés dans les adipocytes du \cle{tissu adipeux}. L'accumulation progressive de ces graisses de réserve entraîne la prise de masse : surpoids puis obésité. Les \SI{110}{g} de lipides quotidiens (soit $110 \times 38 = \SI{4180}{kJ}$, plus de 45 \% des dépenses !) aggravent ce stockage.
\item \textbf{Deux modifications chiffrées :}
\begin{itemize}
\item \emph{Alimentaire :} réduire les apports de \SI{1500}{kJ}, par exemple en supprimant environ $40$ g de lipides ($40 \times 38 = \SI{1520}{kJ} \approx \SI{1500}{kJ}$) : moins d'huile de friture, cuisson grillée ou bouillie.
\item \emph{Activité physique :} dépenser les \SI{1100}{kJ} restants par \textbf{1 heure de marche rapide par jour} ($1 \times 1100 = \SI{1100}{kJ}$).
\end{itemize}
Vérification : nouveaux apports $= 11800 - 1500 = \SI{10300}{kJ}$ ; nouvelles dépenses $= 9200 + 1100 = \SI{10300}{kJ}$. Le bilan est désormais \textbf{nul} : l'équilibre est rétabli.
\end{enumerate}
\end{corrige}

\begin{corrige}[Exercice 9 — Exploitation de données : la journée d'Arnaud]
\tcblower
\begin{enumerate}
\item \textbf{Apports totaux.} Totaux par nutriment : glucides $= 90 + 110 + 80 = \SI{280}{g}$ ; lipides $= 25 + 35 + 15 = \SI{75}{g}$ ; protéines $= 15 + 25 + 10 = \SI{50}{g}$.
\begin{align*}
E_{\text{glucides}} &= 280 \times 17 = \SI{4760}{kJ}\\
E_{\text{lipides}} &= 75 \times 38 = \SI{2850}{kJ}\\
E_{\text{protéines}} &= 50 \times 17 = \SI{850}{kJ}\\
E_{\text{total}} &= 4760 + 2850 + 850 = \SI{8460}{kJ}.
\end{align*}
\item \textbf{Dépenses totales :} $6500 + 2600 + 1000 = \SI{10100}{kJ}$.
\item \textbf{Bilan :} $B = 8460 - 10100 = -\SI{1640}{kJ}$ : bilan \textbf{négatif}.
\item \textbf{Risque encouru :} si cette journée est représentative, le déficit chronique oblige l'organisme à mobiliser le glycogène, puis les graisses du tissu adipeux, puis les protéines musculaires : c'est le mécanisme de la \cle{dénutrition} — amaigrissement, fatigue, baisse de concentration et des résultats scolaires, fragilité face aux infections. C'est d'autant plus préoccupant qu'un adolescent en croissance a des besoins élevés.
\item \textbf{Mesures de rééquilibrage :}
\begin{itemize}
\item enrichir les repas : ajouter une collation énergétique en milieu de matinée (arachides, avocat, banane, bouillie enrichie) pour combler environ \SI{1700}{kJ} ;
\item ne pas sauter le petit-déjeuner et augmenter les portions de féculents (riz, foutou) aux repas, tout en gardant une activité physique modérée bénéfique à la santé.
\end{itemize}
\end{enumerate}
\textbf{Barème indicatif (10 pts) :} totaux par nutriment (2 pts) ; apports calculés avec unités (2 pts) ; dépenses (1 pt) ; bilan avec signe et conclusion (2 pts) ; mécanisme de la dénutrition nommé et expliqué (2 pts) ; mesures justifiées (1 pt).

\textbf{Points de vigilance :} exiger le détail par nutriment (pas seulement le total) ; le bilan doit être \emph{signé} et \emph{interprété} ; le trouble doit être \emph{nommé} (dénutrition) et son mécanisme (mobilisation des réserves) explicité ; les mesures doivent être adaptées au contexte d'un élève.
\end{corrige}

\begin{corrige}[Exercice 10 — Étude documentaire : l'obésité en milieu urbain camerounais]
\tcblower
\begin{enumerate}
\item \textbf{Deux causes de déséquilibre positif :}
\begin{itemize}
\item la consommation de plats de rue très gras et de boissons sucrées : elle agit en \emph{augmentant les apports} énergétiques ;
\item la réduction des déplacements à pied (motos-taxis) et la sédentarité liée aux écrans : elle agit en \emph{diminuant les dépenses} énergétiques.
\end{itemize}
\item \textbf{Du déséquilibre à l'obésité :} lorsque les apports dépassent durablement les dépenses, le bilan énergétique est positif. L'énergie excédentaire n'est pas éliminée : elle est convertie en \cle{triglycérides} et stockée dans les adipocytes du tissu adipeux. L'accumulation progressive de ces graisses fait augmenter la masse corporelle : l'IMC dépasse 25 (surpoids) puis 30 (obésité).
\item \textbf{Situation rurale :} dans les zones en insécurité alimentaire, les apports sont \emph{inférieurs} aux dépenses : le bilan est négatif, de nature opposée. Conséquences chez l'enfant : amaigrissement et arrêt ou ralentissement de la croissance ; fonte musculaire et fatigue ; baisse des défenses immunitaires (sensibilité accrue aux infections comme le paludisme).
\item \textbf{Paragraphe argumenté (proposition) :} Pour prévenir l'obésité chez les élèves d'un lycée urbain, il faut agir sur les deux termes du bilan énergétique. Du côté des apports, la cantine et les kiosques devraient proposer des repas moins gras et moins sucrés (fruits, arachides grillées plutôt que beignets et sodas), car glucides et lipides excédentaires sont stockés sous forme de triglycérides dans le tissu adipeux. Du côté des dépenses, l'activité physique étant le poste de dépense le plus modifiable, le lycée devrait développer le sport (cours d'EPS effectifs, tournois, club de marche) et encourager les trajets à pied : une heure de marche rapide dépense environ \SI{1100}{kJ}. Enfin, une éducation nutritionnelle (calcul de l'IMC, notion de métabolisme de base, lecture des étiquettes) permettrait à chaque élève de surveiller son propre bilan. En rétablissant l'égalité entre apports et dépenses, on maintient une corpulence normale ($18,5 \leq \mathrm{IMC} < 25$) et on prévient le surpoids et ses complications.
\end{enumerate}
\textbf{Barème indicatif (10 pts) :} causes relevées et correctement rattachées aux apports/dépenses (2 pts) ; mécanisme du stockage (triglycérides, tissu adipeux) (2 pts) ; bilan négatif rural et deux conséquences (2 pts) ; paragraphe mobilisant au moins trois savoirs, structuré et argumenté (4 pts).

\textbf{Points de vigilance :} distinguer clairement les causes agissant sur les apports de celles agissant sur les dépenses ; nommer la forme de stockage (triglycérides/adipocytes) ; le paragraphe doit mobiliser des \emph{savoirs de la leçon} (bilan, IMC, activité physique) et non rester un simple discours de bon sens.
\end{corrige}

\begin{corrige}[Exercice 11 — Problème de synthèse : deux situations contrastées]
\tcblower
\begin{enumerate}
\item \textbf{Bilans journaliers :}
\begin{itemize}
\item Mme Ngo Bell : $B_A = 9800 - 12500 = -\SI{2700}{kJ}$ : bilan \textbf{négatif} ;
\item Jean-Claude : $B_B = 12100 - 9000 = +\SI{3100}{kJ}$ : bilan \textbf{positif}.
\end{itemize}
\item \textbf{IMC de Jean-Claude :}
\[ \mathrm{IMC} = \dfrac{95}{1,72^{2}} = \dfrac{95}{2,9584} \approx 32,1. \]
$\mathrm{IMC} \geq 30$ : Jean-Claude est en situation d'\textbf{obésité}.
\item \textbf{Conséquences prévisibles :}
\begin{itemize}
\item \emph{Mme Ngo Bell (bilan négatif chronique)} : épuisement du glycogène, puis fonte des graisses du tissu adipeux, puis dégradation des protéines musculaires : amaigrissement, dénutrition, fatigue chronique, baisse de résistance aux infections — paradoxal pour un travail physiquement exigeant.
\item \emph{Jean-Claude (bilan positif chronique)} : stockage de l'énergie excédentaire sous forme de triglycérides dans le tissu adipeux : prise de masse, obésité installée, avec risques associés (diabète de type 2, hypertension artérielle, maladies cardiovasculaires).
\end{itemize}
\item \textbf{Programme de rééquilibrage chiffré :}
\begin{itemize}
\item Nouvel apport visé : $12100 - 1600 = \SI{10500}{kJ/j}$.
\item Reste à compenser par l'activité : $3100 - 1600 = \SI{1500}{kJ}$.
\item Durée de marche rapide : $\dfrac{1500}{1100} \approx \SI{1,36}{h}$, soit environ \textbf{1 h 20 min à 1 h 30 de marche rapide par jour}.
\end{itemize}
Vérification : dépenses $= 9000 + 1500 = \SI{10500}{kJ}$ $=$ apports : le bilan devient \textbf{nul}, l'équilibre est rétabli. (Pour perdre du poids, il faudrait même viser un léger bilan négatif.)
\item \textbf{L'activité physique, régulateur privilégié :} au-delà de la dépense immédiate, l'activité physique régulière augmente la masse musculaire, or le muscle consomme de l'énergie même au repos : elle élève donc durablement le métabolisme de base. Elle améliore aussi l'utilisation du glucose par les muscles et prévient les maladies associées à l'obésité. Contrairement aux régimes restrictifs seuls, elle agit sur les dépenses sans priver l'organisme des nutriments indispensables.
\end{enumerate}
\textbf{Barème indicatif (12 pts) :} deux bilans signés (2 pts) ; IMC calculé et interprété avec les repères (2 pts) ; conséquences nommant les réserves mobilisées pour chaque situation (3 pts) ; programme chiffré cohérent avec vérification (3 pts) ; explication du rôle régulateur de l'activité physique (2 pts).

\textbf{Points de vigilance :} le signe du bilan doit être explicité et interprété ; l'IMC doit être accompagné de la \emph{conclusion} (obésité) et non du seul nombre ; au 4., vérifier la cohérence finale (apports $=$ dépenses) ; au 5., aller au-delà de « ça fait dépenser des calories » en citant l'effet sur le métabolisme de base.
\end{corrige}

\newpage

\coursec{Fiche bilan}

\courssub{Carte mentale de la leçon}

\begin{popfigure}
\begin{tikzpicture}[mindmap, grow cyclic, every node/.style={concept}, concept color=popblue!70,
  root concept/.append style={concept color=popblue, font=\bfseries, text=white, minimum size=3.2cm},
  level 1/.append style={level distance=3.6cm, sibling angle=60, font=\small\bfseries},
  level 2/.append style={level distance=2.2cm, font=\scriptsize}]
\node[root concept]{Bilan énergétique de l'organisme}
  child[concept color=poporange!80]{ node {Besoins énergétiques}
    child { node {Métabolisme de base} }
    child { node {Activité physique} }
  }
  child[concept color=popgreen!80]{ node {Apports}
    child { node {Glucides, lipides, protéines} }
    child { node {Ration alimentaire} }
  }
  child[concept color=poppurple!80]{ node {Dépenses}
    child { node {MB $+$ effort $+$ thermorégulation} }
  }
  child[concept color=poppink!80]{ node {Déséquilibre $+$}
    child { node {Surpoids, obésité} }
    child { node {IMC $\geq$ 25} }
  }
  child[concept color=popgold!90]{ node {Déséquilibre $-$}
    child { node {Dénutrition, épuisement} }
  }
  child[concept color=popblue!60]{ node {Rééquilibrage}
    child { node {Sport $+$ diététique} }
  };
\end{tikzpicture}
\legende{Carte mentale de la leçon : le bilan énergétique compare les apports alimentaires aux dépenses de l'organisme ; son déséquilibre (positif ou négatif) entraîne des conséquences pathologiques que l'activité physique et l'hygiène alimentaire permettent de corriger.}
\end{popfigure}

\courssub{L'essentiel à retenir}

\begin{bilan}
\textbf{1. Définitions clés}
\begin{itemize}
  \item \cle{Bilan énergétique} : comparaison, sur une période donnée, entre l'\cle{énergie apportée} par les aliments et l'\cle{énergie dépensée} par l'organisme.
  \item \cle{Métabolisme de base (MB)} : dépense énergétique minimale au repos, nécessaire aux fonctions vitales (battements cardiaques, respiration, thermorégulation à \SI{37}{\celsius}). Il représente environ 60 à 70\,\% des dépenses totales.
  \item \cle{Ration alimentaire} : quantité d'aliments consommée en 24 h ; elle doit couvrir les besoins énergétiques et plastiques.
  \item \cle{Surpoids} et \cle{obésité} : excès de masse grasse résultant d'un bilan énergétique durablement positif.
  \item \cle{Dénutrition} : état pathologique résultant d'apports insuffisants (bilan négatif prolongé), pouvant aller jusqu'au \cle{kwashiorkor} (carence en protéines) ou au \cle{marasme} (carence énergétique globale).
\end{itemize}

\textbf{2. Formules et valeurs à connaître par cœur}

\begin{center}
\begin{tabularx}{\linewidth}{|l|X|}
\hline
\rowcolor{popblueL} \textbf{Grandeur} & \textbf{Expression / valeur} \\
\hline
Bilan énergétique & $B = A - D$ \quad ($A$ : apports ; $D$ : dépenses, en kJ ou kcal) \\
\hline
Équilibre & $B = 0$ : poids stable \\
\hline
Déséquilibre positif & $B > 0$ : stockage de graisses $\Rightarrow$ prise de poids \\
\hline
Déséquilibre négatif & $B < 0$ : puisage dans les réserves $\Rightarrow$ amaigrissement \\
\hline
Valeurs énergétiques & glucides : \SI{17}{kJ/g} ; protéines : \SI{17}{kJ/g} ; lipides : \SI{38}{kJ/g} \\
\hline
Besoins d'un adulte & environ \num{7500} à \SI{12500}{kJ/jour} selon sexe, âge, activité \\
\hline
IMC & $\mathrm{IMC} = \dfrac{m}{t^{2}}$ ($m$ en kg, $t$ en m) ; normal : 18,5 à 25 ; surpoids : 25 à 30 ; obésité : $\geq 30$ \\
\hline
Conversion & 1 kcal = \SI{4,18}{kJ} \\
\hline
\end{tabularx}
\end{center}

\textbf{3. Conséquences des déséquilibres}
\begin{itemize}
  \item \textbf{Bilan positif} : surcharge pondérale $\Rightarrow$ surcharge des articulations et du c\oe ur, risque de diabète de type 2, d'hypertension artérielle, de maladies cardiovasculaires.
  \item \textbf{Bilan négatif} : fonte musculaire, baisse des défenses immunitaires, fatigue chronique, retard de croissance chez l'enfant, \cle{épuisement} pouvant entraîner la mort.
\end{itemize}

\textbf{4. Méthodes express}
\begin{itemize}
  \item \textbf{Établir un bilan} : calculer $A$ (aliments $\times$ valeurs énergétiques), calculer $D$ (MB $+$ activités), comparer et conclure sur le signe de $B$.
  \item \textbf{Proposer un rééquilibrage} : si $B>0$ : réduire les aliments gras et sucrés, augmenter l'activité physique ; si $B<0$ : enrichir la ration (protéines, énergie), repos, suivi médical.
  \item \textbf{Rôle de l'activité physique} : elle augmente $D$, mobilise les graisses de réserve, développe la masse musculaire (qui élève le MB) : c'est un régulateur naturel du bilan.
\end{itemize}
\end{bilan}

\begin{aretenir}[Checklist avant le Probatoire]
\begin{itemize}
  \item[$\square$] Je sais définir le bilan énergétique et donner son expression $B = A - D$.
  \item[$\square$] Je connais la définition du métabolisme de base et sa part dans les dépenses.
  \item[$\square$] Je connais les valeurs énergétiques des glucides, protéines (\SI{17}{kJ/g}) et lipides (\SI{38}{kJ/g}).
  \item[$\square$] Je sais convertir kcal et kJ (1 kcal = \SI{4,18}{kJ}).
  \item[$\square$] Je sais calculer un IMC et interpréter le résultat (normal, surpoids, obésité).
  \item[$\square$] Je sais établir un bilan énergétique simple à partir d'un tableau de données.
  \item[$\square$] Je distingue dénutrition, kwashiorkor et marasme.
  \item[$\square$] Je cite au moins trois conséquences de l'obésité (diabète, hypertension, maladies cardiovasculaires).
  \item[$\square$] J'explique le rôle régulateur de l'activité physique sur le bilan énergétique.
  \item[$\square$] Je sais proposer des mesures de rééquilibrage adaptées à une situation concrète (alimentation $+$ exercice).
\end{itemize}
\end{aretenir}

\findecours

\end{document}
